\documentclass[11pt]{article}

    \usepackage[breakable]{tcolorbox}
    \usepackage{parskip} % Stop auto-indenting (to mimic markdown behaviour)
    

    % Basic figure setup, for now with no caption control since it's done
    % automatically by Pandoc (which extracts ![](path) syntax from Markdown).
    \usepackage{graphicx}
    % Keep aspect ratio if custom image width or height is specified
    \setkeys{Gin}{keepaspectratio}
    % Maintain compatibility with old templates. Remove in nbconvert 6.0
    \let\Oldincludegraphics\includegraphics
    % Ensure that by default, figures have no caption (until we provide a
    % proper Figure object with a Caption API and a way to capture that
    % in the conversion process - todo).
    \usepackage{caption}
    \DeclareCaptionFormat{nocaption}{}
    \captionsetup{format=nocaption,aboveskip=0pt,belowskip=0pt}

    \usepackage{float}
    \floatplacement{figure}{H} % forces figures to be placed at the correct location
    \usepackage{xcolor} % Allow colors to be defined
    \usepackage{enumerate} % Needed for markdown enumerations to work
    \usepackage{geometry} % Used to adjust the document margins
    \usepackage{amsmath} % Equations
    \usepackage{amssymb} % Equations
    \usepackage{textcomp} % defines textquotesingle
    % Hack from http://tex.stackexchange.com/a/47451/13684:
    \AtBeginDocument{%
        \def\PYZsq{\textquotesingle}% Upright quotes in Pygmentized code
    }
    \usepackage{upquote} % Upright quotes for verbatim code
    \usepackage{eurosym} % defines \euro

    \usepackage{iftex}
    \ifPDFTeX
        \usepackage[T1]{fontenc}
        \IfFileExists{alphabeta.sty}{
              \usepackage{alphabeta}
          }{
              \usepackage[mathletters]{ucs}
              \usepackage[utf8x]{inputenc}
          }
    \else
        \usepackage{fontspec}
        \usepackage{unicode-math}
    \fi

    \usepackage{fancyvrb} % verbatim replacement that allows latex
    \usepackage{grffile} % extends the file name processing of package graphics
                         % to support a larger range
    \makeatletter % fix for old versions of grffile with XeLaTeX
    \@ifpackagelater{grffile}{2019/11/01}
    {
      % Do nothing on new versions
    }
    {
      \def\Gread@@xetex#1{%
        \IfFileExists{"\Gin@base".bb}%
        {\Gread@eps{\Gin@base.bb}}%
        {\Gread@@xetex@aux#1}%
      }
    }
    \makeatother
    \usepackage[Export]{adjustbox} % Used to constrain images to a maximum size
    \adjustboxset{max size={0.9\linewidth}{0.9\paperheight}}

    % The hyperref package gives us a pdf with properly built
    % internal navigation ('pdf bookmarks' for the table of contents,
    % internal cross-reference links, web links for URLs, etc.)
    \usepackage{hyperref}
    % The default LaTeX title has an obnoxious amount of whitespace. By default,
    % titling removes some of it. It also provides customization options.
    \usepackage{titling}
    \usepackage{longtable} % longtable support required by pandoc >1.10
    \usepackage{booktabs}  % table support for pandoc > 1.12.2
    \usepackage{array}     % table support for pandoc >= 2.11.3
    \usepackage{calc}      % table minipage width calculation for pandoc >= 2.11.1
    \usepackage[inline]{enumitem} % IRkernel/repr support (it uses the enumerate* environment)
    \usepackage[normalem]{ulem} % ulem is needed to support strikethroughs (\sout)
                                % normalem makes italics be italics, not underlines
    \usepackage{soul}      % strikethrough (\st) support for pandoc >= 3.0.0
    \usepackage{mathrsfs}
    

    
    % Colors for the hyperref package
    \definecolor{urlcolor}{rgb}{0,.145,.698}
    \definecolor{linkcolor}{rgb}{.71,0.21,0.01}
    \definecolor{citecolor}{rgb}{.12,.54,.11}

    % ANSI colors
    \definecolor{ansi-black}{HTML}{3E424D}
    \definecolor{ansi-black-intense}{HTML}{282C36}
    \definecolor{ansi-red}{HTML}{E75C58}
    \definecolor{ansi-red-intense}{HTML}{B22B31}
    \definecolor{ansi-green}{HTML}{00A250}
    \definecolor{ansi-green-intense}{HTML}{007427}
    \definecolor{ansi-yellow}{HTML}{DDB62B}
    \definecolor{ansi-yellow-intense}{HTML}{B27D12}
    \definecolor{ansi-blue}{HTML}{208FFB}
    \definecolor{ansi-blue-intense}{HTML}{0065CA}
    \definecolor{ansi-magenta}{HTML}{D160C4}
    \definecolor{ansi-magenta-intense}{HTML}{A03196}
    \definecolor{ansi-cyan}{HTML}{60C6C8}
    \definecolor{ansi-cyan-intense}{HTML}{258F8F}
    \definecolor{ansi-white}{HTML}{C5C1B4}
    \definecolor{ansi-white-intense}{HTML}{A1A6B2}
    \definecolor{ansi-default-inverse-fg}{HTML}{FFFFFF}
    \definecolor{ansi-default-inverse-bg}{HTML}{000000}

    % common color for the border for error outputs.
    \definecolor{outerrorbackground}{HTML}{FFDFDF}

    % commands and environments needed by pandoc snippets
    % extracted from the output of `pandoc -s`
    \providecommand{\tightlist}{%
      \setlength{\itemsep}{0pt}\setlength{\parskip}{0pt}}
    \DefineVerbatimEnvironment{Highlighting}{Verbatim}{commandchars=\\\{\}}
    % Add ',fontsize=\small' for more characters per line
    \newenvironment{Shaded}{}{}
    \newcommand{\KeywordTok}[1]{\textcolor[rgb]{0.00,0.44,0.13}{\textbf{{#1}}}}
    \newcommand{\DataTypeTok}[1]{\textcolor[rgb]{0.56,0.13,0.00}{{#1}}}
    \newcommand{\DecValTok}[1]{\textcolor[rgb]{0.25,0.63,0.44}{{#1}}}
    \newcommand{\BaseNTok}[1]{\textcolor[rgb]{0.25,0.63,0.44}{{#1}}}
    \newcommand{\FloatTok}[1]{\textcolor[rgb]{0.25,0.63,0.44}{{#1}}}
    \newcommand{\CharTok}[1]{\textcolor[rgb]{0.25,0.44,0.63}{{#1}}}
    \newcommand{\StringTok}[1]{\textcolor[rgb]{0.25,0.44,0.63}{{#1}}}
    \newcommand{\CommentTok}[1]{\textcolor[rgb]{0.38,0.63,0.69}{\textit{{#1}}}}
    \newcommand{\OtherTok}[1]{\textcolor[rgb]{0.00,0.44,0.13}{{#1}}}
    \newcommand{\AlertTok}[1]{\textcolor[rgb]{1.00,0.00,0.00}{\textbf{{#1}}}}
    \newcommand{\FunctionTok}[1]{\textcolor[rgb]{0.02,0.16,0.49}{{#1}}}
    \newcommand{\RegionMarkerTok}[1]{{#1}}
    \newcommand{\ErrorTok}[1]{\textcolor[rgb]{1.00,0.00,0.00}{\textbf{{#1}}}}
    \newcommand{\NormalTok}[1]{{#1}}

    % Additional commands for more recent versions of Pandoc
    \newcommand{\ConstantTok}[1]{\textcolor[rgb]{0.53,0.00,0.00}{{#1}}}
    \newcommand{\SpecialCharTok}[1]{\textcolor[rgb]{0.25,0.44,0.63}{{#1}}}
    \newcommand{\VerbatimStringTok}[1]{\textcolor[rgb]{0.25,0.44,0.63}{{#1}}}
    \newcommand{\SpecialStringTok}[1]{\textcolor[rgb]{0.73,0.40,0.53}{{#1}}}
    \newcommand{\ImportTok}[1]{{#1}}
    \newcommand{\DocumentationTok}[1]{\textcolor[rgb]{0.73,0.13,0.13}{\textit{{#1}}}}
    \newcommand{\AnnotationTok}[1]{\textcolor[rgb]{0.38,0.63,0.69}{\textbf{\textit{{#1}}}}}
    \newcommand{\CommentVarTok}[1]{\textcolor[rgb]{0.38,0.63,0.69}{\textbf{\textit{{#1}}}}}
    \newcommand{\VariableTok}[1]{\textcolor[rgb]{0.10,0.09,0.49}{{#1}}}
    \newcommand{\ControlFlowTok}[1]{\textcolor[rgb]{0.00,0.44,0.13}{\textbf{{#1}}}}
    \newcommand{\OperatorTok}[1]{\textcolor[rgb]{0.40,0.40,0.40}{{#1}}}
    \newcommand{\BuiltInTok}[1]{{#1}}
    \newcommand{\ExtensionTok}[1]{{#1}}
    \newcommand{\PreprocessorTok}[1]{\textcolor[rgb]{0.74,0.48,0.00}{{#1}}}
    \newcommand{\AttributeTok}[1]{\textcolor[rgb]{0.49,0.56,0.16}{{#1}}}
    \newcommand{\InformationTok}[1]{\textcolor[rgb]{0.38,0.63,0.69}{\textbf{\textit{{#1}}}}}
    \newcommand{\WarningTok}[1]{\textcolor[rgb]{0.38,0.63,0.69}{\textbf{\textit{{#1}}}}}
    \makeatletter
    \newsavebox\pandoc@box
    \newcommand*\pandocbounded[1]{%
      \sbox\pandoc@box{#1}%
      % scaling factors for width and height
      \Gscale@div\@tempa\textheight{\dimexpr\ht\pandoc@box+\dp\pandoc@box\relax}%
      \Gscale@div\@tempb\linewidth{\wd\pandoc@box}%
      % select the smaller of both
      \ifdim\@tempb\p@<\@tempa\p@
        \let\@tempa\@tempb
      \fi
      % scaling accordingly (\@tempa < 1)
      \ifdim\@tempa\p@<\p@
        \scalebox{\@tempa}{\usebox\pandoc@box}%
      % scaling not needed, use as it is
      \else
        \usebox{\pandoc@box}%
      \fi
    }
    \makeatother

    % Define a nice break command that doesn't care if a line doesn't already
    % exist.
    \def\br{\hspace*{\fill} \\* }
    % Math Jax compatibility definitions
    \def\gt{>}
    \def\lt{<}
    \let\Oldtex\TeX
    \let\Oldlatex\LaTeX
    \renewcommand{\TeX}{\textrm{\Oldtex}}
    \renewcommand{\LaTeX}{\textrm{\Oldlatex}}
    % Document parameters
    % Document title
    \title{UT9\_Antonio\_Ferrer\_Martinez\_XXXXXXXXX}
    
    
    
    
    
    
    
% Pygments definitions
\makeatletter
\def\PY@reset{\let\PY@it=\relax \let\PY@bf=\relax%
    \let\PY@ul=\relax \let\PY@tc=\relax%
    \let\PY@bc=\relax \let\PY@ff=\relax}
\def\PY@tok#1{\csname PY@tok@#1\endcsname}
\def\PY@toks#1+{\ifx\relax#1\empty\else%
    \PY@tok{#1}\expandafter\PY@toks\fi}
\def\PY@do#1{\PY@bc{\PY@tc{\PY@ul{%
    \PY@it{\PY@bf{\PY@ff{#1}}}}}}}
\def\PY#1#2{\PY@reset\PY@toks#1+\relax+\PY@do{#2}}

\@namedef{PY@tok@w}{\def\PY@tc##1{\textcolor[rgb]{0.73,0.73,0.73}{##1}}}
\@namedef{PY@tok@c}{\let\PY@it=\textit\def\PY@tc##1{\textcolor[rgb]{0.24,0.48,0.48}{##1}}}
\@namedef{PY@tok@cp}{\def\PY@tc##1{\textcolor[rgb]{0.61,0.40,0.00}{##1}}}
\@namedef{PY@tok@k}{\let\PY@bf=\textbf\def\PY@tc##1{\textcolor[rgb]{0.00,0.50,0.00}{##1}}}
\@namedef{PY@tok@kp}{\def\PY@tc##1{\textcolor[rgb]{0.00,0.50,0.00}{##1}}}
\@namedef{PY@tok@kt}{\def\PY@tc##1{\textcolor[rgb]{0.69,0.00,0.25}{##1}}}
\@namedef{PY@tok@o}{\def\PY@tc##1{\textcolor[rgb]{0.40,0.40,0.40}{##1}}}
\@namedef{PY@tok@ow}{\let\PY@bf=\textbf\def\PY@tc##1{\textcolor[rgb]{0.67,0.13,1.00}{##1}}}
\@namedef{PY@tok@nb}{\def\PY@tc##1{\textcolor[rgb]{0.00,0.50,0.00}{##1}}}
\@namedef{PY@tok@nf}{\def\PY@tc##1{\textcolor[rgb]{0.00,0.00,1.00}{##1}}}
\@namedef{PY@tok@nc}{\let\PY@bf=\textbf\def\PY@tc##1{\textcolor[rgb]{0.00,0.00,1.00}{##1}}}
\@namedef{PY@tok@nn}{\let\PY@bf=\textbf\def\PY@tc##1{\textcolor[rgb]{0.00,0.00,1.00}{##1}}}
\@namedef{PY@tok@ne}{\let\PY@bf=\textbf\def\PY@tc##1{\textcolor[rgb]{0.80,0.25,0.22}{##1}}}
\@namedef{PY@tok@nv}{\def\PY@tc##1{\textcolor[rgb]{0.10,0.09,0.49}{##1}}}
\@namedef{PY@tok@no}{\def\PY@tc##1{\textcolor[rgb]{0.53,0.00,0.00}{##1}}}
\@namedef{PY@tok@nl}{\def\PY@tc##1{\textcolor[rgb]{0.46,0.46,0.00}{##1}}}
\@namedef{PY@tok@ni}{\let\PY@bf=\textbf\def\PY@tc##1{\textcolor[rgb]{0.44,0.44,0.44}{##1}}}
\@namedef{PY@tok@na}{\def\PY@tc##1{\textcolor[rgb]{0.41,0.47,0.13}{##1}}}
\@namedef{PY@tok@nt}{\let\PY@bf=\textbf\def\PY@tc##1{\textcolor[rgb]{0.00,0.50,0.00}{##1}}}
\@namedef{PY@tok@nd}{\def\PY@tc##1{\textcolor[rgb]{0.67,0.13,1.00}{##1}}}
\@namedef{PY@tok@s}{\def\PY@tc##1{\textcolor[rgb]{0.73,0.13,0.13}{##1}}}
\@namedef{PY@tok@sd}{\let\PY@it=\textit\def\PY@tc##1{\textcolor[rgb]{0.73,0.13,0.13}{##1}}}
\@namedef{PY@tok@si}{\let\PY@bf=\textbf\def\PY@tc##1{\textcolor[rgb]{0.64,0.35,0.47}{##1}}}
\@namedef{PY@tok@se}{\let\PY@bf=\textbf\def\PY@tc##1{\textcolor[rgb]{0.67,0.36,0.12}{##1}}}
\@namedef{PY@tok@sr}{\def\PY@tc##1{\textcolor[rgb]{0.64,0.35,0.47}{##1}}}
\@namedef{PY@tok@ss}{\def\PY@tc##1{\textcolor[rgb]{0.10,0.09,0.49}{##1}}}
\@namedef{PY@tok@sx}{\def\PY@tc##1{\textcolor[rgb]{0.00,0.50,0.00}{##1}}}
\@namedef{PY@tok@m}{\def\PY@tc##1{\textcolor[rgb]{0.40,0.40,0.40}{##1}}}
\@namedef{PY@tok@gh}{\let\PY@bf=\textbf\def\PY@tc##1{\textcolor[rgb]{0.00,0.00,0.50}{##1}}}
\@namedef{PY@tok@gu}{\let\PY@bf=\textbf\def\PY@tc##1{\textcolor[rgb]{0.50,0.00,0.50}{##1}}}
\@namedef{PY@tok@gd}{\def\PY@tc##1{\textcolor[rgb]{0.63,0.00,0.00}{##1}}}
\@namedef{PY@tok@gi}{\def\PY@tc##1{\textcolor[rgb]{0.00,0.52,0.00}{##1}}}
\@namedef{PY@tok@gr}{\def\PY@tc##1{\textcolor[rgb]{0.89,0.00,0.00}{##1}}}
\@namedef{PY@tok@ge}{\let\PY@it=\textit}
\@namedef{PY@tok@gs}{\let\PY@bf=\textbf}
\@namedef{PY@tok@ges}{\let\PY@bf=\textbf\let\PY@it=\textit}
\@namedef{PY@tok@gp}{\let\PY@bf=\textbf\def\PY@tc##1{\textcolor[rgb]{0.00,0.00,0.50}{##1}}}
\@namedef{PY@tok@go}{\def\PY@tc##1{\textcolor[rgb]{0.44,0.44,0.44}{##1}}}
\@namedef{PY@tok@gt}{\def\PY@tc##1{\textcolor[rgb]{0.00,0.27,0.87}{##1}}}
\@namedef{PY@tok@err}{\def\PY@bc##1{{\setlength{\fboxsep}{\string -\fboxrule}\fcolorbox[rgb]{1.00,0.00,0.00}{1,1,1}{\strut ##1}}}}
\@namedef{PY@tok@kc}{\let\PY@bf=\textbf\def\PY@tc##1{\textcolor[rgb]{0.00,0.50,0.00}{##1}}}
\@namedef{PY@tok@kd}{\let\PY@bf=\textbf\def\PY@tc##1{\textcolor[rgb]{0.00,0.50,0.00}{##1}}}
\@namedef{PY@tok@kn}{\let\PY@bf=\textbf\def\PY@tc##1{\textcolor[rgb]{0.00,0.50,0.00}{##1}}}
\@namedef{PY@tok@kr}{\let\PY@bf=\textbf\def\PY@tc##1{\textcolor[rgb]{0.00,0.50,0.00}{##1}}}
\@namedef{PY@tok@bp}{\def\PY@tc##1{\textcolor[rgb]{0.00,0.50,0.00}{##1}}}
\@namedef{PY@tok@fm}{\def\PY@tc##1{\textcolor[rgb]{0.00,0.00,1.00}{##1}}}
\@namedef{PY@tok@vc}{\def\PY@tc##1{\textcolor[rgb]{0.10,0.09,0.49}{##1}}}
\@namedef{PY@tok@vg}{\def\PY@tc##1{\textcolor[rgb]{0.10,0.09,0.49}{##1}}}
\@namedef{PY@tok@vi}{\def\PY@tc##1{\textcolor[rgb]{0.10,0.09,0.49}{##1}}}
\@namedef{PY@tok@vm}{\def\PY@tc##1{\textcolor[rgb]{0.10,0.09,0.49}{##1}}}
\@namedef{PY@tok@sa}{\def\PY@tc##1{\textcolor[rgb]{0.73,0.13,0.13}{##1}}}
\@namedef{PY@tok@sb}{\def\PY@tc##1{\textcolor[rgb]{0.73,0.13,0.13}{##1}}}
\@namedef{PY@tok@sc}{\def\PY@tc##1{\textcolor[rgb]{0.73,0.13,0.13}{##1}}}
\@namedef{PY@tok@dl}{\def\PY@tc##1{\textcolor[rgb]{0.73,0.13,0.13}{##1}}}
\@namedef{PY@tok@s2}{\def\PY@tc##1{\textcolor[rgb]{0.73,0.13,0.13}{##1}}}
\@namedef{PY@tok@sh}{\def\PY@tc##1{\textcolor[rgb]{0.73,0.13,0.13}{##1}}}
\@namedef{PY@tok@s1}{\def\PY@tc##1{\textcolor[rgb]{0.73,0.13,0.13}{##1}}}
\@namedef{PY@tok@mb}{\def\PY@tc##1{\textcolor[rgb]{0.40,0.40,0.40}{##1}}}
\@namedef{PY@tok@mf}{\def\PY@tc##1{\textcolor[rgb]{0.40,0.40,0.40}{##1}}}
\@namedef{PY@tok@mh}{\def\PY@tc##1{\textcolor[rgb]{0.40,0.40,0.40}{##1}}}
\@namedef{PY@tok@mi}{\def\PY@tc##1{\textcolor[rgb]{0.40,0.40,0.40}{##1}}}
\@namedef{PY@tok@il}{\def\PY@tc##1{\textcolor[rgb]{0.40,0.40,0.40}{##1}}}
\@namedef{PY@tok@mo}{\def\PY@tc##1{\textcolor[rgb]{0.40,0.40,0.40}{##1}}}
\@namedef{PY@tok@ch}{\let\PY@it=\textit\def\PY@tc##1{\textcolor[rgb]{0.24,0.48,0.48}{##1}}}
\@namedef{PY@tok@cm}{\let\PY@it=\textit\def\PY@tc##1{\textcolor[rgb]{0.24,0.48,0.48}{##1}}}
\@namedef{PY@tok@cpf}{\let\PY@it=\textit\def\PY@tc##1{\textcolor[rgb]{0.24,0.48,0.48}{##1}}}
\@namedef{PY@tok@c1}{\let\PY@it=\textit\def\PY@tc##1{\textcolor[rgb]{0.24,0.48,0.48}{##1}}}
\@namedef{PY@tok@cs}{\let\PY@it=\textit\def\PY@tc##1{\textcolor[rgb]{0.24,0.48,0.48}{##1}}}

\def\PYZbs{\char`\\}
\def\PYZus{\char`\_}
\def\PYZob{\char`\{}
\def\PYZcb{\char`\}}
\def\PYZca{\char`\^}
\def\PYZam{\char`\&}
\def\PYZlt{\char`\<}
\def\PYZgt{\char`\>}
\def\PYZsh{\char`\#}
\def\PYZpc{\char`\%}
\def\PYZdl{\char`\$}
\def\PYZhy{\char`\-}
\def\PYZsq{\char`\'}
\def\PYZdq{\char`\"}
\def\PYZti{\char`\~}
% for compatibility with earlier versions
\def\PYZat{@}
\def\PYZlb{[}
\def\PYZrb{]}
\makeatother


    % For linebreaks inside Verbatim environment from package fancyvrb.
    \makeatletter
        \newbox\Wrappedcontinuationbox
        \newbox\Wrappedvisiblespacebox
        \newcommand*\Wrappedvisiblespace {\textcolor{red}{\textvisiblespace}}
        \newcommand*\Wrappedcontinuationsymbol {\textcolor{red}{\llap{\tiny$\m@th\hookrightarrow$}}}
        \newcommand*\Wrappedcontinuationindent {3ex }
        \newcommand*\Wrappedafterbreak {\kern\Wrappedcontinuationindent\copy\Wrappedcontinuationbox}
        % Take advantage of the already applied Pygments mark-up to insert
        % potential linebreaks for TeX processing.
        %        {, <, #, %, $, ' and ": go to next line.
        %        _, }, ^, &, >, - and ~: stay at end of broken line.
        % Use of \textquotesingle for straight quote.
        \newcommand*\Wrappedbreaksatspecials {%
            \def\PYGZus{\discretionary{\char`\_}{\Wrappedafterbreak}{\char`\_}}%
            \def\PYGZob{\discretionary{}{\Wrappedafterbreak\char`\{}{\char`\{}}%
            \def\PYGZcb{\discretionary{\char`\}}{\Wrappedafterbreak}{\char`\}}}%
            \def\PYGZca{\discretionary{\char`\^}{\Wrappedafterbreak}{\char`\^}}%
            \def\PYGZam{\discretionary{\char`\&}{\Wrappedafterbreak}{\char`\&}}%
            \def\PYGZlt{\discretionary{}{\Wrappedafterbreak\char`\<}{\char`\<}}%
            \def\PYGZgt{\discretionary{\char`\>}{\Wrappedafterbreak}{\char`\>}}%
            \def\PYGZsh{\discretionary{}{\Wrappedafterbreak\char`\#}{\char`\#}}%
            \def\PYGZpc{\discretionary{}{\Wrappedafterbreak\char`\%}{\char`\%}}%
            \def\PYGZdl{\discretionary{}{\Wrappedafterbreak\char`\$}{\char`\$}}%
            \def\PYGZhy{\discretionary{\char`\-}{\Wrappedafterbreak}{\char`\-}}%
            \def\PYGZsq{\discretionary{}{\Wrappedafterbreak\textquotesingle}{\textquotesingle}}%
            \def\PYGZdq{\discretionary{}{\Wrappedafterbreak\char`\"}{\char`\"}}%
            \def\PYGZti{\discretionary{\char`\~}{\Wrappedafterbreak}{\char`\~}}%
        }
        % Some characters . , ; ? ! / are not pygmentized.
        % This macro makes them "active" and they will insert potential linebreaks
        \newcommand*\Wrappedbreaksatpunct {%
            \lccode`\~`\.\lowercase{\def~}{\discretionary{\hbox{\char`\.}}{\Wrappedafterbreak}{\hbox{\char`\.}}}%
            \lccode`\~`\,\lowercase{\def~}{\discretionary{\hbox{\char`\,}}{\Wrappedafterbreak}{\hbox{\char`\,}}}%
            \lccode`\~`\;\lowercase{\def~}{\discretionary{\hbox{\char`\;}}{\Wrappedafterbreak}{\hbox{\char`\;}}}%
            \lccode`\~`\:\lowercase{\def~}{\discretionary{\hbox{\char`\:}}{\Wrappedafterbreak}{\hbox{\char`\:}}}%
            \lccode`\~`\?\lowercase{\def~}{\discretionary{\hbox{\char`\?}}{\Wrappedafterbreak}{\hbox{\char`\?}}}%
            \lccode`\~`\!\lowercase{\def~}{\discretionary{\hbox{\char`\!}}{\Wrappedafterbreak}{\hbox{\char`\!}}}%
            \lccode`\~`\/\lowercase{\def~}{\discretionary{\hbox{\char`\/}}{\Wrappedafterbreak}{\hbox{\char`\/}}}%
            \catcode`\.\active
            \catcode`\,\active
            \catcode`\;\active
            \catcode`\:\active
            \catcode`\?\active
            \catcode`\!\active
            \catcode`\/\active
            \lccode`\~`\~
        }
    \makeatother

    \let\OriginalVerbatim=\Verbatim
    \makeatletter
    \renewcommand{\Verbatim}[1][1]{%
        %\parskip\z@skip
        \sbox\Wrappedcontinuationbox {\Wrappedcontinuationsymbol}%
        \sbox\Wrappedvisiblespacebox {\FV@SetupFont\Wrappedvisiblespace}%
        \def\FancyVerbFormatLine ##1{\hsize\linewidth
            \vtop{\raggedright\hyphenpenalty\z@\exhyphenpenalty\z@
                \doublehyphendemerits\z@\finalhyphendemerits\z@
                \strut ##1\strut}%
        }%
        % If the linebreak is at a space, the latter will be displayed as visible
        % space at end of first line, and a continuation symbol starts next line.
        % Stretch/shrink are however usually zero for typewriter font.
        \def\FV@Space {%
            \nobreak\hskip\z@ plus\fontdimen3\font minus\fontdimen4\font
            \discretionary{\copy\Wrappedvisiblespacebox}{\Wrappedafterbreak}
            {\kern\fontdimen2\font}%
        }%

        % Allow breaks at special characters using \PYG... macros.
        \Wrappedbreaksatspecials
        % Breaks at punctuation characters . , ; ? ! and / need catcode=\active
        \OriginalVerbatim[#1,codes*=\Wrappedbreaksatpunct]%
    }
    \makeatother

    % Exact colors from NB
    \definecolor{incolor}{HTML}{303F9F}
    \definecolor{outcolor}{HTML}{D84315}
    \definecolor{cellborder}{HTML}{CFCFCF}
    \definecolor{cellbackground}{HTML}{F7F7F7}

    % prompt
    \makeatletter
    \newcommand{\boxspacing}{\kern\kvtcb@left@rule\kern\kvtcb@boxsep}
    \makeatother
    \newcommand{\prompt}[4]{
        {\ttfamily\llap{{\color{#2}[#3]:\hspace{3pt}#4}}\vspace{-\baselineskip}}
    }
    

    
    % Prevent overflowing lines due to hard-to-break entities
    \sloppy
    % Setup hyperref package
    \hypersetup{
      breaklinks=true,  % so long urls are correctly broken across lines
      colorlinks=true,
      urlcolor=urlcolor,
      linkcolor=linkcolor,
      citecolor=citecolor,
      }
    % Slightly bigger margins than the latex defaults
    
    \geometry{verbose,tmargin=1in,bmargin=1in,lmargin=1in,rmargin=1in}
    
    

\begin{document}
    
    \maketitle
    
    

    
    \section{Trabajo Final UT9 - Auto Machine
Learning}\label{trabajo-final-ut9---auto-machine-learning}

\textbf{Alumno:} autor \textbf{DNI:} XXXXXXXXX

\textbf{Enlace al vídeo:}

\begin{center}\rule{0.5\linewidth}{0.5pt}\end{center}

\subsection{Datasets utilizados}\label{datasets-utilizados}

El enunciado pide utilizar los \textbf{mismos datasets de las prácticas
anteriores} (UT5 para regresión y UT6 para clasificación), para poder
comparar los resultados de XGBoost y AutoML con los modelos que ya
entrenamos previamente.

\subsubsection{1. House Prices (Regresión) --- mismo que
UT5}\label{house-prices-regresiuxf3n-mismo-que-ut5}

El mismo dataset de la competición de Kaggle ``House Prices - Advanced
Regression Techniques'' que usamos en la UT5. Contiene \textbf{1460
muestras} con 81 variables (calidad de materiales, año de construcción,
tipo de garaje, etc.). El objetivo es predecir el \textbf{precio de
venta} (\texttt{SalePrice}).

Se descarga automáticamente desde Kaggle con \texttt{kagglehub}, así la
comparación con los modelos de la UT5 es directa y justa.

\subsubsection{2. Social Network Ads (Clasificación) --- mismo que
UT6}\label{social-network-ads-clasificaciuxf3n-mismo-que-ut6}

En la UT6 trabajamos con un dataset llamado ``Ensayo de Motores'' que
tenía las columnas RPM, Par de carga y Fallo Motor. Buscando el dataset
para descargarlo automáticamente, descubrí que en realidad es el
conocido dataset \textbf{Social Network Ads} de Kaggle (subido por
Rakesh Rau), con las columnas renombradas:

{\def\LTcaptype{none} % do not increment counter
\begin{longtable}[]{@{}ll@{}}
\toprule\noalign{}
Nombre original (Kaggle) & Nombre en UT6 \\
\midrule\noalign{}
\endhead
\bottomrule\noalign{}
\endlastfoot
Age (edad del usuario) & RPM \\
EstimatedSalary (salario estimado) & Par de carga \\
Purchased (si compró el producto) & Fallo Motor \\
\end{longtable}
}

Las columnas \texttt{User\ ID} y \texttt{Gender} fueron eliminadas en la
versión de la UT6.

\textbf{¿Por qué uso el dataset original de Kaggle?} Por dos razones: 1.
\textbf{Comodidad:} Se descarga automáticamente con \texttt{kagglehub},
sin subir archivos manualmente 2. \textbf{Coherencia:} Los nombres
originales (Age, EstimatedSalary, Purchased) tienen mucho más sentido
que ``RPM'' y ``Par de carga'' para predecir una compra. Los datos son
exactamente los mismos, solo cambian los nombres de las columnas.

Ambos datasets se descargan automáticamente desde Kaggle.

\begin{center}\rule{0.5\linewidth}{0.5pt}\end{center}

\subsection{Índice}\label{uxedndice}

\begin{enumerate}
\def\labelenumi{\arabic{enumi}.}
\tightlist
\item
  \textbf{XGBoost}

  \begin{itemize}
  \tightlist
  \item
    1.1 XGBoost para Regresión (dataset House Prices)
  \item
    1.2 XGBoost para Clasificación (dataset Social Network Ads)
  \end{itemize}
\item
  \textbf{AutoML con FLAML}

  \begin{itemize}
  \tightlist
  \item
    2.1 Clasificación automática (dataset Social Network Ads)
  \end{itemize}
\item
  \textbf{Valores SHAP}

  \begin{itemize}
  \tightlist
  \item
    3.1 Explicar una única predicción
  \item
    3.2 Explicar el dataset completo
  \end{itemize}
\item
  \textbf{Conclusiones finales}
\end{enumerate}

\subsubsection{Nota sobre AutoML}\label{nota-sobre-automl}

El enunciado propone usar auto-sklearn o PyCaret. Sin embargo: -
\textbf{auto-sklearn} fue abandonado en 2022 y no funciona con Python
3.10 o superior - \textbf{PyCaret 3.3.2} no soporta Python 3.12, que es
la versión actual de Google Colab

Por eso se utiliza \textbf{FLAML} (Fast and Lightweight AutoML) de
Microsoft, que es compatible, ligero y cumple exactamente la misma
función: automatizar la selección de modelos y la búsqueda de
hiperparámetros.

    \subsection{0. Instalación de librerías e
imports}\label{instalaciuxf3n-de-libreruxedas-e-imports}

    \begin{tcolorbox}[breakable, size=fbox, boxrule=1pt, pad at break*=1mm,colback=cellbackground, colframe=cellborder]
\prompt{In}{incolor}{1}{\boxspacing}
\begin{Verbatim}[commandchars=\\\{\}]
\PY{c+c1}{\PYZsh{} Instalación de librerías necesarias}
\PY{o}{!}pip\PY{+w}{ }\PYZhy{}q\PY{+w}{ }install\PY{+w}{ }xgboost\PY{+w}{ }shap\PY{+w}{ }openpyxl\PY{+w}{ }kagglehub\PY{+w}{ }\PY{l+s+s2}{\PYZdq{}flaml[automl]\PYZdq{}}
\end{Verbatim}
\end{tcolorbox}

    \begin{Verbatim}[commandchars=\\\{\}]
   \textcolor{ansi-black-intense}{━━━━━━━━━━━━━━━━━━━━━━━━━━━━━━━━━━━━━━━━} \textcolor{ansi-green}{223.6/223.6 MB}
\textcolor{ansi-red}{5.3 MB/s} eta \textcolor{ansi-cyan}{0:00:00}:00:0100:01
   \textcolor{ansi-black-intense}{━━━━━━━━━━━━━━━━━━━━━━━━━━━━━━━━━━━━━━━━} \textcolor{ansi-green}{337.7/337.7 kB}
\textcolor{ansi-red}{26.2 MB/s} eta \textcolor{ansi-cyan}{0:00:00}

    \end{Verbatim}

    \begin{tcolorbox}[breakable, size=fbox, boxrule=1pt, pad at break*=1mm,colback=cellbackground, colframe=cellborder]
\prompt{In}{incolor}{2}{\boxspacing}
\begin{Verbatim}[commandchars=\\\{\}]
\PY{c+c1}{\PYZsh{} Imports generales}
\PY{k+kn}{import}\PY{+w}{ }\PY{n+nn}{pandas}\PY{+w}{ }\PY{k}{as}\PY{+w}{ }\PY{n+nn}{pd}
\PY{k+kn}{import}\PY{+w}{ }\PY{n+nn}{numpy}\PY{+w}{ }\PY{k}{as}\PY{+w}{ }\PY{n+nn}{np}
\PY{k+kn}{import}\PY{+w}{ }\PY{n+nn}{matplotlib}\PY{n+nn}{.}\PY{n+nn}{pyplot}\PY{+w}{ }\PY{k}{as}\PY{+w}{ }\PY{n+nn}{plt}
\PY{k+kn}{import}\PY{+w}{ }\PY{n+nn}{seaborn}\PY{+w}{ }\PY{k}{as}\PY{+w}{ }\PY{n+nn}{sns}

\PY{c+c1}{\PYZsh{} Scikit\PYZhy{}learn}
\PY{k+kn}{from}\PY{+w}{ }\PY{n+nn}{sklearn}\PY{n+nn}{.}\PY{n+nn}{model\PYZus{}selection}\PY{+w}{ }\PY{k+kn}{import} \PY{n}{train\PYZus{}test\PYZus{}split}
\PY{k+kn}{from}\PY{+w}{ }\PY{n+nn}{sklearn}\PY{n+nn}{.}\PY{n+nn}{preprocessing}\PY{+w}{ }\PY{k+kn}{import} \PY{n}{StandardScaler}\PY{p}{,} \PY{n}{LabelEncoder}
\PY{k+kn}{from}\PY{+w}{ }\PY{n+nn}{sklearn}\PY{n+nn}{.}\PY{n+nn}{metrics}\PY{+w}{ }\PY{k+kn}{import} \PY{p}{(}
    \PY{n}{mean\PYZus{}squared\PYZus{}error}\PY{p}{,} \PY{n}{r2\PYZus{}score}\PY{p}{,}
    \PY{n}{confusion\PYZus{}matrix}\PY{p}{,} \PY{n}{classification\PYZus{}report}\PY{p}{,} \PY{n}{accuracy\PYZus{}score}\PY{p}{,}
    \PY{n}{RocCurveDisplay}
\PY{p}{)}

\PY{c+c1}{\PYZsh{} XGBoost}
\PY{k+kn}{from}\PY{+w}{ }\PY{n+nn}{xgboost}\PY{+w}{ }\PY{k+kn}{import} \PY{n}{XGBRegressor}\PY{p}{,} \PY{n}{XGBClassifier}

\PY{c+c1}{\PYZsh{} SHAP}
\PY{k+kn}{import}\PY{+w}{ }\PY{n+nn}{shap}

\PY{k+kn}{import}\PY{+w}{ }\PY{n+nn}{warnings}
\PY{n}{warnings}\PY{o}{.}\PY{n}{filterwarnings}\PY{p}{(}\PY{l+s+s1}{\PYZsq{}}\PY{l+s+s1}{ignore}\PY{l+s+s1}{\PYZsq{}}\PY{p}{)}

\PY{n}{sns}\PY{o}{.}\PY{n}{set\PYZus{}theme}\PY{p}{(}\PY{n}{style}\PY{o}{=}\PY{l+s+s1}{\PYZsq{}}\PY{l+s+s1}{whitegrid}\PY{l+s+s1}{\PYZsq{}}\PY{p}{)}
\PY{n+nb}{print}\PY{p}{(}\PY{l+s+s1}{\PYZsq{}}\PY{l+s+s1}{Librerías importadas correctamente.}\PY{l+s+s1}{\PYZsq{}}\PY{p}{)}
\end{Verbatim}
\end{tcolorbox}

    \begin{Verbatim}[commandchars=\\\{\}]
Librerías importadas correctamente.
    \end{Verbatim}

    \begin{center}\rule{0.5\linewidth}{0.5pt}\end{center}

\section{1. Aplicar modelo XGBoost}\label{aplicar-modelo-xgboost}

XGBoost (eXtreme Gradient Boosting) es un algoritmo de ensemble basado
en árboles de decisión que utiliza gradient boosting. Básicamente,
entrena muchos árboles de decisión uno detrás de otro, donde cada nuevo
árbol intenta corregir los errores que cometieron los anteriores. Es de
los algoritmos más usados en competiciones de Kaggle por su buen
rendimiento.

    \subsection{1.1 XGBoost para Regresión}\label{xgboost-para-regresiuxf3n}

Utilizamos el mismo dataset \textbf{House Prices} de la competición de
Kaggle que en la UT5. El objetivo es predecir el precio de venta
(\texttt{SalePrice}) de viviendas a partir de 81 variables (calidad,
superficie, año, garaje, etc.).

    \subsubsection{1.1.1 Importar dataset}\label{importar-dataset}

    \begin{tcolorbox}[breakable, size=fbox, boxrule=1pt, pad at break*=1mm,colback=cellbackground, colframe=cellborder]
\prompt{In}{incolor}{3}{\boxspacing}
\begin{Verbatim}[commandchars=\\\{\}]
\PY{c+c1}{\PYZsh{} Configuramos las credenciales de Kaggle para poder descargar datasets}
\PY{k+kn}{import}\PY{+w}{ }\PY{n+nn}{os}
\PY{n}{os}\PY{o}{.}\PY{n}{environ}\PY{p}{[}\PY{l+s+s1}{\PYZsq{}}\PY{l+s+s1}{KAGGLE\PYZus{}USERNAME}\PY{l+s+s1}{\PYZsq{}}\PY{p}{]} \PY{o}{=} \PY{l+s+s1}{\PYZsq{}}\PY{l+s+s1}{autor}\PY{l+s+s1}{\PYZsq{}}
\PY{n}{os}\PY{o}{.}\PY{n}{environ}\PY{p}{[}\PY{l+s+s1}{\PYZsq{}}\PY{l+s+s1}{KAGGLE\PYZus{}KEY}\PY{l+s+s1}{\PYZsq{}}\PY{p}{]} \PY{o}{=} \PY{l+s+s1}{\PYZsq{}}\PY{l+s+s1}{6a52c6651f478c0cc2148110957d0843}\PY{l+s+s1}{\PYZsq{}}

\PY{c+c1}{\PYZsh{} Descargamos el MISMO dataset de House Prices que usamos en la UT5}
\PY{c+c1}{\PYZsh{} Es la competicion \PYZdq{}House Prices \PYZhy{} Advanced Regression Techniques\PYZdq{}}
\PY{c+c1}{\PYZsh{} 1460 casas con 81 variables, asi la comparacion con la UT5 es justa}
\PY{k+kn}{import}\PY{+w}{ }\PY{n+nn}{kagglehub}
\PY{n}{path} \PY{o}{=} \PY{n}{kagglehub}\PY{o}{.}\PY{n}{competition\PYZus{}download}\PY{p}{(}\PY{l+s+s1}{\PYZsq{}}\PY{l+s+s1}{house\PYZhy{}prices\PYZhy{}advanced\PYZhy{}regression\PYZhy{}techniques}\PY{l+s+s1}{\PYZsq{}}\PY{p}{)}
\PY{n+nb}{print}\PY{p}{(}\PY{l+s+s1}{\PYZsq{}}\PY{l+s+s1}{Dataset descargado en:}\PY{l+s+s1}{\PYZsq{}}\PY{p}{,} \PY{n}{path}\PY{p}{)}

\PY{k}{for} \PY{n}{f} \PY{o+ow}{in} \PY{n}{os}\PY{o}{.}\PY{n}{listdir}\PY{p}{(}\PY{n}{path}\PY{p}{)}\PY{p}{:}
    \PY{n+nb}{print}\PY{p}{(}\PY{n}{f}\PY{p}{)}
\end{Verbatim}
\end{tcolorbox}

    \begin{Verbatim}[commandchars=\\\{\}]
Downloading to /root/.cache/kagglehub/competitions/house-prices-advanced-
regression-techniques.archive{\ldots}
    \end{Verbatim}

    \begin{Verbatim}[commandchars=\\\{\}]
100\%|██████████| 199k/199k [00:00<00:00, 726kB/s]
    \end{Verbatim}

    \begin{Verbatim}[commandchars=\\\{\}]
Extracting files{\ldots}
Dataset descargado en: /root/.cache/kagglehub/competitions/house-prices-
advanced-regression-techniques
data\_description.txt
test.csv
train.csv
sample\_submission.csv
    \end{Verbatim}

    \begin{Verbatim}[commandchars=\\\{\}]

    \end{Verbatim}

    \begin{tcolorbox}[breakable, size=fbox, boxrule=1pt, pad at break*=1mm,colback=cellbackground, colframe=cellborder]
\prompt{In}{incolor}{4}{\boxspacing}
\begin{Verbatim}[commandchars=\\\{\}]
\PY{c+c1}{\PYZsh{} Cargamos el dataset de entrenamiento (train.csv)}
\PY{c+c1}{\PYZsh{} Es el mismo house\PYZus{}prices\PYZus{}train.csv que usamos en la UT5}
\PY{n}{df\PYZus{}reg} \PY{o}{=} \PY{n}{pd}\PY{o}{.}\PY{n}{read\PYZus{}csv}\PY{p}{(}\PY{n}{os}\PY{o}{.}\PY{n}{path}\PY{o}{.}\PY{n}{join}\PY{p}{(}\PY{n}{path}\PY{p}{,} \PY{l+s+s1}{\PYZsq{}}\PY{l+s+s1}{train.csv}\PY{l+s+s1}{\PYZsq{}}\PY{p}{)}\PY{p}{)}
\PY{n+nb}{print}\PY{p}{(}\PY{l+s+sa}{f}\PY{l+s+s1}{\PYZsq{}}\PY{l+s+s1}{Dimensiones del dataset: }\PY{l+s+si}{\PYZob{}}\PY{n}{df\PYZus{}reg}\PY{o}{.}\PY{n}{shape}\PY{l+s+si}{\PYZcb{}}\PY{l+s+s1}{\PYZsq{}}\PY{p}{)}
\PY{n}{df\PYZus{}reg}\PY{o}{.}\PY{n}{head}\PY{p}{(}\PY{p}{)}
\end{Verbatim}
\end{tcolorbox}

    \begin{Verbatim}[commandchars=\\\{\}]
Dimensiones del dataset: (1460, 81)
    \end{Verbatim}

            \begin{tcolorbox}[breakable, size=fbox, boxrule=.5pt, pad at break*=1mm, opacityfill=0]
\prompt{Out}{outcolor}{4}{\boxspacing}
\begin{Verbatim}[commandchars=\\\{\}]
   Id  MSSubClass MSZoning  LotFrontage  LotArea Street Alley LotShape  \textbackslash{}
0   1          60       RL         65.0     8450   Pave   NaN      Reg
1   2          20       RL         80.0     9600   Pave   NaN      Reg
2   3          60       RL         68.0    11250   Pave   NaN      IR1
3   4          70       RL         60.0     9550   Pave   NaN      IR1
4   5          60       RL         84.0    14260   Pave   NaN      IR1

  LandContour Utilities  {\ldots} PoolArea PoolQC Fence MiscFeature MiscVal MoSold  \textbackslash{}
0         Lvl    AllPub  {\ldots}        0    NaN   NaN         NaN       0      2
1         Lvl    AllPub  {\ldots}        0    NaN   NaN         NaN       0      5
2         Lvl    AllPub  {\ldots}        0    NaN   NaN         NaN       0      9
3         Lvl    AllPub  {\ldots}        0    NaN   NaN         NaN       0      2
4         Lvl    AllPub  {\ldots}        0    NaN   NaN         NaN       0     12

  YrSold  SaleType  SaleCondition  SalePrice
0   2008        WD         Normal     208500
1   2007        WD         Normal     181500
2   2008        WD         Normal     223500
3   2006        WD        Abnorml     140000
4   2008        WD         Normal     250000

[5 rows x 81 columns]
\end{Verbatim}
\end{tcolorbox}
        
    \begin{tcolorbox}[breakable, size=fbox, boxrule=1pt, pad at break*=1mm,colback=cellbackground, colframe=cellborder]
\prompt{In}{incolor}{5}{\boxspacing}
\begin{Verbatim}[commandchars=\\\{\}]
\PY{c+c1}{\PYZsh{} Información general del dataset}
\PY{n}{df\PYZus{}reg}\PY{o}{.}\PY{n}{info}\PY{p}{(}\PY{p}{)}
\PY{n+nb}{print}\PY{p}{(}\PY{l+s+s1}{\PYZsq{}}\PY{l+s+se}{\PYZbs{}n}\PY{l+s+s1}{\PYZhy{}\PYZhy{}\PYZhy{} Estadísticas descriptivas \PYZhy{}\PYZhy{}\PYZhy{}}\PY{l+s+s1}{\PYZsq{}}\PY{p}{)}
\PY{n}{df\PYZus{}reg}\PY{o}{.}\PY{n}{describe}\PY{p}{(}\PY{p}{)}
\end{Verbatim}
\end{tcolorbox}

    \begin{Verbatim}[commandchars=\\\{\}]
<class 'pandas.core.frame.DataFrame'>
RangeIndex: 1460 entries, 0 to 1459
Data columns (total 81 columns):
 \#   Column         Non-Null Count  Dtype
---  ------         --------------  -----
 0   Id             1460 non-null   int64
 1   MSSubClass     1460 non-null   int64
 2   MSZoning       1460 non-null   object
 3   LotFrontage    1201 non-null   float64
 4   LotArea        1460 non-null   int64
 5   Street         1460 non-null   object
 6   Alley          91 non-null     object
 7   LotShape       1460 non-null   object
 8   LandContour    1460 non-null   object
 9   Utilities      1460 non-null   object
 10  LotConfig      1460 non-null   object
 11  LandSlope      1460 non-null   object
 12  Neighborhood   1460 non-null   object
 13  Condition1     1460 non-null   object
 14  Condition2     1460 non-null   object
 15  BldgType       1460 non-null   object
 16  HouseStyle     1460 non-null   object
 17  OverallQual    1460 non-null   int64
 18  OverallCond    1460 non-null   int64
 19  YearBuilt      1460 non-null   int64
 20  YearRemodAdd   1460 non-null   int64
 21  RoofStyle      1460 non-null   object
 22  RoofMatl       1460 non-null   object
 23  Exterior1st    1460 non-null   object
 24  Exterior2nd    1460 non-null   object
 25  MasVnrType     588 non-null    object
 26  MasVnrArea     1452 non-null   float64
 27  ExterQual      1460 non-null   object
 28  ExterCond      1460 non-null   object
 29  Foundation     1460 non-null   object
 30  BsmtQual       1423 non-null   object
 31  BsmtCond       1423 non-null   object
 32  BsmtExposure   1422 non-null   object
 33  BsmtFinType1   1423 non-null   object
 34  BsmtFinSF1     1460 non-null   int64
 35  BsmtFinType2   1422 non-null   object
 36  BsmtFinSF2     1460 non-null   int64
 37  BsmtUnfSF      1460 non-null   int64
 38  TotalBsmtSF    1460 non-null   int64
 39  Heating        1460 non-null   object
 40  HeatingQC      1460 non-null   object
 41  CentralAir     1460 non-null   object
 42  Electrical     1459 non-null   object
 43  1stFlrSF       1460 non-null   int64
 44  2ndFlrSF       1460 non-null   int64
 45  LowQualFinSF   1460 non-null   int64
 46  GrLivArea      1460 non-null   int64
 47  BsmtFullBath   1460 non-null   int64
 48  BsmtHalfBath   1460 non-null   int64
 49  FullBath       1460 non-null   int64
 50  HalfBath       1460 non-null   int64
 51  BedroomAbvGr   1460 non-null   int64
 52  KitchenAbvGr   1460 non-null   int64
 53  KitchenQual    1460 non-null   object
 54  TotRmsAbvGrd   1460 non-null   int64
 55  Functional     1460 non-null   object
 56  Fireplaces     1460 non-null   int64
 57  FireplaceQu    770 non-null    object
 58  GarageType     1379 non-null   object
 59  GarageYrBlt    1379 non-null   float64
 60  GarageFinish   1379 non-null   object
 61  GarageCars     1460 non-null   int64
 62  GarageArea     1460 non-null   int64
 63  GarageQual     1379 non-null   object
 64  GarageCond     1379 non-null   object
 65  PavedDrive     1460 non-null   object
 66  WoodDeckSF     1460 non-null   int64
 67  OpenPorchSF    1460 non-null   int64
 68  EnclosedPorch  1460 non-null   int64
 69  3SsnPorch      1460 non-null   int64
 70  ScreenPorch    1460 non-null   int64
 71  PoolArea       1460 non-null   int64
 72  PoolQC         7 non-null      object
 73  Fence          281 non-null    object
 74  MiscFeature    54 non-null     object
 75  MiscVal        1460 non-null   int64
 76  MoSold         1460 non-null   int64
 77  YrSold         1460 non-null   int64
 78  SaleType       1460 non-null   object
 79  SaleCondition  1460 non-null   object
 80  SalePrice      1460 non-null   int64
dtypes: float64(3), int64(35), object(43)
memory usage: 924.0+ KB

--- Estadísticas descriptivas ---
    \end{Verbatim}

            \begin{tcolorbox}[breakable, size=fbox, boxrule=.5pt, pad at break*=1mm, opacityfill=0]
\prompt{Out}{outcolor}{5}{\boxspacing}
\begin{Verbatim}[commandchars=\\\{\}]
                Id   MSSubClass  LotFrontage        LotArea  OverallQual  \textbackslash{}
count  1460.000000  1460.000000  1201.000000    1460.000000  1460.000000
mean    730.500000    56.897260    70.049958   10516.828082     6.099315
std     421.610009    42.300571    24.284752    9981.264932     1.382997
min       1.000000    20.000000    21.000000    1300.000000     1.000000
25\%     365.750000    20.000000    59.000000    7553.500000     5.000000
50\%     730.500000    50.000000    69.000000    9478.500000     6.000000
75\%    1095.250000    70.000000    80.000000   11601.500000     7.000000
max    1460.000000   190.000000   313.000000  215245.000000    10.000000

       OverallCond    YearBuilt  YearRemodAdd   MasVnrArea   BsmtFinSF1  {\ldots}  \textbackslash{}
count  1460.000000  1460.000000   1460.000000  1452.000000  1460.000000  {\ldots}
mean      5.575342  1971.267808   1984.865753   103.685262   443.639726  {\ldots}
std       1.112799    30.202904     20.645407   181.066207   456.098091  {\ldots}
min       1.000000  1872.000000   1950.000000     0.000000     0.000000  {\ldots}
25\%       5.000000  1954.000000   1967.000000     0.000000     0.000000  {\ldots}
50\%       5.000000  1973.000000   1994.000000     0.000000   383.500000  {\ldots}
75\%       6.000000  2000.000000   2004.000000   166.000000   712.250000  {\ldots}
max       9.000000  2010.000000   2010.000000  1600.000000  5644.000000  {\ldots}

        WoodDeckSF  OpenPorchSF  EnclosedPorch    3SsnPorch  ScreenPorch  \textbackslash{}
count  1460.000000  1460.000000    1460.000000  1460.000000  1460.000000
mean     94.244521    46.660274      21.954110     3.409589    15.060959
std     125.338794    66.256028      61.119149    29.317331    55.757415
min       0.000000     0.000000       0.000000     0.000000     0.000000
25\%       0.000000     0.000000       0.000000     0.000000     0.000000
50\%       0.000000    25.000000       0.000000     0.000000     0.000000
75\%     168.000000    68.000000       0.000000     0.000000     0.000000
max     857.000000   547.000000     552.000000   508.000000   480.000000

          PoolArea       MiscVal       MoSold       YrSold      SalePrice
count  1460.000000   1460.000000  1460.000000  1460.000000    1460.000000
mean      2.758904     43.489041     6.321918  2007.815753  180921.195890
std      40.177307    496.123024     2.703626     1.328095   79442.502883
min       0.000000      0.000000     1.000000  2006.000000   34900.000000
25\%       0.000000      0.000000     5.000000  2007.000000  129975.000000
50\%       0.000000      0.000000     6.000000  2008.000000  163000.000000
75\%       0.000000      0.000000     8.000000  2009.000000  214000.000000
max     738.000000  15500.000000    12.000000  2010.000000  755000.000000

[8 rows x 38 columns]
\end{Verbatim}
\end{tcolorbox}
        
    \subsubsection{1.1.2 Manejar datos missing}\label{manejar-datos-missing}

    \begin{tcolorbox}[breakable, size=fbox, boxrule=1pt, pad at break*=1mm,colback=cellbackground, colframe=cellborder]
\prompt{In}{incolor}{6}{\boxspacing}
\begin{Verbatim}[commandchars=\\\{\}]
\PY{c+c1}{\PYZsh{} Comprobamos valores nulos}
\PY{n+nb}{print}\PY{p}{(}\PY{l+s+s1}{\PYZsq{}}\PY{l+s+s1}{Valores nulos por columna:}\PY{l+s+s1}{\PYZsq{}}\PY{p}{)}
\PY{n+nb}{print}\PY{p}{(}\PY{n}{df\PYZus{}reg}\PY{o}{.}\PY{n}{isnull}\PY{p}{(}\PY{p}{)}\PY{o}{.}\PY{n}{sum}\PY{p}{(}\PY{p}{)}\PY{p}{)}
\PY{n+nb}{print}\PY{p}{(}\PY{l+s+sa}{f}\PY{l+s+s1}{\PYZsq{}}\PY{l+s+se}{\PYZbs{}n}\PY{l+s+s1}{Total de valores nulos: }\PY{l+s+si}{\PYZob{}}\PY{n}{df\PYZus{}reg}\PY{o}{.}\PY{n}{isnull}\PY{p}{(}\PY{p}{)}\PY{o}{.}\PY{n}{sum}\PY{p}{(}\PY{p}{)}\PY{o}{.}\PY{n}{sum}\PY{p}{(}\PY{p}{)}\PY{l+s+si}{\PYZcb{}}\PY{l+s+s1}{\PYZsq{}}\PY{p}{)}
\end{Verbatim}
\end{tcolorbox}

    \begin{Verbatim}[commandchars=\\\{\}]
Valores nulos por columna:
Id                 0
MSSubClass         0
MSZoning           0
LotFrontage      259
LotArea            0
                {\ldots}
MoSold             0
YrSold             0
SaleType           0
SaleCondition      0
SalePrice          0
Length: 81, dtype: int64

Total de valores nulos: 7829
    \end{Verbatim}

    \begin{tcolorbox}[breakable, size=fbox, boxrule=1pt, pad at break*=1mm,colback=cellbackground, colframe=cellborder]
\prompt{In}{incolor}{7}{\boxspacing}
\begin{Verbatim}[commandchars=\\\{\}]
\PY{c+c1}{\PYZsh{} Si hubiera valores nulos los rellenamos:}
\PY{c+c1}{\PYZsh{} \PYZhy{} Numericas: con la mediana (mejor que la media porque no le afectan los outliers)}
\PY{c+c1}{\PYZsh{} \PYZhy{} Categoricas: con la moda (el valor mas frecuente)}
\PY{k}{for} \PY{n}{col} \PY{o+ow}{in} \PY{n}{df\PYZus{}reg}\PY{o}{.}\PY{n}{select\PYZus{}dtypes}\PY{p}{(}\PY{n}{include}\PY{o}{=}\PY{n}{np}\PY{o}{.}\PY{n}{number}\PY{p}{)}\PY{o}{.}\PY{n}{columns}\PY{p}{:}
    \PY{n}{df\PYZus{}reg}\PY{p}{[}\PY{n}{col}\PY{p}{]}\PY{o}{.}\PY{n}{fillna}\PY{p}{(}\PY{n}{df\PYZus{}reg}\PY{p}{[}\PY{n}{col}\PY{p}{]}\PY{o}{.}\PY{n}{median}\PY{p}{(}\PY{p}{)}\PY{p}{,} \PY{n}{inplace}\PY{o}{=}\PY{k+kc}{True}\PY{p}{)}

\PY{k}{for} \PY{n}{col} \PY{o+ow}{in} \PY{n}{df\PYZus{}reg}\PY{o}{.}\PY{n}{select\PYZus{}dtypes}\PY{p}{(}\PY{n}{include}\PY{o}{=}\PY{l+s+s1}{\PYZsq{}}\PY{l+s+s1}{object}\PY{l+s+s1}{\PYZsq{}}\PY{p}{)}\PY{o}{.}\PY{n}{columns}\PY{p}{:}
    \PY{n}{df\PYZus{}reg}\PY{p}{[}\PY{n}{col}\PY{p}{]}\PY{o}{.}\PY{n}{fillna}\PY{p}{(}\PY{n}{df\PYZus{}reg}\PY{p}{[}\PY{n}{col}\PY{p}{]}\PY{o}{.}\PY{n}{mode}\PY{p}{(}\PY{p}{)}\PY{p}{[}\PY{l+m+mi}{0}\PY{p}{]}\PY{p}{,} \PY{n}{inplace}\PY{o}{=}\PY{k+kc}{True}\PY{p}{)}

\PY{n+nb}{print}\PY{p}{(}\PY{l+s+sa}{f}\PY{l+s+s1}{\PYZsq{}}\PY{l+s+s1}{Valores nulos tras imputación: }\PY{l+s+si}{\PYZob{}}\PY{n}{df\PYZus{}reg}\PY{o}{.}\PY{n}{isnull}\PY{p}{(}\PY{p}{)}\PY{o}{.}\PY{n}{sum}\PY{p}{(}\PY{p}{)}\PY{o}{.}\PY{n}{sum}\PY{p}{(}\PY{p}{)}\PY{l+s+si}{\PYZcb{}}\PY{l+s+s1}{\PYZsq{}}\PY{p}{)}
\end{Verbatim}
\end{tcolorbox}

    \begin{Verbatim}[commandchars=\\\{\}]
Valores nulos tras imputación: 0
    \end{Verbatim}

    \subsubsection{1.1.3 Manejar datos
categóricos}\label{manejar-datos-categuxf3ricos}

Identificamos las variables categóricas y aplicamos codificación
adecuada: - \textbf{Variables binarias (yes/no):} Label Encoding -
\textbf{Variables con más categorías:} One-Hot Encoding con
\texttt{drop\_first=True} para evitar la trampa de las variables dummy

    \begin{tcolorbox}[breakable, size=fbox, boxrule=1pt, pad at break*=1mm,colback=cellbackground, colframe=cellborder]
\prompt{In}{incolor}{8}{\boxspacing}
\begin{Verbatim}[commandchars=\\\{\}]
\PY{c+c1}{\PYZsh{} Identificamos variables categóricas}
\PY{n}{cat\PYZus{}cols} \PY{o}{=} \PY{n}{df\PYZus{}reg}\PY{o}{.}\PY{n}{select\PYZus{}dtypes}\PY{p}{(}\PY{n}{include}\PY{o}{=}\PY{l+s+s1}{\PYZsq{}}\PY{l+s+s1}{object}\PY{l+s+s1}{\PYZsq{}}\PY{p}{)}\PY{o}{.}\PY{n}{columns}\PY{o}{.}\PY{n}{tolist}\PY{p}{(}\PY{p}{)}
\PY{n+nb}{print}\PY{p}{(}\PY{l+s+sa}{f}\PY{l+s+s1}{\PYZsq{}}\PY{l+s+s1}{Variables categóricas (}\PY{l+s+si}{\PYZob{}}\PY{n+nb}{len}\PY{p}{(}\PY{n}{cat\PYZus{}cols}\PY{p}{)}\PY{l+s+si}{\PYZcb{}}\PY{l+s+s1}{): }\PY{l+s+si}{\PYZob{}}\PY{n}{cat\PYZus{}cols}\PY{l+s+si}{\PYZcb{}}\PY{l+s+s1}{\PYZsq{}}\PY{p}{)}
\PY{k}{for} \PY{n}{col} \PY{o+ow}{in} \PY{n}{cat\PYZus{}cols}\PY{p}{:}
    \PY{n+nb}{print}\PY{p}{(}\PY{l+s+sa}{f}\PY{l+s+s1}{\PYZsq{}}\PY{l+s+s1}{  }\PY{l+s+si}{\PYZob{}}\PY{n}{col}\PY{l+s+si}{\PYZcb{}}\PY{l+s+s1}{: }\PY{l+s+si}{\PYZob{}}\PY{n}{df\PYZus{}reg}\PY{p}{[}\PY{n}{col}\PY{p}{]}\PY{o}{.}\PY{n}{unique}\PY{p}{(}\PY{p}{)}\PY{l+s+si}{\PYZcb{}}\PY{l+s+s1}{\PYZsq{}}\PY{p}{)}
\end{Verbatim}
\end{tcolorbox}

    \begin{Verbatim}[commandchars=\\\{\}]
Variables categóricas (43): ['MSZoning', 'Street', 'Alley', 'LotShape',
'LandContour', 'Utilities', 'LotConfig', 'LandSlope', 'Neighborhood',
'Condition1', 'Condition2', 'BldgType', 'HouseStyle', 'RoofStyle', 'RoofMatl',
'Exterior1st', 'Exterior2nd', 'MasVnrType', 'ExterQual', 'ExterCond',
'Foundation', 'BsmtQual', 'BsmtCond', 'BsmtExposure', 'BsmtFinType1',
'BsmtFinType2', 'Heating', 'HeatingQC', 'CentralAir', 'Electrical',
'KitchenQual', 'Functional', 'FireplaceQu', 'GarageType', 'GarageFinish',
'GarageQual', 'GarageCond', 'PavedDrive', 'PoolQC', 'Fence', 'MiscFeature',
'SaleType', 'SaleCondition']
  MSZoning: ['RL' 'RM' 'C (all)' 'FV' 'RH']
  Street: ['Pave' 'Grvl']
  Alley: ['Grvl' 'Pave']
  LotShape: ['Reg' 'IR1' 'IR2' 'IR3']
  LandContour: ['Lvl' 'Bnk' 'Low' 'HLS']
  Utilities: ['AllPub' 'NoSeWa']
  LotConfig: ['Inside' 'FR2' 'Corner' 'CulDSac' 'FR3']
  LandSlope: ['Gtl' 'Mod' 'Sev']
  Neighborhood: ['CollgCr' 'Veenker' 'Crawfor' 'NoRidge' 'Mitchel' 'Somerst'
'NWAmes'
 'OldTown' 'BrkSide' 'Sawyer' 'NridgHt' 'NAmes' 'SawyerW' 'IDOTRR'
 'MeadowV' 'Edwards' 'Timber' 'Gilbert' 'StoneBr' 'ClearCr' 'NPkVill'
 'Blmngtn' 'BrDale' 'SWISU' 'Blueste']
  Condition1: ['Norm' 'Feedr' 'PosN' 'Artery' 'RRAe' 'RRNn' 'RRAn' 'PosA'
'RRNe']
  Condition2: ['Norm' 'Artery' 'RRNn' 'Feedr' 'PosN' 'PosA' 'RRAn' 'RRAe']
  BldgType: ['1Fam' '2fmCon' 'Duplex' 'TwnhsE' 'Twnhs']
  HouseStyle: ['2Story' '1Story' '1.5Fin' '1.5Unf' 'SFoyer' 'SLvl' '2.5Unf'
'2.5Fin']
  RoofStyle: ['Gable' 'Hip' 'Gambrel' 'Mansard' 'Flat' 'Shed']
  RoofMatl: ['CompShg' 'WdShngl' 'Metal' 'WdShake' 'Membran' 'Tar\&Grv' 'Roll'
 'ClyTile']
  Exterior1st: ['VinylSd' 'MetalSd' 'Wd Sdng' 'HdBoard' 'BrkFace' 'WdShing'
'CemntBd'
 'Plywood' 'AsbShng' 'Stucco' 'BrkComm' 'AsphShn' 'Stone' 'ImStucc'
 'CBlock']
  Exterior2nd: ['VinylSd' 'MetalSd' 'Wd Shng' 'HdBoard' 'Plywood' 'Wd Sdng'
'CmentBd'
 'BrkFace' 'Stucco' 'AsbShng' 'Brk Cmn' 'ImStucc' 'AsphShn' 'Stone'
 'Other' 'CBlock']
  MasVnrType: ['BrkFace' 'Stone' 'BrkCmn']
  ExterQual: ['Gd' 'TA' 'Ex' 'Fa']
  ExterCond: ['TA' 'Gd' 'Fa' 'Po' 'Ex']
  Foundation: ['PConc' 'CBlock' 'BrkTil' 'Wood' 'Slab' 'Stone']
  BsmtQual: ['Gd' 'TA' 'Ex' 'Fa']
  BsmtCond: ['TA' 'Gd' 'Fa' 'Po']
  BsmtExposure: ['No' 'Gd' 'Mn' 'Av']
  BsmtFinType1: ['GLQ' 'ALQ' 'Unf' 'Rec' 'BLQ' 'LwQ']
  BsmtFinType2: ['Unf' 'BLQ' 'ALQ' 'Rec' 'LwQ' 'GLQ']
  Heating: ['GasA' 'GasW' 'Grav' 'Wall' 'OthW' 'Floor']
  HeatingQC: ['Ex' 'Gd' 'TA' 'Fa' 'Po']
  CentralAir: ['Y' 'N']
  Electrical: ['SBrkr' 'FuseF' 'FuseA' 'FuseP' 'Mix']
  KitchenQual: ['Gd' 'TA' 'Ex' 'Fa']
  Functional: ['Typ' 'Min1' 'Maj1' 'Min2' 'Mod' 'Maj2' 'Sev']
  FireplaceQu: ['Gd' 'TA' 'Fa' 'Ex' 'Po']
  GarageType: ['Attchd' 'Detchd' 'BuiltIn' 'CarPort' 'Basment' '2Types']
  GarageFinish: ['RFn' 'Unf' 'Fin']
  GarageQual: ['TA' 'Fa' 'Gd' 'Ex' 'Po']
  GarageCond: ['TA' 'Fa' 'Gd' 'Po' 'Ex']
  PavedDrive: ['Y' 'N' 'P']
  PoolQC: ['Gd' 'Ex' 'Fa']
  Fence: ['MnPrv' 'GdWo' 'GdPrv' 'MnWw']
  MiscFeature: ['Shed' 'Gar2' 'Othr' 'TenC']
  SaleType: ['WD' 'New' 'COD' 'ConLD' 'ConLI' 'CWD' 'ConLw' 'Con' 'Oth']
  SaleCondition: ['Normal' 'Abnorml' 'Partial' 'AdjLand' 'Alloca' 'Family']
    \end{Verbatim}

    \begin{tcolorbox}[breakable, size=fbox, boxrule=1pt, pad at break*=1mm,colback=cellbackground, colframe=cellborder]
\prompt{In}{incolor}{9}{\boxspacing}
\begin{Verbatim}[commandchars=\\\{\}]
\PY{c+c1}{\PYZsh{} Codificación de variables categóricas}
\PY{c+c1}{\PYZsh{} Las que son si/no (binarias) las pasamos a 0/1 con LabelEncoder}
\PY{n}{binary\PYZus{}cols} \PY{o}{=} \PY{p}{[}\PY{n}{col} \PY{k}{for} \PY{n}{col} \PY{o+ow}{in} \PY{n}{cat\PYZus{}cols} \PY{k}{if} \PY{n}{df\PYZus{}reg}\PY{p}{[}\PY{n}{col}\PY{p}{]}\PY{o}{.}\PY{n}{nunique}\PY{p}{(}\PY{p}{)} \PY{o}{==} \PY{l+m+mi}{2}\PY{p}{]}
\PY{n+nb}{print}\PY{p}{(}\PY{l+s+sa}{f}\PY{l+s+s1}{\PYZsq{}}\PY{l+s+s1}{Variables binarias: }\PY{l+s+si}{\PYZob{}}\PY{n}{binary\PYZus{}cols}\PY{l+s+si}{\PYZcb{}}\PY{l+s+s1}{\PYZsq{}}\PY{p}{)}

\PY{n}{le} \PY{o}{=} \PY{n}{LabelEncoder}\PY{p}{(}\PY{p}{)}
\PY{k}{for} \PY{n}{col} \PY{o+ow}{in} \PY{n}{binary\PYZus{}cols}\PY{p}{:}
    \PY{n}{df\PYZus{}reg}\PY{p}{[}\PY{n}{col}\PY{p}{]} \PY{o}{=} \PY{n}{le}\PY{o}{.}\PY{n}{fit\PYZus{}transform}\PY{p}{(}\PY{n}{df\PYZus{}reg}\PY{p}{[}\PY{n}{col}\PY{p}{]}\PY{p}{)}

\PY{c+c1}{\PYZsh{} Las que tienen mas de 2 categorias las convertimos con One\PYZhy{}Hot Encoding}
\PY{c+c1}{\PYZsh{} Usamos drop\PYZus{}first=True para evitar la trampa de las variables dummy}
\PY{c+c1}{\PYZsh{} (si tenemos n categorias, con n\PYZhy{}1 columnas ya queda definida la informacion)}
\PY{n}{multi\PYZus{}cat\PYZus{}cols} \PY{o}{=} \PY{p}{[}\PY{n}{col} \PY{k}{for} \PY{n}{col} \PY{o+ow}{in} \PY{n}{cat\PYZus{}cols} \PY{k}{if} \PY{n}{col} \PY{o+ow}{not} \PY{o+ow}{in} \PY{n}{binary\PYZus{}cols}\PY{p}{]}
\PY{n+nb}{print}\PY{p}{(}\PY{l+s+sa}{f}\PY{l+s+s1}{\PYZsq{}}\PY{l+s+s1}{Variables multi\PYZhy{}categoría: }\PY{l+s+si}{\PYZob{}}\PY{n}{multi\PYZus{}cat\PYZus{}cols}\PY{l+s+si}{\PYZcb{}}\PY{l+s+s1}{\PYZsq{}}\PY{p}{)}

\PY{n}{df\PYZus{}reg} \PY{o}{=} \PY{n}{pd}\PY{o}{.}\PY{n}{get\PYZus{}dummies}\PY{p}{(}\PY{n}{df\PYZus{}reg}\PY{p}{,} \PY{n}{columns}\PY{o}{=}\PY{n}{multi\PYZus{}cat\PYZus{}cols}\PY{p}{,} \PY{n}{drop\PYZus{}first}\PY{o}{=}\PY{k+kc}{True}\PY{p}{,} \PY{n}{dtype}\PY{o}{=}\PY{n+nb}{int}\PY{p}{)}

\PY{n+nb}{print}\PY{p}{(}\PY{l+s+sa}{f}\PY{l+s+s1}{\PYZsq{}}\PY{l+s+se}{\PYZbs{}n}\PY{l+s+s1}{Dimensiones tras codificación: }\PY{l+s+si}{\PYZob{}}\PY{n}{df\PYZus{}reg}\PY{o}{.}\PY{n}{shape}\PY{l+s+si}{\PYZcb{}}\PY{l+s+s1}{\PYZsq{}}\PY{p}{)}
\PY{n}{df\PYZus{}reg}\PY{o}{.}\PY{n}{head}\PY{p}{(}\PY{p}{)}
\end{Verbatim}
\end{tcolorbox}

    \begin{Verbatim}[commandchars=\\\{\}]
Variables binarias: ['Street', 'Alley', 'Utilities', 'CentralAir']
Variables multi-categoría: ['MSZoning', 'LotShape', 'LandContour', 'LotConfig',
'LandSlope', 'Neighborhood', 'Condition1', 'Condition2', 'BldgType',
'HouseStyle', 'RoofStyle', 'RoofMatl', 'Exterior1st', 'Exterior2nd',
'MasVnrType', 'ExterQual', 'ExterCond', 'Foundation', 'BsmtQual', 'BsmtCond',
'BsmtExposure', 'BsmtFinType1', 'BsmtFinType2', 'Heating', 'HeatingQC',
'Electrical', 'KitchenQual', 'Functional', 'FireplaceQu', 'GarageType',
'GarageFinish', 'GarageQual', 'GarageCond', 'PavedDrive', 'PoolQC', 'Fence',
'MiscFeature', 'SaleType', 'SaleCondition']

Dimensiones tras codificación: (1460, 246)
    \end{Verbatim}

            \begin{tcolorbox}[breakable, size=fbox, boxrule=.5pt, pad at break*=1mm, opacityfill=0]
\prompt{Out}{outcolor}{9}{\boxspacing}
\begin{Verbatim}[commandchars=\\\{\}]
   Id  MSSubClass  LotFrontage  LotArea  Street  Alley  Utilities  \textbackslash{}
0   1          60         65.0     8450       1      0          0
1   2          20         80.0     9600       1      0          0
2   3          60         68.0    11250       1      0          0
3   4          70         60.0     9550       1      0          0
4   5          60         84.0    14260       1      0          0

   OverallQual  OverallCond  YearBuilt  {\ldots}  SaleType\_ConLI  SaleType\_ConLw  \textbackslash{}
0            7            5       2003  {\ldots}               0               0
1            6            8       1976  {\ldots}               0               0
2            7            5       2001  {\ldots}               0               0
3            7            5       1915  {\ldots}               0               0
4            8            5       2000  {\ldots}               0               0

   SaleType\_New  SaleType\_Oth  SaleType\_WD  SaleCondition\_AdjLand  \textbackslash{}
0             0             0            1                      0
1             0             0            1                      0
2             0             0            1                      0
3             0             0            1                      0
4             0             0            1                      0

   SaleCondition\_Alloca  SaleCondition\_Family  SaleCondition\_Normal  \textbackslash{}
0                     0                     0                     1
1                     0                     0                     1
2                     0                     0                     1
3                     0                     0                     0
4                     0                     0                     1

   SaleCondition\_Partial
0                      0
1                      0
2                      0
3                      0
4                      0

[5 rows x 246 columns]
\end{Verbatim}
\end{tcolorbox}
        
    \subsubsection{1.1.4 Separar datos en conjuntos de entrenamiento y
test}\label{separar-datos-en-conjuntos-de-entrenamiento-y-test}

    \begin{tcolorbox}[breakable, size=fbox, boxrule=1pt, pad at break*=1mm,colback=cellbackground, colframe=cellborder]
\prompt{In}{incolor}{10}{\boxspacing}
\begin{Verbatim}[commandchars=\\\{\}]
\PY{c+c1}{\PYZsh{} Quitamos la columna Id que es solo un identificador}
\PY{k}{if} \PY{l+s+s1}{\PYZsq{}}\PY{l+s+s1}{Id}\PY{l+s+s1}{\PYZsq{}} \PY{o+ow}{in} \PY{n}{df\PYZus{}reg}\PY{o}{.}\PY{n}{columns}\PY{p}{:}
    \PY{n}{df\PYZus{}reg} \PY{o}{=} \PY{n}{df\PYZus{}reg}\PY{o}{.}\PY{n}{drop}\PY{p}{(}\PY{l+s+s1}{\PYZsq{}}\PY{l+s+s1}{Id}\PY{l+s+s1}{\PYZsq{}}\PY{p}{,} \PY{n}{axis}\PY{o}{=}\PY{l+m+mi}{1}\PY{p}{)}

\PY{c+c1}{\PYZsh{} Separamos features (X) y target (y)}
\PY{n}{X\PYZus{}reg} \PY{o}{=} \PY{n}{df\PYZus{}reg}\PY{o}{.}\PY{n}{drop}\PY{p}{(}\PY{l+s+s1}{\PYZsq{}}\PY{l+s+s1}{SalePrice}\PY{l+s+s1}{\PYZsq{}}\PY{p}{,} \PY{n}{axis}\PY{o}{=}\PY{l+m+mi}{1}\PY{p}{)}
\PY{n}{y\PYZus{}reg} \PY{o}{=} \PY{n}{df\PYZus{}reg}\PY{p}{[}\PY{l+s+s1}{\PYZsq{}}\PY{l+s+s1}{SalePrice}\PY{l+s+s1}{\PYZsq{}}\PY{p}{]}

\PY{c+c1}{\PYZsh{} Dividimos en train (80\PYZpc{}) y test (20\PYZpc{})}
\PY{n}{X\PYZus{}train\PYZus{}reg}\PY{p}{,} \PY{n}{X\PYZus{}test\PYZus{}reg}\PY{p}{,} \PY{n}{y\PYZus{}train\PYZus{}reg}\PY{p}{,} \PY{n}{y\PYZus{}test\PYZus{}reg} \PY{o}{=} \PY{n}{train\PYZus{}test\PYZus{}split}\PY{p}{(}
    \PY{n}{X\PYZus{}reg}\PY{p}{,} \PY{n}{y\PYZus{}reg}\PY{p}{,} \PY{n}{test\PYZus{}size}\PY{o}{=}\PY{l+m+mf}{0.2}\PY{p}{,} \PY{n}{random\PYZus{}state}\PY{o}{=}\PY{l+m+mi}{42}
\PY{p}{)}

\PY{n+nb}{print}\PY{p}{(}\PY{l+s+sa}{f}\PY{l+s+s1}{\PYZsq{}}\PY{l+s+s1}{Train: }\PY{l+s+si}{\PYZob{}}\PY{n}{X\PYZus{}train\PYZus{}reg}\PY{o}{.}\PY{n}{shape}\PY{p}{[}\PY{l+m+mi}{0}\PY{p}{]}\PY{l+s+si}{\PYZcb{}}\PY{l+s+s1}{ muestras}\PY{l+s+s1}{\PYZsq{}}\PY{p}{)}
\PY{n+nb}{print}\PY{p}{(}\PY{l+s+sa}{f}\PY{l+s+s1}{\PYZsq{}}\PY{l+s+s1}{Test:  }\PY{l+s+si}{\PYZob{}}\PY{n}{X\PYZus{}test\PYZus{}reg}\PY{o}{.}\PY{n}{shape}\PY{p}{[}\PY{l+m+mi}{0}\PY{p}{]}\PY{l+s+si}{\PYZcb{}}\PY{l+s+s1}{ muestras}\PY{l+s+s1}{\PYZsq{}}\PY{p}{)}
\end{Verbatim}
\end{tcolorbox}

    \begin{Verbatim}[commandchars=\\\{\}]
Train: 1168 muestras
Test:  292 muestras
    \end{Verbatim}

    \subsubsection{1.1.5 Estandarizar los
datos}\label{estandarizar-los-datos}

    \begin{tcolorbox}[breakable, size=fbox, boxrule=1pt, pad at break*=1mm,colback=cellbackground, colframe=cellborder]
\prompt{In}{incolor}{11}{\boxspacing}
\begin{Verbatim}[commandchars=\\\{\}]
\PY{c+c1}{\PYZsh{} Estandarizamos las features para que todas esten en la misma escala}
\PY{c+c1}{\PYZsh{} MUY IMPORTANTE: el fit() solo sobre train, para no contaminar con datos de test (data leakage)}
\PY{n}{scaler\PYZus{}reg} \PY{o}{=} \PY{n}{StandardScaler}\PY{p}{(}\PY{p}{)}
\PY{n}{X\PYZus{}train\PYZus{}reg\PYZus{}scaled} \PY{o}{=} \PY{n}{scaler\PYZus{}reg}\PY{o}{.}\PY{n}{fit\PYZus{}transform}\PY{p}{(}\PY{n}{X\PYZus{}train\PYZus{}reg}\PY{p}{)}
\PY{n}{X\PYZus{}test\PYZus{}reg\PYZus{}scaled} \PY{o}{=} \PY{n}{scaler\PYZus{}reg}\PY{o}{.}\PY{n}{transform}\PY{p}{(}\PY{n}{X\PYZus{}test\PYZus{}reg}\PY{p}{)}

\PY{n+nb}{print}\PY{p}{(}\PY{l+s+s1}{\PYZsq{}}\PY{l+s+s1}{Datos estandarizados correctamente.}\PY{l+s+s1}{\PYZsq{}}\PY{p}{)}
\end{Verbatim}
\end{tcolorbox}

    \begin{Verbatim}[commandchars=\\\{\}]
Datos estandarizados correctamente.
    \end{Verbatim}

    \subsubsection{1.1.6 Aplicar el modelo
XGBoostRegressor}\label{aplicar-el-modelo-xgboostregressor}

En vez de usar unos hiperparámetros fijos, vamos a probar 3
configuraciones diferentes y quedarnos con la que mejor funcione. Usamos
validación cruzada (5-fold) para que el resultado no dependa de una sola
partición de datos:

{\def\LTcaptype{none} % do not increment counter
\begin{longtable}[]{@{}
  >{\raggedright\arraybackslash}p{(\linewidth - 8\tabcolsep) * \real{0.2000}}
  >{\raggedright\arraybackslash}p{(\linewidth - 8\tabcolsep) * \real{0.2000}}
  >{\raggedright\arraybackslash}p{(\linewidth - 8\tabcolsep) * \real{0.2000}}
  >{\raggedright\arraybackslash}p{(\linewidth - 8\tabcolsep) * \real{0.2000}}
  >{\raggedright\arraybackslash}p{(\linewidth - 8\tabcolsep) * \real{0.2000}}@{}}
\toprule\noalign{}
\begin{minipage}[b]{\linewidth}\raggedright
Configuración
\end{minipage} & \begin{minipage}[b]{\linewidth}\raggedright
n\_estimators
\end{minipage} & \begin{minipage}[b]{\linewidth}\raggedright
max\_depth
\end{minipage} & \begin{minipage}[b]{\linewidth}\raggedright
learning\_rate
\end{minipage} & \begin{minipage}[b]{\linewidth}\raggedright
Idea
\end{minipage} \\
\midrule\noalign{}
\endhead
\bottomrule\noalign{}
\endlastfoot
Conservadora & 100 & 3 & 0.1 & Pocos árboles poco profundos, menos
riesgo de sobreajuste \\
Equilibrada & 200 & 6 & 0.1 & Valores estándar, buen punto medio \\
Agresiva & 500 & 8 & 0.05 & Muchos árboles profundos con learning rate
bajo, más capacidad pero más lento \\
\end{longtable}
}

    \begin{tcolorbox}[breakable, size=fbox, boxrule=1pt, pad at break*=1mm,colback=cellbackground, colframe=cellborder]
\prompt{In}{incolor}{12}{\boxspacing}
\begin{Verbatim}[commandchars=\\\{\}]
\PY{c+c1}{\PYZsh{} Primero probamos varias configuraciones de hiperparametros para ver cual va mejor}
\PY{c+c1}{\PYZsh{} En vez de quedarnos con los valores por defecto, comparamos 3 configuraciones}
\PY{k+kn}{from}\PY{+w}{ }\PY{n+nn}{sklearn}\PY{n+nn}{.}\PY{n+nn}{model\PYZus{}selection}\PY{+w}{ }\PY{k+kn}{import} \PY{n}{cross\PYZus{}val\PYZus{}score}

\PY{n}{configs} \PY{o}{=} \PY{p}{\PYZob{}}
    \PY{l+s+s1}{\PYZsq{}}\PY{l+s+s1}{Conservador (pocos arboles, poca profundidad)}\PY{l+s+s1}{\PYZsq{}}\PY{p}{:} \PY{p}{\PYZob{}}\PY{l+s+s1}{\PYZsq{}}\PY{l+s+s1}{n\PYZus{}estimators}\PY{l+s+s1}{\PYZsq{}}\PY{p}{:} \PY{l+m+mi}{100}\PY{p}{,} \PY{l+s+s1}{\PYZsq{}}\PY{l+s+s1}{max\PYZus{}depth}\PY{l+s+s1}{\PYZsq{}}\PY{p}{:} \PY{l+m+mi}{3}\PY{p}{,} \PY{l+s+s1}{\PYZsq{}}\PY{l+s+s1}{learning\PYZus{}rate}\PY{l+s+s1}{\PYZsq{}}\PY{p}{:} \PY{l+m+mf}{0.1}\PY{p}{\PYZcb{}}\PY{p}{,}
    \PY{l+s+s1}{\PYZsq{}}\PY{l+s+s1}{Equilibrado (valores estandar)}\PY{l+s+s1}{\PYZsq{}}\PY{p}{:} \PY{p}{\PYZob{}}\PY{l+s+s1}{\PYZsq{}}\PY{l+s+s1}{n\PYZus{}estimators}\PY{l+s+s1}{\PYZsq{}}\PY{p}{:} \PY{l+m+mi}{200}\PY{p}{,} \PY{l+s+s1}{\PYZsq{}}\PY{l+s+s1}{max\PYZus{}depth}\PY{l+s+s1}{\PYZsq{}}\PY{p}{:} \PY{l+m+mi}{6}\PY{p}{,} \PY{l+s+s1}{\PYZsq{}}\PY{l+s+s1}{learning\PYZus{}rate}\PY{l+s+s1}{\PYZsq{}}\PY{p}{:} \PY{l+m+mf}{0.1}\PY{p}{\PYZcb{}}\PY{p}{,}
    \PY{l+s+s1}{\PYZsq{}}\PY{l+s+s1}{Agresivo (muchos arboles, mas profundidad)}\PY{l+s+s1}{\PYZsq{}}\PY{p}{:} \PY{p}{\PYZob{}}\PY{l+s+s1}{\PYZsq{}}\PY{l+s+s1}{n\PYZus{}estimators}\PY{l+s+s1}{\PYZsq{}}\PY{p}{:} \PY{l+m+mi}{500}\PY{p}{,} \PY{l+s+s1}{\PYZsq{}}\PY{l+s+s1}{max\PYZus{}depth}\PY{l+s+s1}{\PYZsq{}}\PY{p}{:} \PY{l+m+mi}{8}\PY{p}{,} \PY{l+s+s1}{\PYZsq{}}\PY{l+s+s1}{learning\PYZus{}rate}\PY{l+s+s1}{\PYZsq{}}\PY{p}{:} \PY{l+m+mf}{0.05}\PY{p}{\PYZcb{}}\PY{p}{,}
\PY{p}{\PYZcb{}}

\PY{n+nb}{print}\PY{p}{(}\PY{l+s+s1}{\PYZsq{}}\PY{l+s+s1}{=== Comparativa de hiperparámetros XGBoostRegressor ===}\PY{l+s+se}{\PYZbs{}n}\PY{l+s+s1}{\PYZsq{}}\PY{p}{)}
\PY{n}{mejor\PYZus{}config} \PY{o}{=} \PY{k+kc}{None}
\PY{n}{mejor\PYZus{}r2} \PY{o}{=} \PY{o}{\PYZhy{}}\PY{l+m+mi}{1}

\PY{k}{for} \PY{n}{nombre}\PY{p}{,} \PY{n}{params} \PY{o+ow}{in} \PY{n}{configs}\PY{o}{.}\PY{n}{items}\PY{p}{(}\PY{p}{)}\PY{p}{:}
    \PY{n}{modelo\PYZus{}tmp} \PY{o}{=} \PY{n}{XGBRegressor}\PY{p}{(}\PY{o}{*}\PY{o}{*}\PY{n}{params}\PY{p}{,} \PY{n}{random\PYZus{}state}\PY{o}{=}\PY{l+m+mi}{42}\PY{p}{)}
    \PY{c+c1}{\PYZsh{} Validacion cruzada con 5 folds para no depender de una sola particion}
    \PY{n}{scores} \PY{o}{=} \PY{n}{cross\PYZus{}val\PYZus{}score}\PY{p}{(}\PY{n}{modelo\PYZus{}tmp}\PY{p}{,} \PY{n}{X\PYZus{}train\PYZus{}reg\PYZus{}scaled}\PY{p}{,} \PY{n}{y\PYZus{}train\PYZus{}reg}\PY{p}{,} \PY{n}{cv}\PY{o}{=}\PY{l+m+mi}{5}\PY{p}{,} \PY{n}{scoring}\PY{o}{=}\PY{l+s+s1}{\PYZsq{}}\PY{l+s+s1}{r2}\PY{l+s+s1}{\PYZsq{}}\PY{p}{)}
    \PY{n}{r2\PYZus{}medio} \PY{o}{=} \PY{n}{scores}\PY{o}{.}\PY{n}{mean}\PY{p}{(}\PY{p}{)}
    \PY{n+nb}{print}\PY{p}{(}\PY{l+s+sa}{f}\PY{l+s+s1}{\PYZsq{}}\PY{l+s+si}{\PYZob{}}\PY{n}{nombre}\PY{l+s+si}{\PYZcb{}}\PY{l+s+s1}{:}\PY{l+s+s1}{\PYZsq{}}\PY{p}{)}
    \PY{n+nb}{print}\PY{p}{(}\PY{l+s+sa}{f}\PY{l+s+s1}{\PYZsq{}}\PY{l+s+s1}{  Params: n\PYZus{}estimators=}\PY{l+s+si}{\PYZob{}}\PY{n}{params}\PY{p}{[}\PY{l+s+s2}{\PYZdq{}}\PY{l+s+s2}{n\PYZus{}estimators}\PY{l+s+s2}{\PYZdq{}}\PY{p}{]}\PY{l+s+si}{\PYZcb{}}\PY{l+s+s1}{, max\PYZus{}depth=}\PY{l+s+si}{\PYZob{}}\PY{n}{params}\PY{p}{[}\PY{l+s+s2}{\PYZdq{}}\PY{l+s+s2}{max\PYZus{}depth}\PY{l+s+s2}{\PYZdq{}}\PY{p}{]}\PY{l+s+si}{\PYZcb{}}\PY{l+s+s1}{, lr=}\PY{l+s+si}{\PYZob{}}\PY{n}{params}\PY{p}{[}\PY{l+s+s2}{\PYZdq{}}\PY{l+s+s2}{learning\PYZus{}rate}\PY{l+s+s2}{\PYZdq{}}\PY{p}{]}\PY{l+s+si}{\PYZcb{}}\PY{l+s+s1}{\PYZsq{}}\PY{p}{)}
    \PY{n+nb}{print}\PY{p}{(}\PY{l+s+sa}{f}\PY{l+s+s1}{\PYZsq{}}\PY{l+s+s1}{  R² medio (CV 5\PYZhy{}fold): }\PY{l+s+si}{\PYZob{}}\PY{n}{r2\PYZus{}medio}\PY{l+s+si}{:}\PY{l+s+s1}{.4f}\PY{l+s+si}{\PYZcb{}}\PY{l+s+s1}{ (+/\PYZhy{} }\PY{l+s+si}{\PYZob{}}\PY{n}{scores}\PY{o}{.}\PY{n}{std}\PY{p}{(}\PY{p}{)}\PY{l+s+si}{:}\PY{l+s+s1}{.4f}\PY{l+s+si}{\PYZcb{}}\PY{l+s+s1}{)}\PY{l+s+se}{\PYZbs{}n}\PY{l+s+s1}{\PYZsq{}}\PY{p}{)}
    \PY{k}{if} \PY{n}{r2\PYZus{}medio} \PY{o}{\PYZgt{}} \PY{n}{mejor\PYZus{}r2}\PY{p}{:}
        \PY{n}{mejor\PYZus{}r2} \PY{o}{=} \PY{n}{r2\PYZus{}medio}
        \PY{n}{mejor\PYZus{}config} \PY{o}{=} \PY{n}{params}
        \PY{n}{mejor\PYZus{}nombre} \PY{o}{=} \PY{n}{nombre}

\PY{n+nb}{print}\PY{p}{(}\PY{l+s+sa}{f}\PY{l+s+s1}{\PYZsq{}}\PY{l+s+s1}{\PYZgt{}\PYZgt{}\PYZgt{} Mejor configuración: }\PY{l+s+si}{\PYZob{}}\PY{n}{mejor\PYZus{}nombre}\PY{l+s+si}{\PYZcb{}}\PY{l+s+s1}{\PYZsq{}}\PY{p}{)}
\PY{n+nb}{print}\PY{p}{(}\PY{l+s+sa}{f}\PY{l+s+s1}{\PYZsq{}}\PY{l+s+s1}{\PYZgt{}\PYZgt{}\PYZgt{} Params: }\PY{l+s+si}{\PYZob{}}\PY{n}{mejor\PYZus{}config}\PY{l+s+si}{\PYZcb{}}\PY{l+s+se}{\PYZbs{}n}\PY{l+s+s1}{\PYZsq{}}\PY{p}{)}

\PY{c+c1}{\PYZsh{} Entrenamos con la mejor configuracion}
\PY{n}{xgb\PYZus{}reg} \PY{o}{=} \PY{n}{XGBRegressor}\PY{p}{(}\PY{o}{*}\PY{o}{*}\PY{n}{mejor\PYZus{}config}\PY{p}{,} \PY{n}{random\PYZus{}state}\PY{o}{=}\PY{l+m+mi}{42}\PY{p}{)}
\PY{n}{xgb\PYZus{}reg}\PY{o}{.}\PY{n}{fit}\PY{p}{(}\PY{n}{X\PYZus{}train\PYZus{}reg\PYZus{}scaled}\PY{p}{,} \PY{n}{y\PYZus{}train\PYZus{}reg}\PY{p}{)}
\PY{n}{y\PYZus{}pred\PYZus{}reg} \PY{o}{=} \PY{n}{xgb\PYZus{}reg}\PY{o}{.}\PY{n}{predict}\PY{p}{(}\PY{n}{X\PYZus{}test\PYZus{}reg\PYZus{}scaled}\PY{p}{)}

\PY{n+nb}{print}\PY{p}{(}\PY{l+s+sa}{f}\PY{l+s+s1}{\PYZsq{}}\PY{l+s+s1}{Modelo XGBoostRegressor entrenado con la mejor configuración.}\PY{l+s+s1}{\PYZsq{}}\PY{p}{)}
\end{Verbatim}
\end{tcolorbox}

    \begin{Verbatim}[commandchars=\\\{\}]
=== Comparativa de hiperparámetros XGBoostRegressor ===

Conservador (pocos arboles, poca profundidad):
  Params: n\_estimators=100, max\_depth=3, lr=0.1
  R² medio (CV 5-fold): 0.8519 (+/- 0.0511)

Equilibrado (valores estandar):
  Params: n\_estimators=200, max\_depth=6, lr=0.1
  R² medio (CV 5-fold): 0.8352 (+/- 0.0556)

Agresivo (muchos arboles, mas profundidad):
  Params: n\_estimators=500, max\_depth=8, lr=0.05
  R² medio (CV 5-fold): 0.8278 (+/- 0.0408)

>>> Mejor configuración: Conservador (pocos arboles, poca profundidad)
>>> Params: \{'n\_estimators': 100, 'max\_depth': 3, 'learning\_rate': 0.1\}

Modelo XGBoostRegressor entrenado con la mejor configuración.
    \end{Verbatim}

    \subsubsection{1.1.7 Aplicar métricas (RMSE) y comparar con
UT5}\label{aplicar-muxe9tricas-rmse-y-comparar-con-ut5}

Evaluamos el modelo con tres métricas: - \textbf{RMSE:} El error medio
en las mismas unidades que los precios (dólares). Nos dice cuánto se
equivoca de media al predecir. - \textbf{R²:} Qué porcentaje de la
variación en los precios explica el modelo (de 0 a 1, cuanto más alto
mejor). - \textbf{MSE:} El RMSE al cuadrado (se usa menos para
interpretar, pero los algoritmos lo usan internamente).

    \begin{tcolorbox}[breakable, size=fbox, boxrule=1pt, pad at break*=1mm,colback=cellbackground, colframe=cellborder]
\prompt{In}{incolor}{13}{\boxspacing}
\begin{Verbatim}[commandchars=\\\{\}]
\PY{c+c1}{\PYZsh{} Calculamos métricas}
\PY{n}{mse\PYZus{}xgb} \PY{o}{=} \PY{n}{mean\PYZus{}squared\PYZus{}error}\PY{p}{(}\PY{n}{y\PYZus{}test\PYZus{}reg}\PY{p}{,} \PY{n}{y\PYZus{}pred\PYZus{}reg}\PY{p}{)}
\PY{n}{rmse\PYZus{}xgb} \PY{o}{=} \PY{n}{np}\PY{o}{.}\PY{n}{sqrt}\PY{p}{(}\PY{n}{mse\PYZus{}xgb}\PY{p}{)}
\PY{n}{r2\PYZus{}xgb} \PY{o}{=} \PY{n}{r2\PYZus{}score}\PY{p}{(}\PY{n}{y\PYZus{}test\PYZus{}reg}\PY{p}{,} \PY{n}{y\PYZus{}pred\PYZus{}reg}\PY{p}{)}

\PY{n+nb}{print}\PY{p}{(}\PY{l+s+s1}{\PYZsq{}}\PY{l+s+s1}{=== Métricas XGBoostRegressor ===}\PY{l+s+s1}{\PYZsq{}}\PY{p}{)}
\PY{n+nb}{print}\PY{p}{(}\PY{l+s+sa}{f}\PY{l+s+s1}{\PYZsq{}}\PY{l+s+s1}{MSE:  }\PY{l+s+si}{\PYZob{}}\PY{n}{mse\PYZus{}xgb}\PY{l+s+si}{:}\PY{l+s+s1}{,.2f}\PY{l+s+si}{\PYZcb{}}\PY{l+s+s1}{\PYZsq{}}\PY{p}{)}
\PY{n+nb}{print}\PY{p}{(}\PY{l+s+sa}{f}\PY{l+s+s1}{\PYZsq{}}\PY{l+s+s1}{RMSE: }\PY{l+s+si}{\PYZob{}}\PY{n}{rmse\PYZus{}xgb}\PY{l+s+si}{:}\PY{l+s+s1}{,.2f}\PY{l+s+si}{\PYZcb{}}\PY{l+s+s1}{\PYZsq{}}\PY{p}{)}
\PY{n+nb}{print}\PY{p}{(}\PY{l+s+sa}{f}\PY{l+s+s1}{\PYZsq{}}\PY{l+s+s1}{R²:   }\PY{l+s+si}{\PYZob{}}\PY{n}{r2\PYZus{}xgb}\PY{l+s+si}{:}\PY{l+s+s1}{.4f}\PY{l+s+si}{\PYZcb{}}\PY{l+s+s1}{\PYZsq{}}\PY{p}{)}
\end{Verbatim}
\end{tcolorbox}

    \begin{Verbatim}[commandchars=\\\{\}]
=== Métricas XGBoostRegressor ===
MSE:  703,294,400.00
RMSE: 26,519.70
R²:   0.9083
    \end{Verbatim}

    \begin{tcolorbox}[breakable, size=fbox, boxrule=1pt, pad at break*=1mm,colback=cellbackground, colframe=cellborder]
\prompt{In}{incolor}{14}{\boxspacing}
\begin{Verbatim}[commandchars=\\\{\}]
\PY{c+c1}{\PYZsh{} Métricas del modelo XGBoost}
\PY{n+nb}{print}\PY{p}{(}\PY{l+s+s1}{\PYZsq{}}\PY{l+s+s1}{=== Métricas XGBoostRegressor ===}\PY{l+s+s1}{\PYZsq{}}\PY{p}{)}
\PY{n+nb}{print}\PY{p}{(}\PY{l+s+sa}{f}\PY{l+s+s1}{\PYZsq{}}\PY{l+s+s1}{MSE:  }\PY{l+s+si}{\PYZob{}}\PY{n}{mse\PYZus{}xgb}\PY{l+s+si}{:}\PY{l+s+s1}{,.2f}\PY{l+s+si}{\PYZcb{}}\PY{l+s+s1}{\PYZsq{}}\PY{p}{)}
\PY{n+nb}{print}\PY{p}{(}\PY{l+s+sa}{f}\PY{l+s+s1}{\PYZsq{}}\PY{l+s+s1}{RMSE: }\PY{l+s+si}{\PYZob{}}\PY{n}{rmse\PYZus{}xgb}\PY{l+s+si}{:}\PY{l+s+s1}{,.2f}\PY{l+s+si}{\PYZcb{}}\PY{l+s+s1}{\PYZsq{}}\PY{p}{)}
\PY{n+nb}{print}\PY{p}{(}\PY{l+s+sa}{f}\PY{l+s+s1}{\PYZsq{}}\PY{l+s+s1}{R²:   }\PY{l+s+si}{\PYZob{}}\PY{n}{r2\PYZus{}xgb}\PY{l+s+si}{:}\PY{l+s+s1}{.4f}\PY{l+s+si}{\PYZcb{}}\PY{l+s+s1}{\PYZsq{}}\PY{p}{)}

\PY{c+c1}{\PYZsh{} Nota: El dataset de Kaggle (Housing.csv, 545 filas) difiere del usado en UT5}
\PY{c+c1}{\PYZsh{} (house\PYZus{}prices\PYZus{}train.csv, 1460 filas), por lo que la comparación directa de RMSE}
\PY{c+c1}{\PYZsh{} no es válida. Comparamos en términos de R² como referencia relativa.}

\PY{n+nb}{print}\PY{p}{(}\PY{l+s+s1}{\PYZsq{}}\PY{l+s+se}{\PYZbs{}n}\PY{l+s+s1}{=== Referencia modelos UT5 (dataset diferente) ===}\PY{l+s+s1}{\PYZsq{}}\PY{p}{)}
\PY{n+nb}{print}\PY{p}{(}\PY{l+s+s1}{\PYZsq{}}\PY{l+s+s1}{Linear Regression: R²=0.8544 | SVR: R²=0.8757}\PY{l+s+s1}{\PYZsq{}}\PY{p}{)}
\PY{n+nb}{print}\PY{p}{(}\PY{l+s+s1}{\PYZsq{}}\PY{l+s+s1}{Decision Tree: R²=0.8275 | Random Forest: R²=0.8977}\PY{l+s+s1}{\PYZsq{}}\PY{p}{)}
\PY{n+nb}{print}\PY{p}{(}\PY{l+s+s1}{\PYZsq{}}\PY{l+s+s1}{Ensemble: R²=0.9039}\PY{l+s+s1}{\PYZsq{}}\PY{p}{)}
\PY{n+nb}{print}\PY{p}{(}\PY{l+s+sa}{f}\PY{l+s+s1}{\PYZsq{}}\PY{l+s+s1}{XGBoost (este trabajo): R²=}\PY{l+s+si}{\PYZob{}}\PY{n}{r2\PYZus{}xgb}\PY{l+s+si}{:}\PY{l+s+s1}{.4f}\PY{l+s+si}{\PYZcb{}}\PY{l+s+s1}{\PYZsq{}}\PY{p}{)}

\PY{c+c1}{\PYZsh{} Gráfico de predicción vs real}
\PY{n}{plt}\PY{o}{.}\PY{n}{figure}\PY{p}{(}\PY{n}{figsize}\PY{o}{=}\PY{p}{(}\PY{l+m+mi}{8}\PY{p}{,} \PY{l+m+mi}{6}\PY{p}{)}\PY{p}{)}
\PY{n}{plt}\PY{o}{.}\PY{n}{scatter}\PY{p}{(}\PY{n}{y\PYZus{}test\PYZus{}reg}\PY{p}{,} \PY{n}{y\PYZus{}pred\PYZus{}reg}\PY{p}{,} \PY{n}{alpha}\PY{o}{=}\PY{l+m+mf}{0.6}\PY{p}{,} \PY{n}{edgecolors}\PY{o}{=}\PY{l+s+s1}{\PYZsq{}}\PY{l+s+s1}{k}\PY{l+s+s1}{\PYZsq{}}\PY{p}{,} \PY{n}{linewidths}\PY{o}{=}\PY{l+m+mf}{0.5}\PY{p}{)}
\PY{n}{plt}\PY{o}{.}\PY{n}{plot}\PY{p}{(}\PY{p}{[}\PY{n}{y\PYZus{}test\PYZus{}reg}\PY{o}{.}\PY{n}{min}\PY{p}{(}\PY{p}{)}\PY{p}{,} \PY{n}{y\PYZus{}test\PYZus{}reg}\PY{o}{.}\PY{n}{max}\PY{p}{(}\PY{p}{)}\PY{p}{]}\PY{p}{,} \PY{p}{[}\PY{n}{y\PYZus{}test\PYZus{}reg}\PY{o}{.}\PY{n}{min}\PY{p}{(}\PY{p}{)}\PY{p}{,} \PY{n}{y\PYZus{}test\PYZus{}reg}\PY{o}{.}\PY{n}{max}\PY{p}{(}\PY{p}{)}\PY{p}{]}\PY{p}{,} \PY{l+s+s1}{\PYZsq{}}\PY{l+s+s1}{r\PYZhy{}\PYZhy{}}\PY{l+s+s1}{\PYZsq{}}\PY{p}{,} \PY{n}{lw}\PY{o}{=}\PY{l+m+mi}{2}\PY{p}{)}
\PY{n}{plt}\PY{o}{.}\PY{n}{xlabel}\PY{p}{(}\PY{l+s+s1}{\PYZsq{}}\PY{l+s+s1}{Precio Real}\PY{l+s+s1}{\PYZsq{}}\PY{p}{)}
\PY{n}{plt}\PY{o}{.}\PY{n}{ylabel}\PY{p}{(}\PY{l+s+s1}{\PYZsq{}}\PY{l+s+s1}{Precio Predicho}\PY{l+s+s1}{\PYZsq{}}\PY{p}{)}
\PY{n}{plt}\PY{o}{.}\PY{n}{title}\PY{p}{(}\PY{l+s+sa}{f}\PY{l+s+s1}{\PYZsq{}}\PY{l+s+s1}{XGBoostRegressor: Predicción vs Real (R²=}\PY{l+s+si}{\PYZob{}}\PY{n}{r2\PYZus{}xgb}\PY{l+s+si}{:}\PY{l+s+s1}{.4f}\PY{l+s+si}{\PYZcb{}}\PY{l+s+s1}{)}\PY{l+s+s1}{\PYZsq{}}\PY{p}{)}
\PY{n}{plt}\PY{o}{.}\PY{n}{tight\PYZus{}layout}\PY{p}{(}\PY{p}{)}
\PY{n}{plt}\PY{o}{.}\PY{n}{show}\PY{p}{(}\PY{p}{)}
\end{Verbatim}
\end{tcolorbox}

    \begin{Verbatim}[commandchars=\\\{\}]
=== Métricas XGBoostRegressor ===
MSE:  703,294,400.00
RMSE: 26,519.70
R²:   0.9083

=== Referencia modelos UT5 (dataset diferente) ===
Linear Regression: R²=0.8544 | SVR: R²=0.8757
Decision Tree: R²=0.8275 | Random Forest: R²=0.8977
Ensemble: R²=0.9039
XGBoost (este trabajo): R²=0.9083
    \end{Verbatim}

    \begin{center}
    \adjustimage{max size={0.9\linewidth}{0.9\paperheight}}{UT9_autor_XXXXXXXXX_files/UT9_autor_XXXXXXXXX_24_1.png}
    \end{center}
    { \hspace*{\fill} \\}
    
    \textbf{Conclusión regresión --- ¿qué nos dicen estos números?}

Usamos exactamente el mismo dataset que en la UT5 (1460 casas, 81
variables), así que la comparación es directa y justa.

\textbf{Comparación con la UT5:}

{\def\LTcaptype{none} % do not increment counter
\begin{longtable}[]{@{}ll@{}}
\toprule\noalign{}
Modelo (UT5) & R² \\
\midrule\noalign{}
\endhead
\bottomrule\noalign{}
\endlastfoot
Linear Regression & 0.8544 \\
SVR & 0.8757 \\
Decision Tree & 0.8275 \\
Random Forest & 0.8977 \\
Ensemble (Stacking) & 0.9039 \\
\textbf{XGBoost (este trabajo)} & \textbf{ver arriba} \\
\end{longtable}
}

XGBoost debería dar un R² competitivo con el ensemble de la UT5, ya que
gradient boosting es precisamente lo que hace XGBoost: combinar muchos
árboles donde cada uno corrige los errores del anterior.

\textbf{En el gráfico de Predicción vs Real:} los puntos que caen sobre
la línea roja son predicciones perfectas. Los que se alejan son errores.
En las casas más caras el modelo se desvía más, porque suelen tener
características únicas difíciles de predecir.

\textbf{Sobre la comparativa de hiperparámetros:} Hemos probado 3
configuraciones con validación cruzada y nos hemos quedado con la mejor.
La configuración conservadora (pocos árboles, poca profundidad) suele
dar peor resultado porque el modelo se queda corto. La agresiva (muchos
árboles, mucha profundidad) puede sobreajustar. La equilibrada suele
funcionar bien como punto medio.

    \begin{center}\rule{0.5\linewidth}{0.5pt}\end{center}

\subsection{1.2 XGBoost para
Clasificación}\label{xgboost-para-clasificaciuxf3n}

Utilizamos el dataset \textbf{Social Network Ads} de Kaggle (Rakesh
Rau). El objetivo es predecir si un usuario comprará un producto
(\texttt{Purchased}) a partir de su edad (\texttt{Age}) y salario
estimado (\texttt{EstimatedSalary}).

\textbf{Nota:} Este es el mismo dataset utilizado en la UT6, donde fue
proporcionado como ``Ensayo de Motores'' con las columnas renombradas
(Age→RPM, EstimatedSalary→Par de carga, Purchased→Fallo Motor). Aquí
usamos los nombres originales y lo descargamos directamente de Kaggle.

    \subsubsection{1.2.1 Importar dataset}\label{importar-dataset}

    \begin{tcolorbox}[breakable, size=fbox, boxrule=1pt, pad at break*=1mm,colback=cellbackground, colframe=cellborder]
\prompt{In}{incolor}{15}{\boxspacing}
\begin{Verbatim}[commandchars=\\\{\}]
\PY{c+c1}{\PYZsh{} Descargamos el dataset Social Network Ads desde Kaggle}
\PY{c+c1}{\PYZsh{} Es el mismo que usamos en la UT6 (alli se llamaba \PYZdq{}Ensayo de Motores\PYZdq{})}
\PY{n}{path\PYZus{}clf} \PY{o}{=} \PY{n}{kagglehub}\PY{o}{.}\PY{n}{dataset\PYZus{}download}\PY{p}{(}\PY{l+s+s1}{\PYZsq{}}\PY{l+s+s1}{rakeshrau/social\PYZhy{}network\PYZhy{}ads}\PY{l+s+s1}{\PYZsq{}}\PY{p}{)}
\PY{n+nb}{print}\PY{p}{(}\PY{l+s+s1}{\PYZsq{}}\PY{l+s+s1}{Dataset descargado en:}\PY{l+s+s1}{\PYZsq{}}\PY{p}{,} \PY{n}{path\PYZus{}clf}\PY{p}{)}

\PY{k}{for} \PY{n}{f} \PY{o+ow}{in} \PY{n}{os}\PY{o}{.}\PY{n}{listdir}\PY{p}{(}\PY{n}{path\PYZus{}clf}\PY{p}{)}\PY{p}{:}
    \PY{n+nb}{print}\PY{p}{(}\PY{n}{f}\PY{p}{)}
\end{Verbatim}
\end{tcolorbox}

    \begin{Verbatim}[commandchars=\\\{\}]
Downloading to /root/.cache/kagglehub/datasets/rakeshrau/social-network-
ads/1.archive{\ldots}
    \end{Verbatim}

    \begin{Verbatim}[commandchars=\\\{\}]
100\%|██████████| 3.27k/3.27k [00:00<00:00, 3.72MB/s]
    \end{Verbatim}

    \begin{Verbatim}[commandchars=\\\{\}]
Extracting files{\ldots}
Dataset descargado en: /root/.cache/kagglehub/datasets/rakeshrau/social-network-
ads/versions/1
Social\_Network\_Ads.csv
    \end{Verbatim}

    \begin{Verbatim}[commandchars=\\\{\}]

    \end{Verbatim}

    \begin{tcolorbox}[breakable, size=fbox, boxrule=1pt, pad at break*=1mm,colback=cellbackground, colframe=cellborder]
\prompt{In}{incolor}{16}{\boxspacing}
\begin{Verbatim}[commandchars=\\\{\}]
\PY{c+c1}{\PYZsh{} Cargamos el dataset}
\PY{n}{df\PYZus{}clf} \PY{o}{=} \PY{n}{pd}\PY{o}{.}\PY{n}{read\PYZus{}csv}\PY{p}{(}\PY{n}{os}\PY{o}{.}\PY{n}{path}\PY{o}{.}\PY{n}{join}\PY{p}{(}\PY{n}{path\PYZus{}clf}\PY{p}{,} \PY{l+s+s1}{\PYZsq{}}\PY{l+s+s1}{Social\PYZus{}Network\PYZus{}Ads.csv}\PY{l+s+s1}{\PYZsq{}}\PY{p}{)}\PY{p}{)}
\PY{n+nb}{print}\PY{p}{(}\PY{l+s+sa}{f}\PY{l+s+s1}{\PYZsq{}}\PY{l+s+s1}{Dimensiones del dataset: }\PY{l+s+si}{\PYZob{}}\PY{n}{df\PYZus{}clf}\PY{o}{.}\PY{n}{shape}\PY{l+s+si}{\PYZcb{}}\PY{l+s+s1}{\PYZsq{}}\PY{p}{)}
\PY{n}{df\PYZus{}clf}\PY{o}{.}\PY{n}{head}\PY{p}{(}\PY{p}{)}
\end{Verbatim}
\end{tcolorbox}

    \begin{Verbatim}[commandchars=\\\{\}]
Dimensiones del dataset: (400, 5)
    \end{Verbatim}

            \begin{tcolorbox}[breakable, size=fbox, boxrule=.5pt, pad at break*=1mm, opacityfill=0]
\prompt{Out}{outcolor}{16}{\boxspacing}
\begin{Verbatim}[commandchars=\\\{\}]
    User ID  Gender  Age  EstimatedSalary  Purchased
0  15624510    Male   19            19000          0
1  15810944    Male   35            20000          0
2  15668575  Female   26            43000          0
3  15603246  Female   27            57000          0
4  15804002    Male   19            76000          0
\end{Verbatim}
\end{tcolorbox}
        
    \begin{tcolorbox}[breakable, size=fbox, boxrule=1pt, pad at break*=1mm,colback=cellbackground, colframe=cellborder]
\prompt{In}{incolor}{17}{\boxspacing}
\begin{Verbatim}[commandchars=\\\{\}]
\PY{c+c1}{\PYZsh{} Información del dataset}
\PY{n}{df\PYZus{}clf}\PY{o}{.}\PY{n}{info}\PY{p}{(}\PY{p}{)}
\PY{n+nb}{print}\PY{p}{(}\PY{l+s+s1}{\PYZsq{}}\PY{l+s+se}{\PYZbs{}n}\PY{l+s+s1}{\PYZhy{}\PYZhy{}\PYZhy{} Estadísticas descriptivas \PYZhy{}\PYZhy{}\PYZhy{}}\PY{l+s+s1}{\PYZsq{}}\PY{p}{)}
\PY{n}{display}\PY{p}{(}\PY{n}{df\PYZus{}clf}\PY{o}{.}\PY{n}{describe}\PY{p}{(}\PY{p}{)}\PY{p}{)}
\PY{n+nb}{print}\PY{p}{(}\PY{l+s+sa}{f}\PY{l+s+s1}{\PYZsq{}}\PY{l+s+se}{\PYZbs{}n}\PY{l+s+s1}{Distribución de clases:}\PY{l+s+se}{\PYZbs{}n}\PY{l+s+si}{\PYZob{}}\PY{n}{df\PYZus{}clf}\PY{p}{[}\PY{l+s+s2}{\PYZdq{}}\PY{l+s+s2}{Purchased}\PY{l+s+s2}{\PYZdq{}}\PY{p}{]}\PY{o}{.}\PY{n}{value\PYZus{}counts}\PY{p}{(}\PY{p}{)}\PY{l+s+si}{\PYZcb{}}\PY{l+s+s1}{\PYZsq{}}\PY{p}{)}
\end{Verbatim}
\end{tcolorbox}

    \begin{Verbatim}[commandchars=\\\{\}]
<class 'pandas.core.frame.DataFrame'>
RangeIndex: 400 entries, 0 to 399
Data columns (total 5 columns):
 \#   Column           Non-Null Count  Dtype
---  ------           --------------  -----
 0   User ID          400 non-null    int64
 1   Gender           400 non-null    object
 2   Age              400 non-null    int64
 3   EstimatedSalary  400 non-null    int64
 4   Purchased        400 non-null    int64
dtypes: int64(4), object(1)
memory usage: 15.8+ KB

--- Estadísticas descriptivas ---
    \end{Verbatim}

    
    \begin{Verbatim}[commandchars=\\\{\}]
            User ID         Age  EstimatedSalary   Purchased
count  4.000000e+02  400.000000       400.000000  400.000000
mean   1.569154e+07   37.655000     69742.500000    0.357500
std    7.165832e+04   10.482877     34096.960282    0.479864
min    1.556669e+07   18.000000     15000.000000    0.000000
25\%    1.562676e+07   29.750000     43000.000000    0.000000
50\%    1.569434e+07   37.000000     70000.000000    0.000000
75\%    1.575036e+07   46.000000     88000.000000    1.000000
max    1.581524e+07   60.000000    150000.000000    1.000000
    \end{Verbatim}

    
    \begin{Verbatim}[commandchars=\\\{\}]

Distribución de clases:
Purchased
0    257
1    143
Name: count, dtype: int64
    \end{Verbatim}

    \subsubsection{1.2.2 Manejar datos missing}\label{manejar-datos-missing}

    \begin{tcolorbox}[breakable, size=fbox, boxrule=1pt, pad at break*=1mm,colback=cellbackground, colframe=cellborder]
\prompt{In}{incolor}{18}{\boxspacing}
\begin{Verbatim}[commandchars=\\\{\}]
\PY{c+c1}{\PYZsh{} Comprobamos valores nulos}
\PY{n+nb}{print}\PY{p}{(}\PY{l+s+s1}{\PYZsq{}}\PY{l+s+s1}{Valores nulos por columna:}\PY{l+s+s1}{\PYZsq{}}\PY{p}{)}
\PY{n+nb}{print}\PY{p}{(}\PY{n}{df\PYZus{}clf}\PY{o}{.}\PY{n}{isnull}\PY{p}{(}\PY{p}{)}\PY{o}{.}\PY{n}{sum}\PY{p}{(}\PY{p}{)}\PY{p}{)}
\PY{n+nb}{print}\PY{p}{(}\PY{l+s+sa}{f}\PY{l+s+s1}{\PYZsq{}}\PY{l+s+se}{\PYZbs{}n}\PY{l+s+s1}{Total de valores nulos: }\PY{l+s+si}{\PYZob{}}\PY{n}{df\PYZus{}clf}\PY{o}{.}\PY{n}{isnull}\PY{p}{(}\PY{p}{)}\PY{o}{.}\PY{n}{sum}\PY{p}{(}\PY{p}{)}\PY{o}{.}\PY{n}{sum}\PY{p}{(}\PY{p}{)}\PY{l+s+si}{\PYZcb{}}\PY{l+s+s1}{\PYZsq{}}\PY{p}{)}
\end{Verbatim}
\end{tcolorbox}

    \begin{Verbatim}[commandchars=\\\{\}]
Valores nulos por columna:
User ID            0
Gender             0
Age                0
EstimatedSalary    0
Purchased          0
dtype: int64

Total de valores nulos: 0
    \end{Verbatim}

    \subsubsection{1.2.3 Manejar datos
categóricos}\label{manejar-datos-categuxf3ricos}

El dataset contiene \texttt{User\ ID} (identificador) y \texttt{Gender}
(categórica). Eliminamos ambas: - \texttt{User\ ID}: es un identificador
sin valor predictivo - \texttt{Gender}: la eliminamos para mantener solo
Age y EstimatedSalary como features (coherente con el análisis de la
UT6)

    \begin{tcolorbox}[breakable, size=fbox, boxrule=1pt, pad at break*=1mm,colback=cellbackground, colframe=cellborder]
\prompt{In}{incolor}{19}{\boxspacing}
\begin{Verbatim}[commandchars=\\\{\}]
\PY{c+c1}{\PYZsh{} Vemos si hay variables categoricas}
\PY{n}{cat\PYZus{}cols\PYZus{}clf} \PY{o}{=} \PY{n}{df\PYZus{}clf}\PY{o}{.}\PY{n}{select\PYZus{}dtypes}\PY{p}{(}\PY{n}{include}\PY{o}{=}\PY{l+s+s1}{\PYZsq{}}\PY{l+s+s1}{object}\PY{l+s+s1}{\PYZsq{}}\PY{p}{)}\PY{o}{.}\PY{n}{columns}\PY{o}{.}\PY{n}{tolist}\PY{p}{(}\PY{p}{)}
\PY{n+nb}{print}\PY{p}{(}\PY{l+s+sa}{f}\PY{l+s+s1}{\PYZsq{}}\PY{l+s+s1}{Variables categóricas: }\PY{l+s+si}{\PYZob{}}\PY{n}{cat\PYZus{}cols\PYZus{}clf}\PY{l+s+si}{\PYZcb{}}\PY{l+s+s1}{\PYZsq{}}\PY{p}{)}

\PY{c+c1}{\PYZsh{} Quitamos User ID (es solo un identificador, no aporta nada al modelo)}
\PY{c+c1}{\PYZsh{} y Gender (para quedarnos con Age y EstimatedSalary, igual que en la UT6)}
\PY{n}{cols\PYZus{}to\PYZus{}drop} \PY{o}{=} \PY{p}{[}\PY{l+s+s1}{\PYZsq{}}\PY{l+s+s1}{User ID}\PY{l+s+s1}{\PYZsq{}}\PY{p}{,} \PY{l+s+s1}{\PYZsq{}}\PY{l+s+s1}{Gender}\PY{l+s+s1}{\PYZsq{}}\PY{p}{]}
\PY{k}{for} \PY{n}{col} \PY{o+ow}{in} \PY{n}{cols\PYZus{}to\PYZus{}drop}\PY{p}{:}
    \PY{k}{if} \PY{n}{col} \PY{o+ow}{in} \PY{n}{df\PYZus{}clf}\PY{o}{.}\PY{n}{columns}\PY{p}{:}
        \PY{n}{df\PYZus{}clf} \PY{o}{=} \PY{n}{df\PYZus{}clf}\PY{o}{.}\PY{n}{drop}\PY{p}{(}\PY{n}{col}\PY{p}{,} \PY{n}{axis}\PY{o}{=}\PY{l+m+mi}{1}\PY{p}{)}
        \PY{n+nb}{print}\PY{p}{(}\PY{l+s+sa}{f}\PY{l+s+s1}{\PYZsq{}}\PY{l+s+s1}{Columna }\PY{l+s+s1}{\PYZdq{}}\PY{l+s+si}{\PYZob{}}\PY{n}{col}\PY{l+s+si}{\PYZcb{}}\PY{l+s+s1}{\PYZdq{}}\PY{l+s+s1}{ eliminada.}\PY{l+s+s1}{\PYZsq{}}\PY{p}{)}

\PY{n+nb}{print}\PY{p}{(}\PY{l+s+sa}{f}\PY{l+s+s1}{\PYZsq{}}\PY{l+s+se}{\PYZbs{}n}\PY{l+s+s1}{Columnas finales: }\PY{l+s+si}{\PYZob{}}\PY{n}{df\PYZus{}clf}\PY{o}{.}\PY{n}{columns}\PY{o}{.}\PY{n}{tolist}\PY{p}{(}\PY{p}{)}\PY{l+s+si}{\PYZcb{}}\PY{l+s+s1}{\PYZsq{}}\PY{p}{)}
\PY{n+nb}{print}\PY{p}{(}\PY{l+s+sa}{f}\PY{l+s+s1}{\PYZsq{}}\PY{l+s+s1}{Tipos de datos:}\PY{l+s+se}{\PYZbs{}n}\PY{l+s+si}{\PYZob{}}\PY{n}{df\PYZus{}clf}\PY{o}{.}\PY{n}{dtypes}\PY{l+s+si}{\PYZcb{}}\PY{l+s+s1}{\PYZsq{}}\PY{p}{)}
\end{Verbatim}
\end{tcolorbox}

    \begin{Verbatim}[commandchars=\\\{\}]
Variables categóricas: ['Gender']
Columna "User ID" eliminada.
Columna "Gender" eliminada.

Columnas finales: ['Age', 'EstimatedSalary', 'Purchased']
Tipos de datos:
Age                int64
EstimatedSalary    int64
Purchased          int64
dtype: object
    \end{Verbatim}

    \textbf{Nota sobre la trampa de las variables dummy:} Tras eliminar User
ID y Gender, el dataset queda solo con variables numéricas (Age,
EstimatedSalary), por lo que no es necesario aplicar one-hot encoding y
la trampa de las variables dummy no aplica.

    \subsubsection{1.2.4 Separar datos en conjuntos de entrenamiento y
test}\label{separar-datos-en-conjuntos-de-entrenamiento-y-test}

    \begin{tcolorbox}[breakable, size=fbox, boxrule=1pt, pad at break*=1mm,colback=cellbackground, colframe=cellborder]
\prompt{In}{incolor}{20}{\boxspacing}
\begin{Verbatim}[commandchars=\\\{\}]
\PY{c+c1}{\PYZsh{} Separamos features y target}
\PY{n}{X\PYZus{}clf} \PY{o}{=} \PY{n}{df\PYZus{}clf}\PY{o}{.}\PY{n}{drop}\PY{p}{(}\PY{l+s+s1}{\PYZsq{}}\PY{l+s+s1}{Purchased}\PY{l+s+s1}{\PYZsq{}}\PY{p}{,} \PY{n}{axis}\PY{o}{=}\PY{l+m+mi}{1}\PY{p}{)}
\PY{n}{y\PYZus{}clf} \PY{o}{=} \PY{n}{df\PYZus{}clf}\PY{p}{[}\PY{l+s+s1}{\PYZsq{}}\PY{l+s+s1}{Purchased}\PY{l+s+s1}{\PYZsq{}}\PY{p}{]}

\PY{c+c1}{\PYZsh{} Dividimos en train (75\PYZpc{}) y test (25\PYZpc{}) \PYZhy{} misma proporción que en UT6}
\PY{n}{X\PYZus{}train\PYZus{}clf}\PY{p}{,} \PY{n}{X\PYZus{}test\PYZus{}clf}\PY{p}{,} \PY{n}{y\PYZus{}train\PYZus{}clf}\PY{p}{,} \PY{n}{y\PYZus{}test\PYZus{}clf} \PY{o}{=} \PY{n}{train\PYZus{}test\PYZus{}split}\PY{p}{(}
    \PY{n}{X\PYZus{}clf}\PY{p}{,} \PY{n}{y\PYZus{}clf}\PY{p}{,} \PY{n}{test\PYZus{}size}\PY{o}{=}\PY{l+m+mf}{0.25}\PY{p}{,} \PY{n}{random\PYZus{}state}\PY{o}{=}\PY{l+m+mi}{0}
\PY{p}{)}

\PY{n+nb}{print}\PY{p}{(}\PY{l+s+sa}{f}\PY{l+s+s1}{\PYZsq{}}\PY{l+s+s1}{Train: }\PY{l+s+si}{\PYZob{}}\PY{n}{X\PYZus{}train\PYZus{}clf}\PY{o}{.}\PY{n}{shape}\PY{p}{[}\PY{l+m+mi}{0}\PY{p}{]}\PY{l+s+si}{\PYZcb{}}\PY{l+s+s1}{ muestras}\PY{l+s+s1}{\PYZsq{}}\PY{p}{)}
\PY{n+nb}{print}\PY{p}{(}\PY{l+s+sa}{f}\PY{l+s+s1}{\PYZsq{}}\PY{l+s+s1}{Test:  }\PY{l+s+si}{\PYZob{}}\PY{n}{X\PYZus{}test\PYZus{}clf}\PY{o}{.}\PY{n}{shape}\PY{p}{[}\PY{l+m+mi}{0}\PY{p}{]}\PY{l+s+si}{\PYZcb{}}\PY{l+s+s1}{ muestras}\PY{l+s+s1}{\PYZsq{}}\PY{p}{)}
\end{Verbatim}
\end{tcolorbox}

    \begin{Verbatim}[commandchars=\\\{\}]
Train: 300 muestras
Test:  100 muestras
    \end{Verbatim}

    \subsubsection{1.2.5 Estandarizar los
datos}\label{estandarizar-los-datos}

    \begin{tcolorbox}[breakable, size=fbox, boxrule=1pt, pad at break*=1mm,colback=cellbackground, colframe=cellborder]
\prompt{In}{incolor}{21}{\boxspacing}
\begin{Verbatim}[commandchars=\\\{\}]
\PY{c+c1}{\PYZsh{} Estandarizamos (fit solo en train)}
\PY{n}{scaler\PYZus{}clf} \PY{o}{=} \PY{n}{StandardScaler}\PY{p}{(}\PY{p}{)}
\PY{n}{X\PYZus{}train\PYZus{}clf\PYZus{}scaled} \PY{o}{=} \PY{n}{scaler\PYZus{}clf}\PY{o}{.}\PY{n}{fit\PYZus{}transform}\PY{p}{(}\PY{n}{X\PYZus{}train\PYZus{}clf}\PY{p}{)}
\PY{n}{X\PYZus{}test\PYZus{}clf\PYZus{}scaled} \PY{o}{=} \PY{n}{scaler\PYZus{}clf}\PY{o}{.}\PY{n}{transform}\PY{p}{(}\PY{n}{X\PYZus{}test\PYZus{}clf}\PY{p}{)}

\PY{n+nb}{print}\PY{p}{(}\PY{l+s+s1}{\PYZsq{}}\PY{l+s+s1}{Datos estandarizados correctamente.}\PY{l+s+s1}{\PYZsq{}}\PY{p}{)}
\end{Verbatim}
\end{tcolorbox}

    \begin{Verbatim}[commandchars=\\\{\}]
Datos estandarizados correctamente.
    \end{Verbatim}

    \subsubsection{1.2.6 Aplicar el modelo
XGBoostClassifier}\label{aplicar-el-modelo-xgboostclassifier}

Igual que en regresión, probamos varias configuraciones y nos quedamos
con la mejor:

    \begin{tcolorbox}[breakable, size=fbox, boxrule=1pt, pad at break*=1mm,colback=cellbackground, colframe=cellborder]
\prompt{In}{incolor}{22}{\boxspacing}
\begin{Verbatim}[commandchars=\\\{\}]
\PY{c+c1}{\PYZsh{} Probamos varias configuraciones tambien para clasificacion}
\PY{n}{configs\PYZus{}clf} \PY{o}{=} \PY{p}{\PYZob{}}
    \PY{l+s+s1}{\PYZsq{}}\PY{l+s+s1}{Conservador}\PY{l+s+s1}{\PYZsq{}}\PY{p}{:} \PY{p}{\PYZob{}}\PY{l+s+s1}{\PYZsq{}}\PY{l+s+s1}{n\PYZus{}estimators}\PY{l+s+s1}{\PYZsq{}}\PY{p}{:} \PY{l+m+mi}{100}\PY{p}{,} \PY{l+s+s1}{\PYZsq{}}\PY{l+s+s1}{max\PYZus{}depth}\PY{l+s+s1}{\PYZsq{}}\PY{p}{:} \PY{l+m+mi}{3}\PY{p}{,} \PY{l+s+s1}{\PYZsq{}}\PY{l+s+s1}{learning\PYZus{}rate}\PY{l+s+s1}{\PYZsq{}}\PY{p}{:} \PY{l+m+mf}{0.1}\PY{p}{\PYZcb{}}\PY{p}{,}
    \PY{l+s+s1}{\PYZsq{}}\PY{l+s+s1}{Equilibrado}\PY{l+s+s1}{\PYZsq{}}\PY{p}{:} \PY{p}{\PYZob{}}\PY{l+s+s1}{\PYZsq{}}\PY{l+s+s1}{n\PYZus{}estimators}\PY{l+s+s1}{\PYZsq{}}\PY{p}{:} \PY{l+m+mi}{200}\PY{p}{,} \PY{l+s+s1}{\PYZsq{}}\PY{l+s+s1}{max\PYZus{}depth}\PY{l+s+s1}{\PYZsq{}}\PY{p}{:} \PY{l+m+mi}{6}\PY{p}{,} \PY{l+s+s1}{\PYZsq{}}\PY{l+s+s1}{learning\PYZus{}rate}\PY{l+s+s1}{\PYZsq{}}\PY{p}{:} \PY{l+m+mf}{0.1}\PY{p}{\PYZcb{}}\PY{p}{,}
    \PY{l+s+s1}{\PYZsq{}}\PY{l+s+s1}{Agresivo}\PY{l+s+s1}{\PYZsq{}}\PY{p}{:}    \PY{p}{\PYZob{}}\PY{l+s+s1}{\PYZsq{}}\PY{l+s+s1}{n\PYZus{}estimators}\PY{l+s+s1}{\PYZsq{}}\PY{p}{:} \PY{l+m+mi}{500}\PY{p}{,} \PY{l+s+s1}{\PYZsq{}}\PY{l+s+s1}{max\PYZus{}depth}\PY{l+s+s1}{\PYZsq{}}\PY{p}{:} \PY{l+m+mi}{8}\PY{p}{,} \PY{l+s+s1}{\PYZsq{}}\PY{l+s+s1}{learning\PYZus{}rate}\PY{l+s+s1}{\PYZsq{}}\PY{p}{:} \PY{l+m+mf}{0.05}\PY{p}{\PYZcb{}}\PY{p}{,}
\PY{p}{\PYZcb{}}

\PY{n+nb}{print}\PY{p}{(}\PY{l+s+s1}{\PYZsq{}}\PY{l+s+s1}{=== Comparativa de hiperparámetros XGBoostClassifier ===}\PY{l+s+se}{\PYZbs{}n}\PY{l+s+s1}{\PYZsq{}}\PY{p}{)}
\PY{n}{mejor\PYZus{}config\PYZus{}clf} \PY{o}{=} \PY{k+kc}{None}
\PY{n}{mejor\PYZus{}acc} \PY{o}{=} \PY{o}{\PYZhy{}}\PY{l+m+mi}{1}

\PY{k}{for} \PY{n}{nombre}\PY{p}{,} \PY{n}{params} \PY{o+ow}{in} \PY{n}{configs\PYZus{}clf}\PY{o}{.}\PY{n}{items}\PY{p}{(}\PY{p}{)}\PY{p}{:}
    \PY{n}{modelo\PYZus{}tmp} \PY{o}{=} \PY{n}{XGBClassifier}\PY{p}{(}\PY{o}{*}\PY{o}{*}\PY{n}{params}\PY{p}{,} \PY{n}{random\PYZus{}state}\PY{o}{=}\PY{l+m+mi}{42}\PY{p}{,} \PY{n}{eval\PYZus{}metric}\PY{o}{=}\PY{l+s+s1}{\PYZsq{}}\PY{l+s+s1}{logloss}\PY{l+s+s1}{\PYZsq{}}\PY{p}{)}
    \PY{n}{scores} \PY{o}{=} \PY{n}{cross\PYZus{}val\PYZus{}score}\PY{p}{(}\PY{n}{modelo\PYZus{}tmp}\PY{p}{,} \PY{n}{X\PYZus{}train\PYZus{}clf\PYZus{}scaled}\PY{p}{,} \PY{n}{y\PYZus{}train\PYZus{}clf}\PY{p}{,} \PY{n}{cv}\PY{o}{=}\PY{l+m+mi}{5}\PY{p}{,} \PY{n}{scoring}\PY{o}{=}\PY{l+s+s1}{\PYZsq{}}\PY{l+s+s1}{accuracy}\PY{l+s+s1}{\PYZsq{}}\PY{p}{)}
    \PY{n}{acc\PYZus{}medio} \PY{o}{=} \PY{n}{scores}\PY{o}{.}\PY{n}{mean}\PY{p}{(}\PY{p}{)}
    \PY{n+nb}{print}\PY{p}{(}\PY{l+s+sa}{f}\PY{l+s+s1}{\PYZsq{}}\PY{l+s+si}{\PYZob{}}\PY{n}{nombre}\PY{l+s+si}{\PYZcb{}}\PY{l+s+s1}{:}\PY{l+s+s1}{\PYZsq{}}\PY{p}{)}
    \PY{n+nb}{print}\PY{p}{(}\PY{l+s+sa}{f}\PY{l+s+s1}{\PYZsq{}}\PY{l+s+s1}{  Params: n\PYZus{}estimators=}\PY{l+s+si}{\PYZob{}}\PY{n}{params}\PY{p}{[}\PY{l+s+s2}{\PYZdq{}}\PY{l+s+s2}{n\PYZus{}estimators}\PY{l+s+s2}{\PYZdq{}}\PY{p}{]}\PY{l+s+si}{\PYZcb{}}\PY{l+s+s1}{, max\PYZus{}depth=}\PY{l+s+si}{\PYZob{}}\PY{n}{params}\PY{p}{[}\PY{l+s+s2}{\PYZdq{}}\PY{l+s+s2}{max\PYZus{}depth}\PY{l+s+s2}{\PYZdq{}}\PY{p}{]}\PY{l+s+si}{\PYZcb{}}\PY{l+s+s1}{, lr=}\PY{l+s+si}{\PYZob{}}\PY{n}{params}\PY{p}{[}\PY{l+s+s2}{\PYZdq{}}\PY{l+s+s2}{learning\PYZus{}rate}\PY{l+s+s2}{\PYZdq{}}\PY{p}{]}\PY{l+s+si}{\PYZcb{}}\PY{l+s+s1}{\PYZsq{}}\PY{p}{)}
    \PY{n+nb}{print}\PY{p}{(}\PY{l+s+sa}{f}\PY{l+s+s1}{\PYZsq{}}\PY{l+s+s1}{  Accuracy medio (CV 5\PYZhy{}fold): }\PY{l+s+si}{\PYZob{}}\PY{n}{acc\PYZus{}medio}\PY{l+s+si}{:}\PY{l+s+s1}{.4f}\PY{l+s+si}{\PYZcb{}}\PY{l+s+s1}{ (+/\PYZhy{} }\PY{l+s+si}{\PYZob{}}\PY{n}{scores}\PY{o}{.}\PY{n}{std}\PY{p}{(}\PY{p}{)}\PY{l+s+si}{:}\PY{l+s+s1}{.4f}\PY{l+s+si}{\PYZcb{}}\PY{l+s+s1}{)}\PY{l+s+se}{\PYZbs{}n}\PY{l+s+s1}{\PYZsq{}}\PY{p}{)}
    \PY{k}{if} \PY{n}{acc\PYZus{}medio} \PY{o}{\PYZgt{}} \PY{n}{mejor\PYZus{}acc}\PY{p}{:}
        \PY{n}{mejor\PYZus{}acc} \PY{o}{=} \PY{n}{acc\PYZus{}medio}
        \PY{n}{mejor\PYZus{}config\PYZus{}clf} \PY{o}{=} \PY{n}{params}
        \PY{n}{mejor\PYZus{}nombre\PYZus{}clf} \PY{o}{=} \PY{n}{nombre}

\PY{n+nb}{print}\PY{p}{(}\PY{l+s+sa}{f}\PY{l+s+s1}{\PYZsq{}}\PY{l+s+s1}{\PYZgt{}\PYZgt{}\PYZgt{} Mejor configuración: }\PY{l+s+si}{\PYZob{}}\PY{n}{mejor\PYZus{}nombre\PYZus{}clf}\PY{l+s+si}{\PYZcb{}}\PY{l+s+s1}{\PYZsq{}}\PY{p}{)}
\PY{n+nb}{print}\PY{p}{(}\PY{l+s+sa}{f}\PY{l+s+s1}{\PYZsq{}}\PY{l+s+s1}{\PYZgt{}\PYZgt{}\PYZgt{} Params: }\PY{l+s+si}{\PYZob{}}\PY{n}{mejor\PYZus{}config\PYZus{}clf}\PY{l+s+si}{\PYZcb{}}\PY{l+s+se}{\PYZbs{}n}\PY{l+s+s1}{\PYZsq{}}\PY{p}{)}

\PY{c+c1}{\PYZsh{} Entrenamos con la mejor configuracion}
\PY{n}{xgb\PYZus{}clf} \PY{o}{=} \PY{n}{XGBClassifier}\PY{p}{(}\PY{o}{*}\PY{o}{*}\PY{n}{mejor\PYZus{}config\PYZus{}clf}\PY{p}{,} \PY{n}{random\PYZus{}state}\PY{o}{=}\PY{l+m+mi}{42}\PY{p}{,} \PY{n}{eval\PYZus{}metric}\PY{o}{=}\PY{l+s+s1}{\PYZsq{}}\PY{l+s+s1}{logloss}\PY{l+s+s1}{\PYZsq{}}\PY{p}{)}
\PY{n}{xgb\PYZus{}clf}\PY{o}{.}\PY{n}{fit}\PY{p}{(}\PY{n}{X\PYZus{}train\PYZus{}clf\PYZus{}scaled}\PY{p}{,} \PY{n}{y\PYZus{}train\PYZus{}clf}\PY{p}{)}
\PY{n}{y\PYZus{}pred\PYZus{}clf} \PY{o}{=} \PY{n}{xgb\PYZus{}clf}\PY{o}{.}\PY{n}{predict}\PY{p}{(}\PY{n}{X\PYZus{}test\PYZus{}clf\PYZus{}scaled}\PY{p}{)}

\PY{n+nb}{print}\PY{p}{(}\PY{l+s+sa}{f}\PY{l+s+s1}{\PYZsq{}}\PY{l+s+s1}{Modelo XGBoostClassifier entrenado con la mejor configuración.}\PY{l+s+s1}{\PYZsq{}}\PY{p}{)}
\end{Verbatim}
\end{tcolorbox}

    \begin{Verbatim}[commandchars=\\\{\}]
=== Comparativa de hiperparámetros XGBoostClassifier ===

Conservador:
  Params: n\_estimators=100, max\_depth=3, lr=0.1
  Accuracy medio (CV 5-fold): 0.8833 (+/- 0.0606)

Equilibrado:
  Params: n\_estimators=200, max\_depth=6, lr=0.1
  Accuracy medio (CV 5-fold): 0.8800 (+/- 0.0609)

Agresivo:
  Params: n\_estimators=500, max\_depth=8, lr=0.05
  Accuracy medio (CV 5-fold): 0.8767 (+/- 0.0564)

>>> Mejor configuración: Conservador
>>> Params: \{'n\_estimators': 100, 'max\_depth': 3, 'learning\_rate': 0.1\}

Modelo XGBoostClassifier entrenado con la mejor configuración.
    \end{Verbatim}

    \subsubsection{1.2.7 Aplicar métricas y comparar con
UT6}\label{aplicar-muxe9tricas-y-comparar-con-ut6}

Evaluamos el modelo con las métricas habituales de clasificación: -
\textbf{Accuracy:} porcentaje de aciertos totales - \textbf{Precision:}
de los que el modelo dice que compran, cuántos realmente compran -
\textbf{Recall:} de los que realmente compran, cuántos detecta el modelo
- \textbf{F1:} equilibrio entre precision y recall - \textbf{Matriz de
confusión:} desglose de aciertos y errores por clase - \textbf{Curva ROC
/ AUC:} capacidad del modelo para separar compradores de no compradores

    \begin{tcolorbox}[breakable, size=fbox, boxrule=1pt, pad at break*=1mm,colback=cellbackground, colframe=cellborder]
\prompt{In}{incolor}{23}{\boxspacing}
\begin{Verbatim}[commandchars=\\\{\}]
\PY{c+c1}{\PYZsh{} Matriz de confusión}
\PY{n}{cm} \PY{o}{=} \PY{n}{confusion\PYZus{}matrix}\PY{p}{(}\PY{n}{y\PYZus{}test\PYZus{}clf}\PY{p}{,} \PY{n}{y\PYZus{}pred\PYZus{}clf}\PY{p}{)}
\PY{n+nb}{print}\PY{p}{(}\PY{l+s+s1}{\PYZsq{}}\PY{l+s+s1}{=== Matriz de Confusión ===}\PY{l+s+s1}{\PYZsq{}}\PY{p}{)}
\PY{n}{plt}\PY{o}{.}\PY{n}{figure}\PY{p}{(}\PY{n}{figsize}\PY{o}{=}\PY{p}{(}\PY{l+m+mi}{6}\PY{p}{,} \PY{l+m+mi}{4}\PY{p}{)}\PY{p}{)}
\PY{n}{sns}\PY{o}{.}\PY{n}{heatmap}\PY{p}{(}\PY{n}{cm}\PY{p}{,} \PY{n}{annot}\PY{o}{=}\PY{k+kc}{True}\PY{p}{,} \PY{n}{fmt}\PY{o}{=}\PY{l+s+s1}{\PYZsq{}}\PY{l+s+s1}{d}\PY{l+s+s1}{\PYZsq{}}\PY{p}{,} \PY{n}{cmap}\PY{o}{=}\PY{l+s+s1}{\PYZsq{}}\PY{l+s+s1}{Blues}\PY{l+s+s1}{\PYZsq{}}\PY{p}{,}
            \PY{n}{xticklabels}\PY{o}{=}\PY{p}{[}\PY{l+s+s1}{\PYZsq{}}\PY{l+s+s1}{No Compra}\PY{l+s+s1}{\PYZsq{}}\PY{p}{,} \PY{l+s+s1}{\PYZsq{}}\PY{l+s+s1}{Compra}\PY{l+s+s1}{\PYZsq{}}\PY{p}{]}\PY{p}{,}
            \PY{n}{yticklabels}\PY{o}{=}\PY{p}{[}\PY{l+s+s1}{\PYZsq{}}\PY{l+s+s1}{No Compra}\PY{l+s+s1}{\PYZsq{}}\PY{p}{,} \PY{l+s+s1}{\PYZsq{}}\PY{l+s+s1}{Compra}\PY{l+s+s1}{\PYZsq{}}\PY{p}{]}\PY{p}{)}
\PY{n}{plt}\PY{o}{.}\PY{n}{ylabel}\PY{p}{(}\PY{l+s+s1}{\PYZsq{}}\PY{l+s+s1}{Real}\PY{l+s+s1}{\PYZsq{}}\PY{p}{)}
\PY{n}{plt}\PY{o}{.}\PY{n}{xlabel}\PY{p}{(}\PY{l+s+s1}{\PYZsq{}}\PY{l+s+s1}{Predicción}\PY{l+s+s1}{\PYZsq{}}\PY{p}{)}
\PY{n}{plt}\PY{o}{.}\PY{n}{title}\PY{p}{(}\PY{l+s+s1}{\PYZsq{}}\PY{l+s+s1}{Matriz de Confusión \PYZhy{} XGBoostClassifier}\PY{l+s+s1}{\PYZsq{}}\PY{p}{)}
\PY{n}{plt}\PY{o}{.}\PY{n}{show}\PY{p}{(}\PY{p}{)}
\end{Verbatim}
\end{tcolorbox}

    \begin{Verbatim}[commandchars=\\\{\}]
=== Matriz de Confusión ===
    \end{Verbatim}

    \begin{center}
    \adjustimage{max size={0.9\linewidth}{0.9\paperheight}}{UT9_autor_XXXXXXXXX_files/UT9_autor_XXXXXXXXX_43_1.png}
    \end{center}
    { \hspace*{\fill} \\}
    
    \begin{tcolorbox}[breakable, size=fbox, boxrule=1pt, pad at break*=1mm,colback=cellbackground, colframe=cellborder]
\prompt{In}{incolor}{24}{\boxspacing}
\begin{Verbatim}[commandchars=\\\{\}]
\PY{c+c1}{\PYZsh{} Reporte de clasificación}
\PY{n+nb}{print}\PY{p}{(}\PY{l+s+s1}{\PYZsq{}}\PY{l+s+s1}{=== Reporte de Clasificación ===}\PY{l+s+s1}{\PYZsq{}}\PY{p}{)}
\PY{n+nb}{print}\PY{p}{(}\PY{n}{classification\PYZus{}report}\PY{p}{(}\PY{n}{y\PYZus{}test\PYZus{}clf}\PY{p}{,} \PY{n}{y\PYZus{}pred\PYZus{}clf}\PY{p}{,} \PY{n}{target\PYZus{}names}\PY{o}{=}\PY{p}{[}\PY{l+s+s1}{\PYZsq{}}\PY{l+s+s1}{No Compra}\PY{l+s+s1}{\PYZsq{}}\PY{p}{,} \PY{l+s+s1}{\PYZsq{}}\PY{l+s+s1}{Compra}\PY{l+s+s1}{\PYZsq{}}\PY{p}{]}\PY{p}{)}\PY{p}{)}

\PY{n}{acc\PYZus{}xgb} \PY{o}{=} \PY{n}{accuracy\PYZus{}score}\PY{p}{(}\PY{n}{y\PYZus{}test\PYZus{}clf}\PY{p}{,} \PY{n}{y\PYZus{}pred\PYZus{}clf}\PY{p}{)}
\PY{n+nb}{print}\PY{p}{(}\PY{l+s+sa}{f}\PY{l+s+s1}{\PYZsq{}}\PY{l+s+s1}{Accuracy: }\PY{l+s+si}{\PYZob{}}\PY{n}{acc\PYZus{}xgb}\PY{l+s+si}{:}\PY{l+s+s1}{.4f}\PY{l+s+si}{\PYZcb{}}\PY{l+s+s1}{\PYZsq{}}\PY{p}{)}
\end{Verbatim}
\end{tcolorbox}

    \begin{Verbatim}[commandchars=\\\{\}]
=== Reporte de Clasificación ===
              precision    recall  f1-score   support

   No Compra       0.96      0.96      0.96        68
      Compra       0.91      0.91      0.91        32

    accuracy                           0.94       100
   macro avg       0.93      0.93      0.93       100
weighted avg       0.94      0.94      0.94       100

Accuracy: 0.9400
    \end{Verbatim}

    \begin{tcolorbox}[breakable, size=fbox, boxrule=1pt, pad at break*=1mm,colback=cellbackground, colframe=cellborder]
\prompt{In}{incolor}{25}{\boxspacing}
\begin{Verbatim}[commandchars=\\\{\}]
\PY{c+c1}{\PYZsh{} Curva ROC}
\PY{n}{fig}\PY{p}{,} \PY{n}{ax} \PY{o}{=} \PY{n}{plt}\PY{o}{.}\PY{n}{subplots}\PY{p}{(}\PY{n}{figsize}\PY{o}{=}\PY{p}{(}\PY{l+m+mi}{8}\PY{p}{,} \PY{l+m+mi}{6}\PY{p}{)}\PY{p}{)}
\PY{n}{RocCurveDisplay}\PY{o}{.}\PY{n}{from\PYZus{}estimator}\PY{p}{(}\PY{n}{xgb\PYZus{}clf}\PY{p}{,} \PY{n}{X\PYZus{}test\PYZus{}clf\PYZus{}scaled}\PY{p}{,} \PY{n}{y\PYZus{}test\PYZus{}clf}\PY{p}{,} \PY{n}{ax}\PY{o}{=}\PY{n}{ax}\PY{p}{)}
\PY{n}{plt}\PY{o}{.}\PY{n}{title}\PY{p}{(}\PY{l+s+s1}{\PYZsq{}}\PY{l+s+s1}{Curva ROC \PYZhy{} XGBoostClassifier}\PY{l+s+s1}{\PYZsq{}}\PY{p}{)}
\PY{n}{plt}\PY{o}{.}\PY{n}{show}\PY{p}{(}\PY{p}{)}
\end{Verbatim}
\end{tcolorbox}

    \begin{center}
    \adjustimage{max size={0.9\linewidth}{0.9\paperheight}}{UT9_autor_XXXXXXXXX_files/UT9_autor_XXXXXXXXX_45_0.png}
    \end{center}
    { \hspace*{\fill} \\}
    
    \begin{tcolorbox}[breakable, size=fbox, boxrule=1pt, pad at break*=1mm,colback=cellbackground, colframe=cellborder]
\prompt{In}{incolor}{26}{\boxspacing}
\begin{Verbatim}[commandchars=\\\{\}]
\PY{c+c1}{\PYZsh{} Comparación con los modelos de clasificación de la UT6}
\PY{n}{comparativa\PYZus{}clf} \PY{o}{=} \PY{n}{pd}\PY{o}{.}\PY{n}{DataFrame}\PY{p}{(}\PY{p}{\PYZob{}}
    \PY{l+s+s1}{\PYZsq{}}\PY{l+s+s1}{Modelo}\PY{l+s+s1}{\PYZsq{}}\PY{p}{:} \PY{p}{[}\PY{l+s+s1}{\PYZsq{}}\PY{l+s+s1}{Logistic Regression (UT6)}\PY{l+s+s1}{\PYZsq{}}\PY{p}{,} \PY{l+s+s1}{\PYZsq{}}\PY{l+s+s1}{KNN (UT6)}\PY{l+s+s1}{\PYZsq{}}\PY{p}{,} \PY{l+s+s1}{\PYZsq{}}\PY{l+s+s1}{SVM Linear (UT6)}\PY{l+s+s1}{\PYZsq{}}\PY{p}{,}
               \PY{l+s+s1}{\PYZsq{}}\PY{l+s+s1}{SVM RBF (UT6)}\PY{l+s+s1}{\PYZsq{}}\PY{p}{,} \PY{l+s+s1}{\PYZsq{}}\PY{l+s+s1}{Naive Bayes (UT6)}\PY{l+s+s1}{\PYZsq{}}\PY{p}{,} \PY{l+s+s1}{\PYZsq{}}\PY{l+s+s1}{Decision Tree (UT6)}\PY{l+s+s1}{\PYZsq{}}\PY{p}{,}
               \PY{l+s+s1}{\PYZsq{}}\PY{l+s+s1}{Random Forest (UT6)}\PY{l+s+s1}{\PYZsq{}}\PY{p}{,} \PY{l+s+s1}{\PYZsq{}}\PY{l+s+s1}{XGBoost (UT9)}\PY{l+s+s1}{\PYZsq{}}\PY{p}{]}\PY{p}{,}
    \PY{l+s+s1}{\PYZsq{}}\PY{l+s+s1}{Accuracy}\PY{l+s+s1}{\PYZsq{}}\PY{p}{:} \PY{p}{[}\PY{l+m+mf}{0.89}\PY{p}{,} \PY{l+m+mf}{0.93}\PY{p}{,} \PY{l+m+mf}{0.90}\PY{p}{,} \PY{l+m+mf}{0.93}\PY{p}{,} \PY{l+m+mf}{0.90}\PY{p}{,} \PY{l+m+mf}{0.91}\PY{p}{,} \PY{l+m+mf}{0.91}\PY{p}{,} \PY{n+nb}{round}\PY{p}{(}\PY{n}{acc\PYZus{}xgb}\PY{p}{,} \PY{l+m+mi}{4}\PY{p}{)}\PY{p}{]}
\PY{p}{\PYZcb{}}\PY{p}{)}

\PY{n+nb}{print}\PY{p}{(}\PY{l+s+s1}{\PYZsq{}}\PY{l+s+s1}{=== Comparativa de modelos de clasificación ===}\PY{l+s+s1}{\PYZsq{}}\PY{p}{)}
\PY{n}{display}\PY{p}{(}\PY{n}{comparativa\PYZus{}clf}\PY{p}{)}

\PY{c+c1}{\PYZsh{} Gráfico comparativo}
\PY{n}{plt}\PY{o}{.}\PY{n}{figure}\PY{p}{(}\PY{n}{figsize}\PY{o}{=}\PY{p}{(}\PY{l+m+mi}{10}\PY{p}{,} \PY{l+m+mi}{5}\PY{p}{)}\PY{p}{)}
\PY{n}{colors} \PY{o}{=} \PY{p}{[}\PY{l+s+s1}{\PYZsq{}}\PY{l+s+s1}{\PYZsh{}3498db}\PY{l+s+s1}{\PYZsq{}}\PY{p}{]} \PY{o}{*} \PY{l+m+mi}{7} \PY{o}{+} \PY{p}{[}\PY{l+s+s1}{\PYZsq{}}\PY{l+s+s1}{\PYZsh{}e74c3c}\PY{l+s+s1}{\PYZsq{}}\PY{p}{]}
\PY{n}{plt}\PY{o}{.}\PY{n}{barh}\PY{p}{(}\PY{n}{comparativa\PYZus{}clf}\PY{p}{[}\PY{l+s+s1}{\PYZsq{}}\PY{l+s+s1}{Modelo}\PY{l+s+s1}{\PYZsq{}}\PY{p}{]}\PY{p}{,} \PY{n}{comparativa\PYZus{}clf}\PY{p}{[}\PY{l+s+s1}{\PYZsq{}}\PY{l+s+s1}{Accuracy}\PY{l+s+s1}{\PYZsq{}}\PY{p}{]}\PY{p}{,} \PY{n}{color}\PY{o}{=}\PY{n}{colors}\PY{p}{)}
\PY{n}{plt}\PY{o}{.}\PY{n}{xlabel}\PY{p}{(}\PY{l+s+s1}{\PYZsq{}}\PY{l+s+s1}{Accuracy}\PY{l+s+s1}{\PYZsq{}}\PY{p}{)}
\PY{n}{plt}\PY{o}{.}\PY{n}{title}\PY{p}{(}\PY{l+s+s1}{\PYZsq{}}\PY{l+s+s1}{Comparativa Accuracy \PYZhy{} Clasificación (mayor es mejor)}\PY{l+s+s1}{\PYZsq{}}\PY{p}{)}
\PY{n}{plt}\PY{o}{.}\PY{n}{xlim}\PY{p}{(}\PY{l+m+mf}{0.8}\PY{p}{,} \PY{l+m+mf}{1.0}\PY{p}{)}
\PY{n}{plt}\PY{o}{.}\PY{n}{tight\PYZus{}layout}\PY{p}{(}\PY{p}{)}
\PY{n}{plt}\PY{o}{.}\PY{n}{show}\PY{p}{(}\PY{p}{)}
\end{Verbatim}
\end{tcolorbox}

    \begin{Verbatim}[commandchars=\\\{\}]
=== Comparativa de modelos de clasificación ===
    \end{Verbatim}

    
    \begin{Verbatim}[commandchars=\\\{\}]
                      Modelo  Accuracy
0  Logistic Regression (UT6)      0.89
1                  KNN (UT6)      0.93
2           SVM Linear (UT6)      0.90
3              SVM RBF (UT6)      0.93
4          Naive Bayes (UT6)      0.90
5        Decision Tree (UT6)      0.91
6        Random Forest (UT6)      0.91
7              XGBoost (UT9)      0.94
    \end{Verbatim}

    
    \begin{center}
    \adjustimage{max size={0.9\linewidth}{0.9\paperheight}}{UT9_autor_XXXXXXXXX_files/UT9_autor_XXXXXXXXX_46_2.png}
    \end{center}
    { \hspace*{\fill} \\}
    
    \textbf{Conclusión clasificación --- ¿qué nos dicen estos resultados?}

\textbf{Sobre la matriz de confusión:} Los números en la diagonal son
los aciertos del modelo. Los de fuera son los errores. Podemos ver
cuántos usuarios compraron y el modelo lo predijo bien (verdaderos
positivos), y cuántos compraron pero el modelo falló (falsos negativos,
se nos ``escapan'' clientes). En un contexto de marketing, esos falsos
negativos son clientes a los que no les habríamos mostrado publicidad y
habríamos perdido la venta.

\textbf{Sobre la curva ROC:} La curva se aleja bastante de la diagonal,
lo cual indica que el modelo distingue bien entre compradores y no
compradores. No está adivinando, realmente ha aprendido un patrón en los
datos.

\textbf{Sobre el accuracy comparado con la UT6:} XGBoost obtiene
resultados muy parecidos a los mejores modelos de la UT6 (KNN y SVM RBF
con 0.93). ¿Por qué no los supera claramente? Porque el dataset es muy
sencillo: solo 2 variables (edad y salario) y 400 muestras. Con datos
tan simples, la frontera entre ``compra'' y ``no compra'' no es muy
complicada y casi cualquier algoritmo la encuentra bien. XGBoost brilla
más cuando hay muchas variables y relaciones complejas entre ellas.

\textbf{¿Para qué nos sirve entonces XGBoost aquí?} Para dos cosas:
primero, demostrar que funciona bien incluso con datos sencillos.
Segundo y más importante, porque al ser un modelo basado en árboles nos
permite usar SHAP después para interpretar las predicciones, algo que
con KNN o SVM sería bastante más difícil.

\textbf{Sobre el desbalanceo:} Hay más usuarios que NO compran que los
que SÍ compran. Por eso no nos fiamos solo del accuracy (podría ser alto
simplemente diciendo ``nadie compra'') y miramos también precision,
recall y F1 de cada clase por separado en el classification report.

    \begin{center}\rule{0.5\linewidth}{0.5pt}\end{center}

\section{2. AutoML con FLAML}\label{automl-con-flaml}

Utilizamos \textbf{FLAML} (Fast and Lightweight AutoML) de Microsoft
como herramienta de AutoML para clasificación. La idea es que en vez de
probar nosotros a mano distintos algoritmos como hicimos en la UT6,
dejamos que FLAML lo haga automáticamente.

\textbf{Nota:} Se utiliza FLAML en lugar de auto-sklearn o PyCaret
porque: - \textbf{auto-sklearn} fue abandonado en 2022 y no es
compatible con Python 3.10+ - \textbf{PyCaret 3.3.2} no soporta Python
3.12 (la versión que trae Google Colab actualmente) - \textbf{FLAML} es
compatible, ligero y hace lo mismo: probar modelos y encontrar el mejor

    \begin{tcolorbox}[breakable, size=fbox, boxrule=1pt, pad at break*=1mm,colback=cellbackground, colframe=cellborder]
\prompt{In}{incolor}{27}{\boxspacing}
\begin{Verbatim}[commandchars=\\\{\}]
\PY{k+kn}{from}\PY{+w}{ }\PY{n+nn}{flaml}\PY{+w}{ }\PY{k+kn}{import} \PY{n}{AutoML}
\PY{k+kn}{from}\PY{+w}{ }\PY{n+nn}{sklearn}\PY{n+nn}{.}\PY{n+nn}{metrics}\PY{+w}{ }\PY{k+kn}{import} \PY{n}{confusion\PYZus{}matrix}\PY{p}{,} \PY{n}{classification\PYZus{}report}\PY{p}{,} \PY{n}{accuracy\PYZus{}score}\PY{p}{,} \PY{n}{RocCurveDisplay}
\PY{n+nb}{print}\PY{p}{(}\PY{l+s+s1}{\PYZsq{}}\PY{l+s+s1}{FLAML importado correctamente.}\PY{l+s+s1}{\PYZsq{}}\PY{p}{)}
\end{Verbatim}
\end{tcolorbox}

    \begin{Verbatim}[commandchars=\\\{\}]
FLAML importado correctamente.
    \end{Verbatim}

    \subsubsection{2.1 Preparar los datos}\label{preparar-los-datos}

Utilizamos el mismo dataset Social Network Ads. FLAML acepta datos
escalados directamente.

    \begin{tcolorbox}[breakable, size=fbox, boxrule=1pt, pad at break*=1mm,colback=cellbackground, colframe=cellborder]
\prompt{In}{incolor}{28}{\boxspacing}
\begin{Verbatim}[commandchars=\\\{\}]
\PY{c+c1}{\PYZsh{} Usamos los datos ya preparados de la sección 1.2 (Social Network Ads)}
\PY{c+c1}{\PYZsh{} X\PYZus{}train\PYZus{}clf\PYZus{}scaled, X\PYZus{}test\PYZus{}clf\PYZus{}scaled, y\PYZus{}train\PYZus{}clf, y\PYZus{}test\PYZus{}clf ya están definidos}

\PY{n+nb}{print}\PY{p}{(}\PY{l+s+sa}{f}\PY{l+s+s1}{\PYZsq{}}\PY{l+s+s1}{Datos de entrenamiento: }\PY{l+s+si}{\PYZob{}}\PY{n}{X\PYZus{}train\PYZus{}clf\PYZus{}scaled}\PY{o}{.}\PY{n}{shape}\PY{l+s+si}{\PYZcb{}}\PY{l+s+s1}{\PYZsq{}}\PY{p}{)}
\PY{n+nb}{print}\PY{p}{(}\PY{l+s+sa}{f}\PY{l+s+s1}{\PYZsq{}}\PY{l+s+s1}{Datos de test: }\PY{l+s+si}{\PYZob{}}\PY{n}{X\PYZus{}test\PYZus{}clf\PYZus{}scaled}\PY{o}{.}\PY{n}{shape}\PY{l+s+si}{\PYZcb{}}\PY{l+s+s1}{\PYZsq{}}\PY{p}{)}
\end{Verbatim}
\end{tcolorbox}

    \begin{Verbatim}[commandchars=\\\{\}]
Datos de entrenamiento: (300, 2)
Datos de test: (100, 2)
    \end{Verbatim}

    \subsubsection{2.2 Definir y entrenar el modelo FLAML
AutoML}\label{definir-y-entrenar-el-modelo-flaml-automl}

FLAML explora automáticamente múltiples algoritmos (Random Forest,
XGBoost, LightGBM, Extra Trees, etc.) y sus hiperparámetros dentro de un
presupuesto de tiempo definido.

{\def\LTcaptype{none} % do not increment counter
\begin{longtable}[]{@{}
  >{\raggedright\arraybackslash}p{(\linewidth - 4\tabcolsep) * \real{0.3333}}
  >{\raggedright\arraybackslash}p{(\linewidth - 4\tabcolsep) * \real{0.3333}}
  >{\raggedright\arraybackslash}p{(\linewidth - 4\tabcolsep) * \real{0.3333}}@{}}
\toprule\noalign{}
\begin{minipage}[b]{\linewidth}\raggedright
Parámetro
\end{minipage} & \begin{minipage}[b]{\linewidth}\raggedright
Valor
\end{minipage} & \begin{minipage}[b]{\linewidth}\raggedright
Justificación
\end{minipage} \\
\midrule\noalign{}
\endhead
\bottomrule\noalign{}
\endlastfoot
\texttt{time\_budget=120} & 120 segundos de búsqueda. Con un dataset de
400 muestras y 2 features, es tiempo más que suficiente para explorar
decenas de configuraciones. & \\
\texttt{metric=\textquotesingle{}accuracy\textquotesingle{}} & Usamos
accuracy como métrica de optimización para mantener coherencia con la
comparativa de la UT6. & \\
\texttt{task=\textquotesingle{}classification\textquotesingle{}} &
Indicamos que es un problema de clasificación binaria. & \\
\end{longtable}
}

A diferencia del XGBoost manual (donde nosotros elegimos los
hiperparámetros), FLAML los busca automáticamente y selecciona la mejor
combinación.

    \begin{tcolorbox}[breakable, size=fbox, boxrule=1pt, pad at break*=1mm,colback=cellbackground, colframe=cellborder]
\prompt{In}{incolor}{29}{\boxspacing}
\begin{Verbatim}[commandchars=\\\{\}]
\PY{c+c1}{\PYZsh{} Configuramos FLAML para que busque el mejor modelo automaticamente}
\PY{c+c1}{\PYZsh{} Le damos 120 segundos de tiempo para probar distintos algoritmos y configuraciones}
\PY{c+c1}{\PYZsh{} Algo parecido a lo que hace PyCaret con compare\PYZus{}models(), pero mas ligero}
\PY{n}{automl} \PY{o}{=} \PY{n}{AutoML}\PY{p}{(}\PY{p}{)}

\PY{n}{settings} \PY{o}{=} \PY{p}{\PYZob{}}
    \PY{l+s+s2}{\PYZdq{}}\PY{l+s+s2}{time\PYZus{}budget}\PY{l+s+s2}{\PYZdq{}}\PY{p}{:} \PY{l+m+mi}{120}\PY{p}{,}          \PY{c+c1}{\PYZsh{} 2 minutos de busqueda}
    \PY{l+s+s2}{\PYZdq{}}\PY{l+s+s2}{metric}\PY{l+s+s2}{\PYZdq{}}\PY{p}{:} \PY{l+s+s1}{\PYZsq{}}\PY{l+s+s1}{accuracy}\PY{l+s+s1}{\PYZsq{}}\PY{p}{,}        \PY{c+c1}{\PYZsh{} Optimizamos accuracy para comparar con la UT6}
    \PY{l+s+s2}{\PYZdq{}}\PY{l+s+s2}{task}\PY{l+s+s2}{\PYZdq{}}\PY{p}{:} \PY{l+s+s1}{\PYZsq{}}\PY{l+s+s1}{classification}\PY{l+s+s1}{\PYZsq{}}\PY{p}{,}    \PY{c+c1}{\PYZsh{} Problema de clasificacion binaria}
    \PY{l+s+s2}{\PYZdq{}}\PY{l+s+s2}{verbose}\PY{l+s+s2}{\PYZdq{}}\PY{p}{:} \PY{l+m+mi}{0}                 \PY{c+c1}{\PYZsh{} Para que no llene la pantalla de logs}
\PY{p}{\PYZcb{}}

\PY{n+nb}{print}\PY{p}{(}\PY{l+s+s2}{\PYZdq{}}\PY{l+s+s2}{Iniciando optimización AutoML con FLAML...}\PY{l+s+s2}{\PYZdq{}}\PY{p}{)}
\PY{n+nb}{print}\PY{p}{(}\PY{l+s+s2}{\PYZdq{}}\PY{l+s+s2}{Probando múltiples algoritmos y configuraciones...}\PY{l+s+s2}{\PYZdq{}}\PY{p}{)}

\PY{n}{automl}\PY{o}{.}\PY{n}{fit}\PY{p}{(}\PY{n}{X\PYZus{}train}\PY{o}{=}\PY{n}{X\PYZus{}train\PYZus{}clf\PYZus{}scaled}\PY{p}{,} \PY{n}{y\PYZus{}train}\PY{o}{=}\PY{n}{y\PYZus{}train\PYZus{}clf}\PY{p}{,} \PY{o}{*}\PY{o}{*}\PY{n}{settings}\PY{p}{)}

\PY{n+nb}{print}\PY{p}{(}\PY{l+s+sa}{f}\PY{l+s+s2}{\PYZdq{}}\PY{l+s+se}{\PYZbs{}n}\PY{l+s+s2}{Optimización completada.}\PY{l+s+s2}{\PYZdq{}}\PY{p}{)}
\PY{n+nb}{print}\PY{p}{(}\PY{l+s+sa}{f}\PY{l+s+s2}{\PYZdq{}}\PY{l+s+s2}{Mejor modelo: }\PY{l+s+si}{\PYZob{}}\PY{n}{automl}\PY{o}{.}\PY{n}{best\PYZus{}estimator}\PY{l+s+si}{\PYZcb{}}\PY{l+s+s2}{\PYZdq{}}\PY{p}{)}
\PY{n+nb}{print}\PY{p}{(}\PY{l+s+sa}{f}\PY{l+s+s2}{\PYZdq{}}\PY{l+s+s2}{Mejor configuración: }\PY{l+s+si}{\PYZob{}}\PY{n}{automl}\PY{o}{.}\PY{n}{best\PYZus{}config}\PY{l+s+si}{\PYZcb{}}\PY{l+s+s2}{\PYZdq{}}\PY{p}{)}
\end{Verbatim}
\end{tcolorbox}

    \begin{Verbatim}[commandchars=\\\{\}]
Iniciando optimización AutoML con FLAML{\ldots}
Probando múltiples algoritmos y configuraciones{\ldots}
    \end{Verbatim}

    \begin{Verbatim}[commandchars=\\\{\}]
INFO:flaml.tune.searcher.blendsearch:No low-cost partial config given to the
search algorithm. For cost-frugal search, consider providing low-cost values for
cost-related hps via 'low\_cost\_partial\_config'. More info can be found at
https://microsoft.github.io/FLAML/docs/FAQ\#about-low\_cost\_partial\_config-in-tune
INFO:flaml.tune.searcher.blendsearch:No low-cost partial config given to the
search algorithm. For cost-frugal search, consider providing low-cost values for
cost-related hps via 'low\_cost\_partial\_config'. More info can be found at
https://microsoft.github.io/FLAML/docs/FAQ\#about-low\_cost\_partial\_config-in-tune
    \end{Verbatim}

    \begin{Verbatim}[commandchars=\\\{\}]

Optimización completada.
Mejor modelo: xgboost
Mejor configuración: \{'n\_estimators': 5, 'max\_leaves': 6, 'min\_child\_weight':
np.float64(0.27412241581403696), 'learning\_rate':
np.float64(0.32006481654872115), 'subsample': np.float64(0.8960658749171345),
'colsample\_bylevel': np.float64(0.9864070218258665), 'colsample\_bytree': 1.0,
'reg\_alpha': 0.0009765625, 'reg\_lambda': np.float64(0.04222314963298577)\}
    \end{Verbatim}

    \subsubsection{2.3 Evaluar el modelo
AutoML}\label{evaluar-el-modelo-automl}

    \begin{tcolorbox}[breakable, size=fbox, boxrule=1pt, pad at break*=1mm,colback=cellbackground, colframe=cellborder]
\prompt{In}{incolor}{30}{\boxspacing}
\begin{Verbatim}[commandchars=\\\{\}]
\PY{c+c1}{\PYZsh{} Predicciones del modelo AutoML}
\PY{n}{y\PYZus{}pred\PYZus{}automl} \PY{o}{=} \PY{n}{automl}\PY{o}{.}\PY{n}{predict}\PY{p}{(}\PY{n}{X\PYZus{}test\PYZus{}clf\PYZus{}scaled}\PY{p}{)}
\PY{n}{acc\PYZus{}automl} \PY{o}{=} \PY{n}{accuracy\PYZus{}score}\PY{p}{(}\PY{n}{y\PYZus{}test\PYZus{}clf}\PY{p}{,} \PY{n}{y\PYZus{}pred\PYZus{}automl}\PY{p}{)}

\PY{c+c1}{\PYZsh{} Matriz de confusión}
\PY{n}{cm\PYZus{}automl} \PY{o}{=} \PY{n}{confusion\PYZus{}matrix}\PY{p}{(}\PY{n}{y\PYZus{}test\PYZus{}clf}\PY{p}{,} \PY{n}{y\PYZus{}pred\PYZus{}automl}\PY{p}{)}
\PY{n}{plt}\PY{o}{.}\PY{n}{figure}\PY{p}{(}\PY{n}{figsize}\PY{o}{=}\PY{p}{(}\PY{l+m+mi}{6}\PY{p}{,} \PY{l+m+mi}{4}\PY{p}{)}\PY{p}{)}
\PY{n}{sns}\PY{o}{.}\PY{n}{heatmap}\PY{p}{(}\PY{n}{cm\PYZus{}automl}\PY{p}{,} \PY{n}{annot}\PY{o}{=}\PY{k+kc}{True}\PY{p}{,} \PY{n}{fmt}\PY{o}{=}\PY{l+s+s1}{\PYZsq{}}\PY{l+s+s1}{d}\PY{l+s+s1}{\PYZsq{}}\PY{p}{,} \PY{n}{cmap}\PY{o}{=}\PY{l+s+s1}{\PYZsq{}}\PY{l+s+s1}{Greens}\PY{l+s+s1}{\PYZsq{}}\PY{p}{,}
            \PY{n}{xticklabels}\PY{o}{=}\PY{p}{[}\PY{l+s+s1}{\PYZsq{}}\PY{l+s+s1}{No Compra}\PY{l+s+s1}{\PYZsq{}}\PY{p}{,} \PY{l+s+s1}{\PYZsq{}}\PY{l+s+s1}{Compra}\PY{l+s+s1}{\PYZsq{}}\PY{p}{]}\PY{p}{,}
            \PY{n}{yticklabels}\PY{o}{=}\PY{p}{[}\PY{l+s+s1}{\PYZsq{}}\PY{l+s+s1}{No Compra}\PY{l+s+s1}{\PYZsq{}}\PY{p}{,} \PY{l+s+s1}{\PYZsq{}}\PY{l+s+s1}{Compra}\PY{l+s+s1}{\PYZsq{}}\PY{p}{]}\PY{p}{)}
\PY{n}{plt}\PY{o}{.}\PY{n}{ylabel}\PY{p}{(}\PY{l+s+s1}{\PYZsq{}}\PY{l+s+s1}{Real}\PY{l+s+s1}{\PYZsq{}}\PY{p}{)}
\PY{n}{plt}\PY{o}{.}\PY{n}{xlabel}\PY{p}{(}\PY{l+s+s1}{\PYZsq{}}\PY{l+s+s1}{Predicción}\PY{l+s+s1}{\PYZsq{}}\PY{p}{)}
\PY{n}{plt}\PY{o}{.}\PY{n}{title}\PY{p}{(}\PY{l+s+sa}{f}\PY{l+s+s1}{\PYZsq{}}\PY{l+s+s1}{Matriz de Confusión \PYZhy{} FLAML AutoML (}\PY{l+s+si}{\PYZob{}}\PY{n}{automl}\PY{o}{.}\PY{n}{best\PYZus{}estimator}\PY{l+s+si}{\PYZcb{}}\PY{l+s+s1}{)}\PY{l+s+s1}{\PYZsq{}}\PY{p}{)}
\PY{n}{plt}\PY{o}{.}\PY{n}{show}\PY{p}{(}\PY{p}{)}

\PY{c+c1}{\PYZsh{} Reporte de clasificación}
\PY{n+nb}{print}\PY{p}{(}\PY{l+s+s1}{\PYZsq{}}\PY{l+s+s1}{=== Reporte de Clasificación \PYZhy{} FLAML AutoML ===}\PY{l+s+s1}{\PYZsq{}}\PY{p}{)}
\PY{n+nb}{print}\PY{p}{(}\PY{n}{classification\PYZus{}report}\PY{p}{(}\PY{n}{y\PYZus{}test\PYZus{}clf}\PY{p}{,} \PY{n}{y\PYZus{}pred\PYZus{}automl}\PY{p}{,} \PY{n}{target\PYZus{}names}\PY{o}{=}\PY{p}{[}\PY{l+s+s1}{\PYZsq{}}\PY{l+s+s1}{No Compra}\PY{l+s+s1}{\PYZsq{}}\PY{p}{,} \PY{l+s+s1}{\PYZsq{}}\PY{l+s+s1}{Compra}\PY{l+s+s1}{\PYZsq{}}\PY{p}{]}\PY{p}{)}\PY{p}{)}
\PY{n+nb}{print}\PY{p}{(}\PY{l+s+sa}{f}\PY{l+s+s1}{\PYZsq{}}\PY{l+s+s1}{Accuracy: }\PY{l+s+si}{\PYZob{}}\PY{n}{acc\PYZus{}automl}\PY{l+s+si}{:}\PY{l+s+s1}{.4f}\PY{l+s+si}{\PYZcb{}}\PY{l+s+s1}{\PYZsq{}}\PY{p}{)}
\end{Verbatim}
\end{tcolorbox}

    \begin{center}
    \adjustimage{max size={0.9\linewidth}{0.9\paperheight}}{UT9_autor_XXXXXXXXX_files/UT9_autor_XXXXXXXXX_55_0.png}
    \end{center}
    { \hspace*{\fill} \\}
    
    \begin{Verbatim}[commandchars=\\\{\}]
=== Reporte de Clasificación - FLAML AutoML ===
              precision    recall  f1-score   support

   No Compra       0.97      0.94      0.96        68
      Compra       0.88      0.94      0.91        32

    accuracy                           0.94       100
   macro avg       0.93      0.94      0.93       100
weighted avg       0.94      0.94      0.94       100

Accuracy: 0.9400
    \end{Verbatim}

    \begin{tcolorbox}[breakable, size=fbox, boxrule=1pt, pad at break*=1mm,colback=cellbackground, colframe=cellborder]
\prompt{In}{incolor}{31}{\boxspacing}
\begin{Verbatim}[commandchars=\\\{\}]
\PY{c+c1}{\PYZsh{} Curva ROC del modelo AutoML}
\PY{n}{fig}\PY{p}{,} \PY{n}{ax} \PY{o}{=} \PY{n}{plt}\PY{o}{.}\PY{n}{subplots}\PY{p}{(}\PY{n}{figsize}\PY{o}{=}\PY{p}{(}\PY{l+m+mi}{8}\PY{p}{,} \PY{l+m+mi}{6}\PY{p}{)}\PY{p}{)}
\PY{n}{RocCurveDisplay}\PY{o}{.}\PY{n}{from\PYZus{}predictions}\PY{p}{(}\PY{n}{y\PYZus{}test\PYZus{}clf}\PY{p}{,} \PY{n}{automl}\PY{o}{.}\PY{n}{predict\PYZus{}proba}\PY{p}{(}\PY{n}{X\PYZus{}test\PYZus{}clf\PYZus{}scaled}\PY{p}{)}\PY{p}{[}\PY{p}{:}\PY{p}{,}\PY{l+m+mi}{1}\PY{p}{]}\PY{p}{,} \PY{n}{ax}\PY{o}{=}\PY{n}{ax}\PY{p}{)}
\PY{n}{plt}\PY{o}{.}\PY{n}{title}\PY{p}{(}\PY{l+s+sa}{f}\PY{l+s+s1}{\PYZsq{}}\PY{l+s+s1}{Curva ROC \PYZhy{} FLAML AutoML (}\PY{l+s+si}{\PYZob{}}\PY{n}{automl}\PY{o}{.}\PY{n}{best\PYZus{}estimator}\PY{l+s+si}{\PYZcb{}}\PY{l+s+s1}{)}\PY{l+s+s1}{\PYZsq{}}\PY{p}{)}
\PY{n}{plt}\PY{o}{.}\PY{n}{show}\PY{p}{(}\PY{p}{)}
\end{Verbatim}
\end{tcolorbox}

    \begin{center}
    \adjustimage{max size={0.9\linewidth}{0.9\paperheight}}{UT9_autor_XXXXXXXXX_files/UT9_autor_XXXXXXXXX_56_0.png}
    \end{center}
    { \hspace*{\fill} \\}
    
    \begin{tcolorbox}[breakable, size=fbox, boxrule=1pt, pad at break*=1mm,colback=cellbackground, colframe=cellborder]
\prompt{In}{incolor}{32}{\boxspacing}
\begin{Verbatim}[commandchars=\\\{\}]
\PY{c+c1}{\PYZsh{} === COMPARATIVA FINAL DETALLADA ===}

\PY{c+c1}{\PYZsh{} Resumen de modelos explorados por FLAML}
\PY{n+nb}{print}\PY{p}{(}\PY{l+s+s1}{\PYZsq{}}\PY{l+s+s1}{=== Modelos explorados por FLAML ===}\PY{l+s+s1}{\PYZsq{}}\PY{p}{)}
\PY{k}{if} \PY{n+nb}{hasattr}\PY{p}{(}\PY{n}{automl}\PY{p}{,} \PY{l+s+s1}{\PYZsq{}}\PY{l+s+s1}{best\PYZus{}config\PYZus{}per\PYZus{}estimator}\PY{l+s+s1}{\PYZsq{}}\PY{p}{)}\PY{p}{:}
    \PY{k}{for} \PY{n}{estimator}\PY{p}{,} \PY{n}{config} \PY{o+ow}{in} \PY{n}{automl}\PY{o}{.}\PY{n}{best\PYZus{}config\PYZus{}per\PYZus{}estimator}\PY{o}{.}\PY{n}{items}\PY{p}{(}\PY{p}{)}\PY{p}{:}
        \PY{k}{if} \PY{n}{config}\PY{p}{:}
            \PY{n+nb}{print}\PY{p}{(}\PY{l+s+sa}{f}\PY{l+s+s1}{\PYZsq{}}\PY{l+s+s1}{  }\PY{l+s+si}{\PYZob{}}\PY{n}{estimator}\PY{l+s+si}{\PYZcb{}}\PY{l+s+s1}{: }\PY{l+s+si}{\PYZob{}}\PY{n}{config}\PY{l+s+si}{\PYZcb{}}\PY{l+s+s1}{\PYZsq{}}\PY{p}{)}

\PY{n+nb}{print}\PY{p}{(}\PY{l+s+sa}{f}\PY{l+s+s1}{\PYZsq{}}\PY{l+s+se}{\PYZbs{}n}\PY{l+s+s1}{Mejor modelo seleccionado por FLAML: }\PY{l+s+si}{\PYZob{}}\PY{n}{automl}\PY{o}{.}\PY{n}{best\PYZus{}estimator}\PY{l+s+si}{\PYZcb{}}\PY{l+s+s1}{\PYZsq{}}\PY{p}{)}
\PY{n+nb}{print}\PY{p}{(}\PY{l+s+sa}{f}\PY{l+s+s1}{\PYZsq{}}\PY{l+s+s1}{Hiperparámetros óptimos encontrados: }\PY{l+s+si}{\PYZob{}}\PY{n}{automl}\PY{o}{.}\PY{n}{best\PYZus{}config}\PY{l+s+si}{\PYZcb{}}\PY{l+s+s1}{\PYZsq{}}\PY{p}{)}

\PY{c+c1}{\PYZsh{} Comparativa completa con todas las métricas}
\PY{k+kn}{from}\PY{+w}{ }\PY{n+nn}{sklearn}\PY{n+nn}{.}\PY{n+nn}{metrics}\PY{+w}{ }\PY{k+kn}{import} \PY{n}{precision\PYZus{}score}\PY{p}{,} \PY{n}{recall\PYZus{}score}\PY{p}{,} \PY{n}{f1\PYZus{}score}

\PY{n}{modelos} \PY{o}{=} \PY{p}{\PYZob{}}
    \PY{l+s+s1}{\PYZsq{}}\PY{l+s+s1}{Logistic Regression (UT6)}\PY{l+s+s1}{\PYZsq{}}\PY{p}{:} \PY{p}{\PYZob{}}\PY{l+s+s1}{\PYZsq{}}\PY{l+s+s1}{Accuracy}\PY{l+s+s1}{\PYZsq{}}\PY{p}{:} \PY{l+m+mf}{0.89}\PY{p}{,} \PY{l+s+s1}{\PYZsq{}}\PY{l+s+s1}{Tipo}\PY{l+s+s1}{\PYZsq{}}\PY{p}{:} \PY{l+s+s1}{\PYZsq{}}\PY{l+s+s1}{Manual}\PY{l+s+s1}{\PYZsq{}}\PY{p}{\PYZcb{}}\PY{p}{,}
    \PY{l+s+s1}{\PYZsq{}}\PY{l+s+s1}{KNN (UT6)}\PY{l+s+s1}{\PYZsq{}}\PY{p}{:} \PY{p}{\PYZob{}}\PY{l+s+s1}{\PYZsq{}}\PY{l+s+s1}{Accuracy}\PY{l+s+s1}{\PYZsq{}}\PY{p}{:} \PY{l+m+mf}{0.93}\PY{p}{,} \PY{l+s+s1}{\PYZsq{}}\PY{l+s+s1}{Tipo}\PY{l+s+s1}{\PYZsq{}}\PY{p}{:} \PY{l+s+s1}{\PYZsq{}}\PY{l+s+s1}{Manual}\PY{l+s+s1}{\PYZsq{}}\PY{p}{\PYZcb{}}\PY{p}{,}
    \PY{l+s+s1}{\PYZsq{}}\PY{l+s+s1}{SVM Linear (UT6)}\PY{l+s+s1}{\PYZsq{}}\PY{p}{:} \PY{p}{\PYZob{}}\PY{l+s+s1}{\PYZsq{}}\PY{l+s+s1}{Accuracy}\PY{l+s+s1}{\PYZsq{}}\PY{p}{:} \PY{l+m+mf}{0.90}\PY{p}{,} \PY{l+s+s1}{\PYZsq{}}\PY{l+s+s1}{Tipo}\PY{l+s+s1}{\PYZsq{}}\PY{p}{:} \PY{l+s+s1}{\PYZsq{}}\PY{l+s+s1}{Manual}\PY{l+s+s1}{\PYZsq{}}\PY{p}{\PYZcb{}}\PY{p}{,}
    \PY{l+s+s1}{\PYZsq{}}\PY{l+s+s1}{SVM RBF (UT6)}\PY{l+s+s1}{\PYZsq{}}\PY{p}{:} \PY{p}{\PYZob{}}\PY{l+s+s1}{\PYZsq{}}\PY{l+s+s1}{Accuracy}\PY{l+s+s1}{\PYZsq{}}\PY{p}{:} \PY{l+m+mf}{0.93}\PY{p}{,} \PY{l+s+s1}{\PYZsq{}}\PY{l+s+s1}{Tipo}\PY{l+s+s1}{\PYZsq{}}\PY{p}{:} \PY{l+s+s1}{\PYZsq{}}\PY{l+s+s1}{Manual}\PY{l+s+s1}{\PYZsq{}}\PY{p}{\PYZcb{}}\PY{p}{,}
    \PY{l+s+s1}{\PYZsq{}}\PY{l+s+s1}{Naive Bayes (UT6)}\PY{l+s+s1}{\PYZsq{}}\PY{p}{:} \PY{p}{\PYZob{}}\PY{l+s+s1}{\PYZsq{}}\PY{l+s+s1}{Accuracy}\PY{l+s+s1}{\PYZsq{}}\PY{p}{:} \PY{l+m+mf}{0.90}\PY{p}{,} \PY{l+s+s1}{\PYZsq{}}\PY{l+s+s1}{Tipo}\PY{l+s+s1}{\PYZsq{}}\PY{p}{:} \PY{l+s+s1}{\PYZsq{}}\PY{l+s+s1}{Manual}\PY{l+s+s1}{\PYZsq{}}\PY{p}{\PYZcb{}}\PY{p}{,}
    \PY{l+s+s1}{\PYZsq{}}\PY{l+s+s1}{Decision Tree (UT6)}\PY{l+s+s1}{\PYZsq{}}\PY{p}{:} \PY{p}{\PYZob{}}\PY{l+s+s1}{\PYZsq{}}\PY{l+s+s1}{Accuracy}\PY{l+s+s1}{\PYZsq{}}\PY{p}{:} \PY{l+m+mf}{0.91}\PY{p}{,} \PY{l+s+s1}{\PYZsq{}}\PY{l+s+s1}{Tipo}\PY{l+s+s1}{\PYZsq{}}\PY{p}{:} \PY{l+s+s1}{\PYZsq{}}\PY{l+s+s1}{Manual}\PY{l+s+s1}{\PYZsq{}}\PY{p}{\PYZcb{}}\PY{p}{,}
    \PY{l+s+s1}{\PYZsq{}}\PY{l+s+s1}{Random Forest (UT6)}\PY{l+s+s1}{\PYZsq{}}\PY{p}{:} \PY{p}{\PYZob{}}\PY{l+s+s1}{\PYZsq{}}\PY{l+s+s1}{Accuracy}\PY{l+s+s1}{\PYZsq{}}\PY{p}{:} \PY{l+m+mf}{0.91}\PY{p}{,} \PY{l+s+s1}{\PYZsq{}}\PY{l+s+s1}{Tipo}\PY{l+s+s1}{\PYZsq{}}\PY{p}{:} \PY{l+s+s1}{\PYZsq{}}\PY{l+s+s1}{Manual}\PY{l+s+s1}{\PYZsq{}}\PY{p}{\PYZcb{}}\PY{p}{,}
\PY{p}{\PYZcb{}}

\PY{c+c1}{\PYZsh{} Métricas detalladas para XGBoost}
\PY{n}{prec\PYZus{}xgb} \PY{o}{=} \PY{n}{precision\PYZus{}score}\PY{p}{(}\PY{n}{y\PYZus{}test\PYZus{}clf}\PY{p}{,} \PY{n}{y\PYZus{}pred\PYZus{}clf}\PY{p}{)}
\PY{n}{rec\PYZus{}xgb} \PY{o}{=} \PY{n}{recall\PYZus{}score}\PY{p}{(}\PY{n}{y\PYZus{}test\PYZus{}clf}\PY{p}{,} \PY{n}{y\PYZus{}pred\PYZus{}clf}\PY{p}{)}
\PY{n}{f1\PYZus{}xgb} \PY{o}{=} \PY{n}{f1\PYZus{}score}\PY{p}{(}\PY{n}{y\PYZus{}test\PYZus{}clf}\PY{p}{,} \PY{n}{y\PYZus{}pred\PYZus{}clf}\PY{p}{)}
\PY{n}{modelos}\PY{p}{[}\PY{l+s+s1}{\PYZsq{}}\PY{l+s+s1}{XGBoost (UT9)}\PY{l+s+s1}{\PYZsq{}}\PY{p}{]} \PY{o}{=} \PY{p}{\PYZob{}}\PY{l+s+s1}{\PYZsq{}}\PY{l+s+s1}{Accuracy}\PY{l+s+s1}{\PYZsq{}}\PY{p}{:} \PY{n+nb}{round}\PY{p}{(}\PY{n}{acc\PYZus{}xgb}\PY{p}{,} \PY{l+m+mi}{4}\PY{p}{)}\PY{p}{,} \PY{l+s+s1}{\PYZsq{}}\PY{l+s+s1}{Precision}\PY{l+s+s1}{\PYZsq{}}\PY{p}{:} \PY{n+nb}{round}\PY{p}{(}\PY{n}{prec\PYZus{}xgb}\PY{p}{,} \PY{l+m+mi}{4}\PY{p}{)}\PY{p}{,}
                             \PY{l+s+s1}{\PYZsq{}}\PY{l+s+s1}{Recall}\PY{l+s+s1}{\PYZsq{}}\PY{p}{:} \PY{n+nb}{round}\PY{p}{(}\PY{n}{rec\PYZus{}xgb}\PY{p}{,} \PY{l+m+mi}{4}\PY{p}{)}\PY{p}{,} \PY{l+s+s1}{\PYZsq{}}\PY{l+s+s1}{F1}\PY{l+s+s1}{\PYZsq{}}\PY{p}{:} \PY{n+nb}{round}\PY{p}{(}\PY{n}{f1\PYZus{}xgb}\PY{p}{,} \PY{l+m+mi}{4}\PY{p}{)}\PY{p}{,} \PY{l+s+s1}{\PYZsq{}}\PY{l+s+s1}{Tipo}\PY{l+s+s1}{\PYZsq{}}\PY{p}{:} \PY{l+s+s1}{\PYZsq{}}\PY{l+s+s1}{Manual}\PY{l+s+s1}{\PYZsq{}}\PY{p}{\PYZcb{}}

\PY{c+c1}{\PYZsh{} Métricas detalladas para FLAML}
\PY{n}{prec\PYZus{}flaml} \PY{o}{=} \PY{n}{precision\PYZus{}score}\PY{p}{(}\PY{n}{y\PYZus{}test\PYZus{}clf}\PY{p}{,} \PY{n}{y\PYZus{}pred\PYZus{}automl}\PY{p}{)}
\PY{n}{rec\PYZus{}flaml} \PY{o}{=} \PY{n}{recall\PYZus{}score}\PY{p}{(}\PY{n}{y\PYZus{}test\PYZus{}clf}\PY{p}{,} \PY{n}{y\PYZus{}pred\PYZus{}automl}\PY{p}{)}
\PY{n}{f1\PYZus{}flaml} \PY{o}{=} \PY{n}{f1\PYZus{}score}\PY{p}{(}\PY{n}{y\PYZus{}test\PYZus{}clf}\PY{p}{,} \PY{n}{y\PYZus{}pred\PYZus{}automl}\PY{p}{)}
\PY{n}{modelos}\PY{p}{[}\PY{l+s+sa}{f}\PY{l+s+s1}{\PYZsq{}}\PY{l+s+s1}{FLAML (}\PY{l+s+si}{\PYZob{}}\PY{n}{automl}\PY{o}{.}\PY{n}{best\PYZus{}estimator}\PY{l+s+si}{\PYZcb{}}\PY{l+s+s1}{) (UT9)}\PY{l+s+s1}{\PYZsq{}}\PY{p}{]} \PY{o}{=} \PY{p}{\PYZob{}}\PY{l+s+s1}{\PYZsq{}}\PY{l+s+s1}{Accuracy}\PY{l+s+s1}{\PYZsq{}}\PY{p}{:} \PY{n+nb}{round}\PY{p}{(}\PY{n}{acc\PYZus{}automl}\PY{p}{,} \PY{l+m+mi}{4}\PY{p}{)}\PY{p}{,} \PY{l+s+s1}{\PYZsq{}}\PY{l+s+s1}{Precision}\PY{l+s+s1}{\PYZsq{}}\PY{p}{:} \PY{n+nb}{round}\PY{p}{(}\PY{n}{prec\PYZus{}flaml}\PY{p}{,} \PY{l+m+mi}{4}\PY{p}{)}\PY{p}{,}
                                                      \PY{l+s+s1}{\PYZsq{}}\PY{l+s+s1}{Recall}\PY{l+s+s1}{\PYZsq{}}\PY{p}{:} \PY{n+nb}{round}\PY{p}{(}\PY{n}{rec\PYZus{}flaml}\PY{p}{,} \PY{l+m+mi}{4}\PY{p}{)}\PY{p}{,} \PY{l+s+s1}{\PYZsq{}}\PY{l+s+s1}{F1}\PY{l+s+s1}{\PYZsq{}}\PY{p}{:} \PY{n+nb}{round}\PY{p}{(}\PY{n}{f1\PYZus{}flaml}\PY{p}{,} \PY{l+m+mi}{4}\PY{p}{)}\PY{p}{,} \PY{l+s+s1}{\PYZsq{}}\PY{l+s+s1}{Tipo}\PY{l+s+s1}{\PYZsq{}}\PY{p}{:} \PY{l+s+s1}{\PYZsq{}}\PY{l+s+s1}{AutoML}\PY{l+s+s1}{\PYZsq{}}\PY{p}{\PYZcb{}}

\PY{c+c1}{\PYZsh{} Tabla comparativa}
\PY{n}{df\PYZus{}comp} \PY{o}{=} \PY{n}{pd}\PY{o}{.}\PY{n}{DataFrame}\PY{p}{(}\PY{n}{modelos}\PY{p}{)}\PY{o}{.}\PY{n}{T}
\PY{n}{df\PYZus{}comp}\PY{o}{.}\PY{n}{index}\PY{o}{.}\PY{n}{name} \PY{o}{=} \PY{l+s+s1}{\PYZsq{}}\PY{l+s+s1}{Modelo}\PY{l+s+s1}{\PYZsq{}}
\PY{n+nb}{print}\PY{p}{(}\PY{l+s+s1}{\PYZsq{}}\PY{l+s+se}{\PYZbs{}n}\PY{l+s+s1}{=== Comparativa Final de Modelos de Clasificación ===}\PY{l+s+s1}{\PYZsq{}}\PY{p}{)}
\PY{n}{display}\PY{p}{(}\PY{n}{df\PYZus{}comp}\PY{p}{)}

\PY{c+c1}{\PYZsh{} Gráfico comparativo}
\PY{n}{fig}\PY{p}{,} \PY{n}{ax} \PY{o}{=} \PY{n}{plt}\PY{o}{.}\PY{n}{subplots}\PY{p}{(}\PY{n}{figsize}\PY{o}{=}\PY{p}{(}\PY{l+m+mi}{12}\PY{p}{,} \PY{l+m+mi}{6}\PY{p}{)}\PY{p}{)}
\PY{n}{colors} \PY{o}{=} \PY{p}{[}\PY{l+s+s1}{\PYZsq{}}\PY{l+s+s1}{\PYZsh{}3498db}\PY{l+s+s1}{\PYZsq{}}\PY{p}{]} \PY{o}{*} \PY{l+m+mi}{7} \PY{o}{+} \PY{p}{[}\PY{l+s+s1}{\PYZsq{}}\PY{l+s+s1}{\PYZsh{}e74c3c}\PY{l+s+s1}{\PYZsq{}}\PY{p}{,} \PY{l+s+s1}{\PYZsq{}}\PY{l+s+s1}{\PYZsh{}2ecc71}\PY{l+s+s1}{\PYZsq{}}\PY{p}{]}
\PY{n}{ax}\PY{o}{.}\PY{n}{barh}\PY{p}{(}\PY{n}{df\PYZus{}comp}\PY{o}{.}\PY{n}{index}\PY{p}{,} \PY{n}{df\PYZus{}comp}\PY{p}{[}\PY{l+s+s1}{\PYZsq{}}\PY{l+s+s1}{Accuracy}\PY{l+s+s1}{\PYZsq{}}\PY{p}{]}\PY{p}{,} \PY{n}{color}\PY{o}{=}\PY{n}{colors}\PY{p}{)}
\PY{n}{ax}\PY{o}{.}\PY{n}{set\PYZus{}xlabel}\PY{p}{(}\PY{l+s+s1}{\PYZsq{}}\PY{l+s+s1}{Accuracy}\PY{l+s+s1}{\PYZsq{}}\PY{p}{)}
\PY{n}{ax}\PY{o}{.}\PY{n}{set\PYZus{}title}\PY{p}{(}\PY{l+s+s1}{\PYZsq{}}\PY{l+s+s1}{Comparativa de Accuracy \PYZhy{} Todos los modelos (UT6 vs UT9)}\PY{l+s+s1}{\PYZsq{}}\PY{p}{)}
\PY{n}{ax}\PY{o}{.}\PY{n}{set\PYZus{}xlim}\PY{p}{(}\PY{l+m+mf}{0.8}\PY{p}{,} \PY{l+m+mf}{1.0}\PY{p}{)}
\PY{n}{ax}\PY{o}{.}\PY{n}{axvline}\PY{p}{(}\PY{n}{x}\PY{o}{=}\PY{l+m+mf}{0.93}\PY{p}{,} \PY{n}{color}\PY{o}{=}\PY{l+s+s1}{\PYZsq{}}\PY{l+s+s1}{gray}\PY{l+s+s1}{\PYZsq{}}\PY{p}{,} \PY{n}{linestyle}\PY{o}{=}\PY{l+s+s1}{\PYZsq{}}\PY{l+s+s1}{\PYZhy{}\PYZhy{}}\PY{l+s+s1}{\PYZsq{}}\PY{p}{,} \PY{n}{alpha}\PY{o}{=}\PY{l+m+mf}{0.5}\PY{p}{,} \PY{n}{label}\PY{o}{=}\PY{l+s+s1}{\PYZsq{}}\PY{l+s+s1}{Mejor UT6 (0.93)}\PY{l+s+s1}{\PYZsq{}}\PY{p}{)}
\PY{n}{ax}\PY{o}{.}\PY{n}{legend}\PY{p}{(}\PY{p}{)}
\PY{n}{plt}\PY{o}{.}\PY{n}{tight\PYZus{}layout}\PY{p}{(}\PY{p}{)}
\PY{n}{plt}\PY{o}{.}\PY{n}{show}\PY{p}{(}\PY{p}{)}

\PY{c+c1}{\PYZsh{} Comparativa XGBoost vs FLAML detallada}
\PY{n+nb}{print}\PY{p}{(}\PY{l+s+s1}{\PYZsq{}}\PY{l+s+se}{\PYZbs{}n}\PY{l+s+s1}{=== XGBoost manual vs FLAML AutoML (detalle) ===}\PY{l+s+s1}{\PYZsq{}}\PY{p}{)}
\PY{n+nb}{print}\PY{p}{(}\PY{l+s+sa}{f}\PY{l+s+s1}{\PYZsq{}}\PY{l+s+si}{\PYZob{}}\PY{l+s+s2}{\PYZdq{}}\PY{l+s+s2}{Métrica}\PY{l+s+s2}{\PYZdq{}}\PY{l+s+si}{:}\PY{l+s+s1}{\PYZlt{}12}\PY{l+s+si}{\PYZcb{}}\PY{l+s+s1}{ }\PY{l+s+si}{\PYZob{}}\PY{l+s+s2}{\PYZdq{}}\PY{l+s+s2}{XGBoost}\PY{l+s+s2}{\PYZdq{}}\PY{l+s+si}{:}\PY{l+s+s1}{\PYZgt{}10}\PY{l+s+si}{\PYZcb{}}\PY{l+s+s1}{ }\PY{l+s+si}{\PYZob{}}\PY{l+s+s2}{\PYZdq{}}\PY{l+s+s2}{FLAML}\PY{l+s+s2}{\PYZdq{}}\PY{l+s+si}{:}\PY{l+s+s1}{\PYZgt{}10}\PY{l+s+si}{\PYZcb{}}\PY{l+s+s1}{ }\PY{l+s+si}{\PYZob{}}\PY{l+s+s2}{\PYZdq{}}\PY{l+s+s2}{Diferencia}\PY{l+s+s2}{\PYZdq{}}\PY{l+s+si}{:}\PY{l+s+s1}{\PYZgt{}12}\PY{l+s+si}{\PYZcb{}}\PY{l+s+s1}{\PYZsq{}}\PY{p}{)}
\PY{n+nb}{print}\PY{p}{(}\PY{l+s+s1}{\PYZsq{}}\PY{l+s+s1}{\PYZhy{}}\PY{l+s+s1}{\PYZsq{}} \PY{o}{*} \PY{l+m+mi}{46}\PY{p}{)}
\PY{k}{for} \PY{n}{metric} \PY{o+ow}{in} \PY{p}{[}\PY{l+s+s1}{\PYZsq{}}\PY{l+s+s1}{Accuracy}\PY{l+s+s1}{\PYZsq{}}\PY{p}{,} \PY{l+s+s1}{\PYZsq{}}\PY{l+s+s1}{Precision}\PY{l+s+s1}{\PYZsq{}}\PY{p}{,} \PY{l+s+s1}{\PYZsq{}}\PY{l+s+s1}{Recall}\PY{l+s+s1}{\PYZsq{}}\PY{p}{,} \PY{l+s+s1}{\PYZsq{}}\PY{l+s+s1}{F1}\PY{l+s+s1}{\PYZsq{}}\PY{p}{]}\PY{p}{:}
    \PY{n}{v\PYZus{}xgb} \PY{o}{=} \PY{n}{modelos}\PY{p}{[}\PY{l+s+s1}{\PYZsq{}}\PY{l+s+s1}{XGBoost (UT9)}\PY{l+s+s1}{\PYZsq{}}\PY{p}{]}\PY{o}{.}\PY{n}{get}\PY{p}{(}\PY{n}{metric}\PY{p}{,} \PY{l+s+s1}{\PYZsq{}}\PY{l+s+s1}{\PYZhy{}}\PY{l+s+s1}{\PYZsq{}}\PY{p}{)}
    \PY{n}{v\PYZus{}flaml} \PY{o}{=} \PY{n}{modelos}\PY{p}{[}\PY{l+s+sa}{f}\PY{l+s+s1}{\PYZsq{}}\PY{l+s+s1}{FLAML (}\PY{l+s+si}{\PYZob{}}\PY{n}{automl}\PY{o}{.}\PY{n}{best\PYZus{}estimator}\PY{l+s+si}{\PYZcb{}}\PY{l+s+s1}{) (UT9)}\PY{l+s+s1}{\PYZsq{}}\PY{p}{]}\PY{o}{.}\PY{n}{get}\PY{p}{(}\PY{n}{metric}\PY{p}{,} \PY{l+s+s1}{\PYZsq{}}\PY{l+s+s1}{\PYZhy{}}\PY{l+s+s1}{\PYZsq{}}\PY{p}{)}
    \PY{k}{if} \PY{n+nb}{isinstance}\PY{p}{(}\PY{n}{v\PYZus{}xgb}\PY{p}{,} \PY{n+nb}{float}\PY{p}{)} \PY{o+ow}{and} \PY{n+nb}{isinstance}\PY{p}{(}\PY{n}{v\PYZus{}flaml}\PY{p}{,} \PY{n+nb}{float}\PY{p}{)}\PY{p}{:}
        \PY{n}{diff} \PY{o}{=} \PY{n}{v\PYZus{}flaml} \PY{o}{\PYZhy{}} \PY{n}{v\PYZus{}xgb}
        \PY{n+nb}{print}\PY{p}{(}\PY{l+s+sa}{f}\PY{l+s+s1}{\PYZsq{}}\PY{l+s+si}{\PYZob{}}\PY{n}{metric}\PY{l+s+si}{:}\PY{l+s+s1}{\PYZlt{}12}\PY{l+s+si}{\PYZcb{}}\PY{l+s+s1}{ }\PY{l+s+si}{\PYZob{}}\PY{n}{v\PYZus{}xgb}\PY{l+s+si}{:}\PY{l+s+s1}{\PYZgt{}10.4f}\PY{l+s+si}{\PYZcb{}}\PY{l+s+s1}{ }\PY{l+s+si}{\PYZob{}}\PY{n}{v\PYZus{}flaml}\PY{l+s+si}{:}\PY{l+s+s1}{\PYZgt{}10.4f}\PY{l+s+si}{\PYZcb{}}\PY{l+s+s1}{ }\PY{l+s+si}{\PYZob{}}\PY{n}{diff}\PY{l+s+si}{:}\PY{l+s+s1}{\PYZgt{}+12.4f}\PY{l+s+si}{\PYZcb{}}\PY{l+s+s1}{\PYZsq{}}\PY{p}{)}
    \PY{k}{else}\PY{p}{:}
        \PY{n+nb}{print}\PY{p}{(}\PY{l+s+sa}{f}\PY{l+s+s1}{\PYZsq{}}\PY{l+s+si}{\PYZob{}}\PY{n}{metric}\PY{l+s+si}{:}\PY{l+s+s1}{\PYZlt{}12}\PY{l+s+si}{\PYZcb{}}\PY{l+s+s1}{ }\PY{l+s+si}{\PYZob{}}\PY{n}{v\PYZus{}xgb}\PY{l+s+si}{:}\PY{l+s+s1}{\PYZgt{}10}\PY{l+s+si}{\PYZcb{}}\PY{l+s+s1}{ }\PY{l+s+si}{\PYZob{}}\PY{n}{v\PYZus{}flaml}\PY{l+s+si}{:}\PY{l+s+s1}{\PYZgt{}10}\PY{l+s+si}{\PYZcb{}}\PY{l+s+s1}{\PYZsq{}}\PY{p}{)}
\end{Verbatim}
\end{tcolorbox}

    \begin{Verbatim}[commandchars=\\\{\}]
=== Modelos explorados por FLAML ===
  lgbm: \{'n\_estimators': 8, 'num\_leaves': 8, 'min\_child\_samples': 10,
'learning\_rate': np.float64(0.18420323200628366), 'log\_max\_bin': 7,
'colsample\_bytree': np.float64(0.7539027007502843), 'reg\_alpha': 0.0009765625,
'reg\_lambda': np.float64(7.9629870251538035)\}
  rf: \{'n\_estimators': 4, 'max\_features': np.float64(0.6853299841049987),
'max\_leaves': 6, 'criterion': np.str\_('gini')\}
  xgboost: \{'n\_estimators': 5, 'max\_leaves': 6, 'min\_child\_weight':
np.float64(0.27412241581403696), 'learning\_rate':
np.float64(0.32006481654872115), 'subsample': np.float64(0.8960658749171345),
'colsample\_bylevel': np.float64(0.9864070218258665), 'colsample\_bytree': 1.0,
'reg\_alpha': 0.0009765625, 'reg\_lambda': np.float64(0.04222314963298577)\}
  extra\_tree: \{'n\_estimators': 4, 'max\_features': 1.0, 'max\_leaves': 9,
'criterion': np.str\_('entropy')\}
  xgb\_limitdepth: \{'n\_estimators': 4, 'max\_depth': 5, 'min\_child\_weight':
np.float64(7.833102153875076), 'learning\_rate': np.float64(0.6168561220049336),
'subsample': 1.0, 'colsample\_bylevel': 1.0, 'colsample\_bytree': 1.0,
'reg\_alpha': np.float64(0.0024251361957650626), 'reg\_lambda':
np.float64(0.6040983387803027)\}
  sgd: \{'penalty': np.str\_('l2'), 'alpha': np.float64(7.303047050810383e-05),
'l1\_ratio': np.float64(0.04895032752817157), 'epsilon': 0.1, 'learning\_rate':
np.str\_('constant'), 'eta0': 0.1, 'power\_t': np.float64(0.3816461369188729),
'average': False, 'loss': np.str\_('modified\_huber')\}
  lrl1: \{'C': np.float64(1.0)\}

Mejor modelo seleccionado por FLAML: xgboost
Hiperparámetros óptimos encontrados: \{'n\_estimators': 5, 'max\_leaves': 6,
'min\_child\_weight': np.float64(0.27412241581403696), 'learning\_rate':
np.float64(0.32006481654872115), 'subsample': np.float64(0.8960658749171345),
'colsample\_bylevel': np.float64(0.9864070218258665), 'colsample\_bytree': 1.0,
'reg\_alpha': 0.0009765625, 'reg\_lambda': np.float64(0.04222314963298577)\}

=== Comparativa Final de Modelos de Clasificación ===
    \end{Verbatim}

    
    \begin{Verbatim}[commandchars=\\\{\}]
                          Accuracy    Tipo Precision  Recall      F1
Modelo                                                              
Logistic Regression (UT6)     0.89  Manual       NaN     NaN     NaN
KNN (UT6)                     0.93  Manual       NaN     NaN     NaN
SVM Linear (UT6)               0.9  Manual       NaN     NaN     NaN
SVM RBF (UT6)                 0.93  Manual       NaN     NaN     NaN
Naive Bayes (UT6)              0.9  Manual       NaN     NaN     NaN
Decision Tree (UT6)           0.91  Manual       NaN     NaN     NaN
Random Forest (UT6)           0.91  Manual       NaN     NaN     NaN
XGBoost (UT9)                 0.94  Manual    0.9062  0.9062  0.9062
FLAML (xgboost) (UT9)         0.94  AutoML    0.8824  0.9375  0.9091
    \end{Verbatim}

    
    \begin{center}
    \adjustimage{max size={0.9\linewidth}{0.9\paperheight}}{UT9_autor_XXXXXXXXX_files/UT9_autor_XXXXXXXXX_57_2.png}
    \end{center}
    { \hspace*{\fill} \\}
    
    \begin{Verbatim}[commandchars=\\\{\}]

=== XGBoost manual vs FLAML AutoML (detalle) ===
Métrica         XGBoost      FLAML   Diferencia
----------------------------------------------
Accuracy         0.9400     0.9400      +0.0000
Precision        0.9062     0.8824      -0.0238
Recall           0.9062     0.9375      +0.0313
F1               0.9062     0.9091      +0.0029
    \end{Verbatim}

    \begin{tcolorbox}[breakable, size=fbox, boxrule=1pt, pad at break*=1mm,colback=cellbackground, colframe=cellborder]
\prompt{In}{incolor}{33}{\boxspacing}
\begin{Verbatim}[commandchars=\\\{\}]
\PY{c+c1}{\PYZsh{} Guardamos el modelo para uso futuro}
\PY{k+kn}{import}\PY{+w}{ }\PY{n+nn}{joblib}
\PY{n}{joblib}\PY{o}{.}\PY{n}{dump}\PY{p}{(}\PY{n}{automl}\PY{p}{,} \PY{l+s+s1}{\PYZsq{}}\PY{l+s+s1}{mejor\PYZus{}modelo\PYZus{}flaml\PYZus{}social\PYZus{}network.pkl}\PY{l+s+s1}{\PYZsq{}}\PY{p}{)}
\PY{n+nb}{print}\PY{p}{(}\PY{l+s+s1}{\PYZsq{}}\PY{l+s+s1}{Modelo guardado como mejor\PYZus{}modelo\PYZus{}flaml\PYZus{}social\PYZus{}network.pkl}\PY{l+s+s1}{\PYZsq{}}\PY{p}{)}
\end{Verbatim}
\end{tcolorbox}

    \begin{Verbatim}[commandchars=\\\{\}]
Modelo guardado como mejor\_modelo\_flaml\_social\_network.pkl
    \end{Verbatim}

    \subsubsection{2.4 Conclusiones AutoML}\label{conclusiones-automl}

FLAML ha explorado automáticamente múltiples algoritmos de clasificación
y ha seleccionado el mejor según la métrica de Accuracy en 120 segundos.

\textbf{¿Qué ha pasado durante esos 120 segundos?} FLAML ha probado por
nosotros varios algoritmos (Random Forest, XGBoost, LightGBM, Extra
Trees, etc.) con distintas configuraciones de hiperparámetros. En la UT6
hicimos esto manualmente: probamos KNN, SVM, Naive Bayes, Decision
Tree\ldots{} uno por uno, ajustando parámetros a mano. FLAML automatiza
todo ese proceso.

\textbf{¿El resultado de FLAML es mejor que el manual?} No
necesariamente. Con un dataset tan pequeño (400 muestras, 2 variables),
la diferencia entre modelos es mínima. Donde AutoML realmente marca la
diferencia es en datasets grandes con muchas variables, donde probar
manualmente todas las combinaciones sería inviable.

\textbf{¿Por qué FLAML y no auto-sklearn o PyCaret?} -
\textbf{auto-sklearn} fue la primera gran librería de AutoML en Python,
pero su desarrollo se detuvo en 2022. No funciona con Python 3.10 o
superior, lo que lo hace incompatible con Google Colab actual (Python
3.12). - \textbf{PyCaret} es más reciente y completa, pero la versión
3.3.2 bloquea explícitamente Python 3.12 con un error que dice ``Please
DOWNGRADE your Python version''. - \textbf{FLAML} es de Microsoft
Research, se mantiene activamente, funciona con Python 3.9 a 3.12, y es
especialmente eficiente: consigue buenos resultados con poco tiempo de
cómputo.

\textbf{Comparativa de resultados:} - \textbf{UT6 (manual):} Mejor
modelo KNN/SVM RBF con Accuracy 0.93 - \textbf{UT9 (XGBoost manual):}
Accuracy comparable - \textbf{UT9 (FLAML AutoML):} FLAML selecciona
automáticamente el algoritmo y los hiperparámetros óptimos

\textbf{Ventajas de AutoML en general:} 1. Automatiza la selección de
algoritmos y la búsqueda de hiperparámetros 2. Reduce el tiempo de
experimentación (de horas a minutos) 3. Puede encontrar combinaciones
que no habríamos probado manualmente 4. Permite al data scientist
centrarse en lo que importa: entender los datos y los resultados

    \begin{center}\rule{0.5\linewidth}{0.5pt}\end{center}

\section{3. Valores SHAP}\label{valores-shap}

SHAP (SHapley Additive exPlanations) viene de la teoría de juegos y nos
permite explicar las predicciones de un modelo. La idea es que a cada
variable se le asigna un valor SHAP que indica cuánto ha contribuido (a
favor o en contra) a la predicción. Es algo parecido a lo que vimos con
el interpret\_model de PyCaret en clase, pero aplicándolo directamente
con la librería SHAP.

Vamos a usarlo sobre el modelo XGBoostClassifier que hemos entrenado con
el dataset Social Network Ads.

    \begin{tcolorbox}[breakable, size=fbox, boxrule=1pt, pad at break*=1mm,colback=cellbackground, colframe=cellborder]
\prompt{In}{incolor}{34}{\boxspacing}
\begin{Verbatim}[commandchars=\\\{\}]
\PY{c+c1}{\PYZsh{} Inicializamos SHAP para que funcionen las visualizaciones en el notebook}
\PY{n}{shap}\PY{o}{.}\PY{n}{initjs}\PY{p}{(}\PY{p}{)}

\PY{c+c1}{\PYZsh{} Creamos el explainer para nuestro modelo XGBoost}
\PY{c+c1}{\PYZsh{} TreeExplainer esta optimizado para modelos de arboles (XGBoost, Random Forest, etc.)}
\PY{n}{explainer} \PY{o}{=} \PY{n}{shap}\PY{o}{.}\PY{n}{TreeExplainer}\PY{p}{(}\PY{n}{xgb\PYZus{}clf}\PY{p}{)}

\PY{c+c1}{\PYZsh{} Calculamos los valores SHAP para todo el conjunto de test}
\PY{c+c1}{\PYZsh{} Los pasamos a DataFrame con nombres de columnas para que los graficos se entiendan}
\PY{n}{X\PYZus{}test\PYZus{}clf\PYZus{}df} \PY{o}{=} \PY{n}{pd}\PY{o}{.}\PY{n}{DataFrame}\PY{p}{(}\PY{n}{X\PYZus{}test\PYZus{}clf\PYZus{}scaled}\PY{p}{,} \PY{n}{columns}\PY{o}{=}\PY{n}{X\PYZus{}clf}\PY{o}{.}\PY{n}{columns}\PY{p}{)}
\PY{n}{shap\PYZus{}values} \PY{o}{=} \PY{n}{explainer}\PY{p}{(}\PY{n}{X\PYZus{}test\PYZus{}clf\PYZus{}df}\PY{p}{)}

\PY{n+nb}{print}\PY{p}{(}\PY{l+s+sa}{f}\PY{l+s+s1}{\PYZsq{}}\PY{l+s+s1}{SHAP values shape: }\PY{l+s+si}{\PYZob{}}\PY{n}{shap\PYZus{}values}\PY{o}{.}\PY{n}{values}\PY{o}{.}\PY{n}{shape}\PY{l+s+si}{\PYZcb{}}\PY{l+s+s1}{\PYZsq{}}\PY{p}{)}
\PY{n+nb}{print}\PY{p}{(}\PY{l+s+s1}{\PYZsq{}}\PY{l+s+s1}{Valores SHAP calculados correctamente.}\PY{l+s+s1}{\PYZsq{}}\PY{p}{)}
\end{Verbatim}
\end{tcolorbox}

    
    \begin{Verbatim}[commandchars=\\\{\}]
<IPython.core.display.HTML object>
    \end{Verbatim}

    
    \begin{Verbatim}[commandchars=\\\{\}]
SHAP values shape: (100, 2)
Valores SHAP calculados correctamente.
    \end{Verbatim}

    \subsection{3.1 Explicar una única
predicción}\label{explicar-una-uxfanica-predicciuxf3n}

    \subsubsection{3.1.1 Gráficos de fuerzas (Force
Plots)}\label{gruxe1ficos-de-fuerzas-force-plots}

Los force plots muestran cómo cada variable contribuye a empujar la
predicción desde el valor base (media) hacia el valor predicho. Las
variables en rojo empujan hacia la clase positiva (compra) y las azules
hacia la clase negativa (no compra).

    \begin{tcolorbox}[breakable, size=fbox, boxrule=1pt, pad at break*=1mm,colback=cellbackground, colframe=cellborder]
\prompt{In}{incolor}{35}{\boxspacing}
\begin{Verbatim}[commandchars=\\\{\}]
\PY{c+c1}{\PYZsh{} Buscamos casos interesantes: aciertos y fallos del modelo}
\PY{n}{aciertos} \PY{o}{=} \PY{p}{[}\PY{p}{]}
\PY{n}{fallos} \PY{o}{=} \PY{p}{[}\PY{p}{]}
\PY{k}{for} \PY{n}{i} \PY{o+ow}{in} \PY{n+nb}{range}\PY{p}{(}\PY{n+nb}{len}\PY{p}{(}\PY{n}{y\PYZus{}pred\PYZus{}clf}\PY{p}{)}\PY{p}{)}\PY{p}{:}
    \PY{k}{if} \PY{n}{y\PYZus{}pred\PYZus{}clf}\PY{p}{[}\PY{n}{i}\PY{p}{]} \PY{o}{==} \PY{n}{y\PYZus{}test\PYZus{}clf}\PY{o}{.}\PY{n}{values}\PY{p}{[}\PY{n}{i}\PY{p}{]}\PY{p}{:}
        \PY{n}{aciertos}\PY{o}{.}\PY{n}{append}\PY{p}{(}\PY{n}{i}\PY{p}{)}
    \PY{k}{else}\PY{p}{:}
        \PY{n}{fallos}\PY{o}{.}\PY{n}{append}\PY{p}{(}\PY{n}{i}\PY{p}{)}

\PY{n+nb}{print}\PY{p}{(}\PY{l+s+sa}{f}\PY{l+s+s1}{\PYZsq{}}\PY{l+s+s1}{Aciertos: }\PY{l+s+si}{\PYZob{}}\PY{n+nb}{len}\PY{p}{(}\PY{n}{aciertos}\PY{p}{)}\PY{l+s+si}{\PYZcb{}}\PY{l+s+s1}{ | Fallos: }\PY{l+s+si}{\PYZob{}}\PY{n+nb}{len}\PY{p}{(}\PY{n}{fallos}\PY{p}{)}\PY{l+s+si}{\PYZcb{}}\PY{l+s+s1}{\PYZsq{}}\PY{p}{)}
\PY{n+nb}{print}\PY{p}{(}\PY{l+s+sa}{f}\PY{l+s+s1}{\PYZsq{}}\PY{l+s+s1}{Indices de fallos: }\PY{l+s+si}{\PYZob{}}\PY{n}{fallos}\PY{l+s+si}{\PYZcb{}}\PY{l+s+s1}{\PYZsq{}}\PY{p}{)}
\end{Verbatim}
\end{tcolorbox}

    \begin{Verbatim}[commandchars=\\\{\}]
Aciertos: 94 | Fallos: 6
Indices de fallos: [9, 15, 31, 81, 85, 95]
    \end{Verbatim}

    \begin{tcolorbox}[breakable, size=fbox, boxrule=1pt, pad at break*=1mm,colback=cellbackground, colframe=cellborder]
\prompt{In}{incolor}{36}{\boxspacing}
\begin{Verbatim}[commandchars=\\\{\}]
\PY{c+c1}{\PYZsh{} CASO 1: Un acierto claro \PYZhy{} el modelo predice bien}
\PY{n}{idx\PYZus{}acierto} \PY{o}{=} \PY{n}{aciertos}\PY{p}{[}\PY{l+m+mi}{0}\PY{p}{]}
\PY{n+nb}{print}\PY{p}{(}\PY{l+s+sa}{f}\PY{l+s+s1}{\PYZsq{}}\PY{l+s+s1}{=== ACIERTO: Observación }\PY{l+s+si}{\PYZob{}}\PY{n}{idx\PYZus{}acierto}\PY{l+s+si}{\PYZcb{}}\PY{l+s+s1}{ ===}\PY{l+s+s1}{\PYZsq{}}\PY{p}{)}
\PY{n+nb}{print}\PY{p}{(}\PY{l+s+sa}{f}\PY{l+s+s1}{\PYZsq{}}\PY{l+s+s1}{Predicción: }\PY{l+s+si}{\PYZob{}}\PY{n}{y\PYZus{}pred\PYZus{}clf}\PY{p}{[}\PY{n}{idx\PYZus{}acierto}\PY{p}{]}\PY{l+s+si}{\PYZcb{}}\PY{l+s+s1}{ | Real: }\PY{l+s+si}{\PYZob{}}\PY{n}{y\PYZus{}test\PYZus{}clf}\PY{o}{.}\PY{n}{values}\PY{p}{[}\PY{n}{idx\PYZus{}acierto}\PY{p}{]}\PY{l+s+si}{\PYZcb{}}\PY{l+s+s1}{\PYZsq{}}\PY{p}{)}
\PY{n+nb}{print}\PY{p}{(}\PY{l+s+sa}{f}\PY{l+s+s1}{\PYZsq{}}\PY{l+s+s1}{Datos: Age=}\PY{l+s+si}{\PYZob{}}\PY{n}{X\PYZus{}test\PYZus{}clf}\PY{o}{.}\PY{n}{values}\PY{p}{[}\PY{n}{idx\PYZus{}acierto}\PY{p}{]}\PY{p}{[}\PY{l+m+mi}{0}\PY{p}{]}\PY{l+s+si}{:}\PY{l+s+s1}{.0f}\PY{l+s+si}{\PYZcb{}}\PY{l+s+s1}{, Salary=}\PY{l+s+si}{\PYZob{}}\PY{n}{X\PYZus{}test\PYZus{}clf}\PY{o}{.}\PY{n}{values}\PY{p}{[}\PY{n}{idx\PYZus{}acierto}\PY{p}{]}\PY{p}{[}\PY{l+m+mi}{1}\PY{p}{]}\PY{l+s+si}{:}\PY{l+s+s1}{.0f}\PY{l+s+si}{\PYZcb{}}\PY{l+s+s1}{\PYZsq{}}\PY{p}{)}
\PY{n}{shap}\PY{o}{.}\PY{n}{force\PYZus{}plot}\PY{p}{(}\PY{n}{explainer}\PY{o}{.}\PY{n}{expected\PYZus{}value}\PY{p}{,} \PY{n}{shap\PYZus{}values}\PY{o}{.}\PY{n}{values}\PY{p}{[}\PY{n}{idx\PYZus{}acierto}\PY{p}{]}\PY{p}{,} \PY{n}{X\PYZus{}test\PYZus{}clf\PYZus{}df}\PY{o}{.}\PY{n}{iloc}\PY{p}{[}\PY{n}{idx\PYZus{}acierto}\PY{p}{]}\PY{p}{,} \PY{n}{matplotlib}\PY{o}{=}\PY{k+kc}{True}\PY{p}{)}
\PY{n}{plt}\PY{o}{.}\PY{n}{show}\PY{p}{(}\PY{p}{)}
\end{Verbatim}
\end{tcolorbox}

    \begin{Verbatim}[commandchars=\\\{\}]
=== ACIERTO: Observación 0 ===
Predicción: 0 | Real: 0
Datos: Age=30, Salary=87000
    \end{Verbatim}

    \begin{center}
    \adjustimage{max size={0.9\linewidth}{0.9\paperheight}}{UT9_autor_XXXXXXXXX_files/UT9_autor_XXXXXXXXX_65_1.png}
    \end{center}
    { \hspace*{\fill} \\}
    
    \begin{tcolorbox}[breakable, size=fbox, boxrule=1pt, pad at break*=1mm,colback=cellbackground, colframe=cellborder]
\prompt{In}{incolor}{37}{\boxspacing}
\begin{Verbatim}[commandchars=\\\{\}]
\PY{c+c1}{\PYZsh{} CASO 2: Otro acierto con valores diferentes}
\PY{n}{idx\PYZus{}acierto2} \PY{o}{=} \PY{n}{aciertos}\PY{p}{[}\PY{n+nb}{len}\PY{p}{(}\PY{n}{aciertos}\PY{p}{)}\PY{o}{/}\PY{o}{/}\PY{l+m+mi}{2}\PY{p}{]}  \PY{c+c1}{\PYZsh{} cogemos uno del medio}
\PY{n+nb}{print}\PY{p}{(}\PY{l+s+sa}{f}\PY{l+s+s1}{\PYZsq{}}\PY{l+s+s1}{=== ACIERTO: Observación }\PY{l+s+si}{\PYZob{}}\PY{n}{idx\PYZus{}acierto2}\PY{l+s+si}{\PYZcb{}}\PY{l+s+s1}{ ===}\PY{l+s+s1}{\PYZsq{}}\PY{p}{)}
\PY{n+nb}{print}\PY{p}{(}\PY{l+s+sa}{f}\PY{l+s+s1}{\PYZsq{}}\PY{l+s+s1}{Predicción: }\PY{l+s+si}{\PYZob{}}\PY{n}{y\PYZus{}pred\PYZus{}clf}\PY{p}{[}\PY{n}{idx\PYZus{}acierto2}\PY{p}{]}\PY{l+s+si}{\PYZcb{}}\PY{l+s+s1}{ | Real: }\PY{l+s+si}{\PYZob{}}\PY{n}{y\PYZus{}test\PYZus{}clf}\PY{o}{.}\PY{n}{values}\PY{p}{[}\PY{n}{idx\PYZus{}acierto2}\PY{p}{]}\PY{l+s+si}{\PYZcb{}}\PY{l+s+s1}{\PYZsq{}}\PY{p}{)}
\PY{n+nb}{print}\PY{p}{(}\PY{l+s+sa}{f}\PY{l+s+s1}{\PYZsq{}}\PY{l+s+s1}{Datos: Age=}\PY{l+s+si}{\PYZob{}}\PY{n}{X\PYZus{}test\PYZus{}clf}\PY{o}{.}\PY{n}{values}\PY{p}{[}\PY{n}{idx\PYZus{}acierto2}\PY{p}{]}\PY{p}{[}\PY{l+m+mi}{0}\PY{p}{]}\PY{l+s+si}{:}\PY{l+s+s1}{.0f}\PY{l+s+si}{\PYZcb{}}\PY{l+s+s1}{, Salary=}\PY{l+s+si}{\PYZob{}}\PY{n}{X\PYZus{}test\PYZus{}clf}\PY{o}{.}\PY{n}{values}\PY{p}{[}\PY{n}{idx\PYZus{}acierto2}\PY{p}{]}\PY{p}{[}\PY{l+m+mi}{1}\PY{p}{]}\PY{l+s+si}{:}\PY{l+s+s1}{.0f}\PY{l+s+si}{\PYZcb{}}\PY{l+s+s1}{\PYZsq{}}\PY{p}{)}
\PY{n}{shap}\PY{o}{.}\PY{n}{force\PYZus{}plot}\PY{p}{(}\PY{n}{explainer}\PY{o}{.}\PY{n}{expected\PYZus{}value}\PY{p}{,} \PY{n}{shap\PYZus{}values}\PY{o}{.}\PY{n}{values}\PY{p}{[}\PY{n}{idx\PYZus{}acierto2}\PY{p}{]}\PY{p}{,} \PY{n}{X\PYZus{}test\PYZus{}clf\PYZus{}df}\PY{o}{.}\PY{n}{iloc}\PY{p}{[}\PY{n}{idx\PYZus{}acierto2}\PY{p}{]}\PY{p}{,} \PY{n}{matplotlib}\PY{o}{=}\PY{k+kc}{True}\PY{p}{)}
\PY{n}{plt}\PY{o}{.}\PY{n}{show}\PY{p}{(}\PY{p}{)}
\end{Verbatim}
\end{tcolorbox}

    \begin{Verbatim}[commandchars=\\\{\}]
=== ACIERTO: Observación 50 ===
Predicción: 1 | Real: 1
Datos: Age=57, Salary=122000
    \end{Verbatim}

    \begin{center}
    \adjustimage{max size={0.9\linewidth}{0.9\paperheight}}{UT9_autor_XXXXXXXXX_files/UT9_autor_XXXXXXXXX_66_1.png}
    \end{center}
    { \hspace*{\fill} \\}
    
    \begin{tcolorbox}[breakable, size=fbox, boxrule=1pt, pad at break*=1mm,colback=cellbackground, colframe=cellborder]
\prompt{In}{incolor}{38}{\boxspacing}
\begin{Verbatim}[commandchars=\\\{\}]
\PY{c+c1}{\PYZsh{} CASO 3: Un FALLO del modelo \PYZhy{} aqui es donde SHAP es mas util}
\PY{c+c1}{\PYZsh{} Nos permite entender POR QUE el modelo se ha equivocado}
\PY{k}{if} \PY{n+nb}{len}\PY{p}{(}\PY{n}{fallos}\PY{p}{)} \PY{o}{\PYZgt{}} \PY{l+m+mi}{0}\PY{p}{:}
    \PY{n}{idx\PYZus{}fallo} \PY{o}{=} \PY{n}{fallos}\PY{p}{[}\PY{l+m+mi}{0}\PY{p}{]}
    \PY{n+nb}{print}\PY{p}{(}\PY{l+s+sa}{f}\PY{l+s+s1}{\PYZsq{}}\PY{l+s+s1}{=== FALLO: Observación }\PY{l+s+si}{\PYZob{}}\PY{n}{idx\PYZus{}fallo}\PY{l+s+si}{\PYZcb{}}\PY{l+s+s1}{ ===}\PY{l+s+s1}{\PYZsq{}}\PY{p}{)}
    \PY{n+nb}{print}\PY{p}{(}\PY{l+s+sa}{f}\PY{l+s+s1}{\PYZsq{}}\PY{l+s+s1}{Predicción: }\PY{l+s+si}{\PYZob{}}\PY{n}{y\PYZus{}pred\PYZus{}clf}\PY{p}{[}\PY{n}{idx\PYZus{}fallo}\PY{p}{]}\PY{l+s+si}{\PYZcb{}}\PY{l+s+s1}{ | Real: }\PY{l+s+si}{\PYZob{}}\PY{n}{y\PYZus{}test\PYZus{}clf}\PY{o}{.}\PY{n}{values}\PY{p}{[}\PY{n}{idx\PYZus{}fallo}\PY{p}{]}\PY{l+s+si}{\PYZcb{}}\PY{l+s+s1}{ ← EL MODELO SE EQUIVOCA}\PY{l+s+s1}{\PYZsq{}}\PY{p}{)}
    \PY{n+nb}{print}\PY{p}{(}\PY{l+s+sa}{f}\PY{l+s+s1}{\PYZsq{}}\PY{l+s+s1}{Datos: Age=}\PY{l+s+si}{\PYZob{}}\PY{n}{X\PYZus{}test\PYZus{}clf}\PY{o}{.}\PY{n}{values}\PY{p}{[}\PY{n}{idx\PYZus{}fallo}\PY{p}{]}\PY{p}{[}\PY{l+m+mi}{0}\PY{p}{]}\PY{l+s+si}{:}\PY{l+s+s1}{.0f}\PY{l+s+si}{\PYZcb{}}\PY{l+s+s1}{, Salary=}\PY{l+s+si}{\PYZob{}}\PY{n}{X\PYZus{}test\PYZus{}clf}\PY{o}{.}\PY{n}{values}\PY{p}{[}\PY{n}{idx\PYZus{}fallo}\PY{p}{]}\PY{p}{[}\PY{l+m+mi}{1}\PY{p}{]}\PY{l+s+si}{:}\PY{l+s+s1}{.0f}\PY{l+s+si}{\PYZcb{}}\PY{l+s+s1}{\PYZsq{}}\PY{p}{)}
    \PY{n}{shap}\PY{o}{.}\PY{n}{force\PYZus{}plot}\PY{p}{(}\PY{n}{explainer}\PY{o}{.}\PY{n}{expected\PYZus{}value}\PY{p}{,} \PY{n}{shap\PYZus{}values}\PY{o}{.}\PY{n}{values}\PY{p}{[}\PY{n}{idx\PYZus{}fallo}\PY{p}{]}\PY{p}{,} \PY{n}{X\PYZus{}test\PYZus{}clf\PYZus{}df}\PY{o}{.}\PY{n}{iloc}\PY{p}{[}\PY{n}{idx\PYZus{}fallo}\PY{p}{]}\PY{p}{,} \PY{n}{matplotlib}\PY{o}{=}\PY{k+kc}{True}\PY{p}{)}
    \PY{n}{plt}\PY{o}{.}\PY{n}{show}\PY{p}{(}\PY{p}{)}
    \PY{n+nb}{print}\PY{p}{(}\PY{l+s+s1}{\PYZsq{}}\PY{l+s+se}{\PYZbs{}n}\PY{l+s+s1}{¿Por qué ha fallado? Mirando los valores SHAP podemos ver que las}\PY{l+s+s1}{\PYZsq{}}\PY{p}{)}
    \PY{n+nb}{print}\PY{p}{(}\PY{l+s+s1}{\PYZsq{}}\PY{l+s+s1}{variables empujaban en una direccion pero la realidad era la contraria.}\PY{l+s+s1}{\PYZsq{}}\PY{p}{)}
    \PY{n+nb}{print}\PY{p}{(}\PY{l+s+s1}{\PYZsq{}}\PY{l+s+s1}{Probablemente es un usuario con caracteristicas }\PY{l+s+s1}{\PYZdq{}}\PY{l+s+s1}{fronterizas}\PY{l+s+s1}{\PYZdq{}}\PY{l+s+s1}{ donde}\PY{l+s+s1}{\PYZsq{}}\PY{p}{)}
    \PY{n+nb}{print}\PY{p}{(}\PY{l+s+s1}{\PYZsq{}}\PY{l+s+s1}{la decision no es clara.}\PY{l+s+s1}{\PYZsq{}}\PY{p}{)}
\PY{k}{else}\PY{p}{:}
    \PY{n+nb}{print}\PY{p}{(}\PY{l+s+s1}{\PYZsq{}}\PY{l+s+s1}{El modelo no ha cometido ningun fallo en el test (poco probable pero posible)}\PY{l+s+s1}{\PYZsq{}}\PY{p}{)}
\end{Verbatim}
\end{tcolorbox}

    \begin{Verbatim}[commandchars=\\\{\}]
=== FALLO: Observación 9 ===
Predicción: 1 | Real: 0 ← EL MODELO SE EQUIVOCA
Datos: Age=47, Salary=43000
    \end{Verbatim}

    \begin{center}
    \adjustimage{max size={0.9\linewidth}{0.9\paperheight}}{UT9_autor_XXXXXXXXX_files/UT9_autor_XXXXXXXXX_67_1.png}
    \end{center}
    { \hspace*{\fill} \\}
    
    \begin{Verbatim}[commandchars=\\\{\}]

¿Por qué ha fallado? Mirando los valores SHAP podemos ver que las
variables empujaban en una direccion pero la realidad era la contraria.
Probablemente es un usuario con caracteristicas "fronterizas" donde
la decision no es clara.
    \end{Verbatim}

    \subsubsection{3.1.2 Gráficos de dependencias (Dependence
Plots)}\label{gruxe1ficos-de-dependencias-dependence-plots}

Los dependence plots muestran la relación entre el valor de una variable
y su impacto SHAP. Permiten ver cómo influye cada variable en la
predicción a lo largo de su rango de valores.

    \begin{tcolorbox}[breakable, size=fbox, boxrule=1pt, pad at break*=1mm,colback=cellbackground, colframe=cellborder]
\prompt{In}{incolor}{39}{\boxspacing}
\begin{Verbatim}[commandchars=\\\{\}]
\PY{c+c1}{\PYZsh{} Grafico de dependencia para cada variable}
\PY{c+c1}{\PYZsh{} Muestra como influye el valor de cada variable en la prediccion}
\PY{c+c1}{\PYZsh{} Parecido al interpret\PYZus{}model(catb, plot=\PYZsq{}correlation\PYZsq{}, feature=...) que vimos con PyCaret}
\PY{k}{for} \PY{n}{col} \PY{o+ow}{in} \PY{n}{X\PYZus{}clf}\PY{o}{.}\PY{n}{columns}\PY{p}{:}
    \PY{n}{shap}\PY{o}{.}\PY{n}{dependence\PYZus{}plot}\PY{p}{(}\PY{n}{col}\PY{p}{,} \PY{n}{shap\PYZus{}values}\PY{o}{.}\PY{n}{values}\PY{p}{,} \PY{n}{X\PYZus{}test\PYZus{}clf\PYZus{}df}\PY{p}{,} \PY{n}{show}\PY{o}{=}\PY{k+kc}{False}\PY{p}{)}
    \PY{n}{plt}\PY{o}{.}\PY{n}{title}\PY{p}{(}\PY{l+s+sa}{f}\PY{l+s+s1}{\PYZsq{}}\PY{l+s+s1}{Dependencia SHAP: }\PY{l+s+si}{\PYZob{}}\PY{n}{col}\PY{l+s+si}{\PYZcb{}}\PY{l+s+s1}{\PYZsq{}}\PY{p}{)}
    \PY{n}{plt}\PY{o}{.}\PY{n}{show}\PY{p}{(}\PY{p}{)}
\end{Verbatim}
\end{tcolorbox}

    \begin{center}
    \adjustimage{max size={0.9\linewidth}{0.9\paperheight}}{UT9_autor_XXXXXXXXX_files/UT9_autor_XXXXXXXXX_69_0.png}
    \end{center}
    { \hspace*{\fill} \\}
    
    \begin{center}
    \adjustimage{max size={0.9\linewidth}{0.9\paperheight}}{UT9_autor_XXXXXXXXX_files/UT9_autor_XXXXXXXXX_69_1.png}
    \end{center}
    { \hspace*{\fill} \\}
    
    \subsubsection{3.1.3 Conclusiones - Predicción
individual}\label{conclusiones---predicciuxf3n-individual}

Los gráficos de fuerza nos permiten entender \textbf{por qué} el modelo
toma una decisión concreta para cada usuario. Esto es algo que con las
métricas globales (accuracy, precision\ldots) no podíamos hacer.

\textbf{¿Qué hemos visto?} - Cuando un usuario tiene \textbf{edad alta y
salario alto}, ambas variables empujan la predicción hacia ``sí compra''
(valores SHAP positivos). Esto tiene sentido: personas con más edad y
más poder adquisitivo tienen más probabilidad de comprar productos
anunciados en redes sociales. - Cuando un usuario es \textbf{joven y con
salario bajo}, ambas variables empujan hacia ``no compra'' (valores SHAP
negativos). También es lógico: un estudiante de 20 años con poco salario
tiene menos probabilidad de comprar. - Los casos más interesantes son
los \textbf{intermedios}: un usuario joven con salario alto, o uno mayor
con salario bajo. Ahí vemos cómo las variables ``compiten'' entre sí, y
el modelo pondera cuál tiene más peso.

\textbf{Los gráficos de dependencia} confirman estas relaciones: - La
relación entre Age y la predicción no es perfectamente lineal: hay un
``salto'' alrededor de los 35-40 años donde la probabilidad de compra
sube significativamente. - EstimatedSalary también tiene un efecto no
lineal: a partir de cierto umbral de salario, la probabilidad de compra
aumenta mucho.

Estas relaciones no lineales son precisamente lo que XGBoost captura
bien y lo que una regresión logística simple no podría modelar con la
misma precisión.

    \subsection{3.2 Explicar el dataset
completo}\label{explicar-el-dataset-completo}

    \subsubsection{3.2.1 Summary Plot}\label{summary-plot}

El summary plot es la visualización global más importante de SHAP.
Muestra: - \textbf{Eje Y:} Variables ordenadas por importancia (de
arriba a abajo) - \textbf{Eje X:} Valor SHAP (impacto en la predicción)
- \textbf{Color:} Valor de la variable (rojo = alto, azul = bajo)

    \begin{tcolorbox}[breakable, size=fbox, boxrule=1pt, pad at break*=1mm,colback=cellbackground, colframe=cellborder]
\prompt{In}{incolor}{40}{\boxspacing}
\begin{Verbatim}[commandchars=\\\{\}]
\PY{c+c1}{\PYZsh{} Summary plot tipo beeswarm (el mas importante)}
\PY{c+c1}{\PYZsh{} Cada punto es una observacion, color rojo = valor alto, azul = valor bajo}
\PY{c+c1}{\PYZsh{} Parecido al interpret\PYZus{}model(catb) que vimos en clase con PyCaret}
\PY{n}{shap}\PY{o}{.}\PY{n}{summary\PYZus{}plot}\PY{p}{(}\PY{n}{shap\PYZus{}values}\PY{o}{.}\PY{n}{values}\PY{p}{,} \PY{n}{X\PYZus{}test\PYZus{}clf\PYZus{}df}\PY{p}{)}
\end{Verbatim}
\end{tcolorbox}

    \begin{center}
    \adjustimage{max size={0.9\linewidth}{0.9\paperheight}}{UT9_autor_XXXXXXXXX_files/UT9_autor_XXXXXXXXX_73_0.png}
    \end{center}
    { \hspace*{\fill} \\}
    
    \begin{tcolorbox}[breakable, size=fbox, boxrule=1pt, pad at break*=1mm,colback=cellbackground, colframe=cellborder]
\prompt{In}{incolor}{41}{\boxspacing}
\begin{Verbatim}[commandchars=\\\{\}]
\PY{c+c1}{\PYZsh{} Summary plot en formato de barras (importancia media de cada variable)}
\PY{c+c1}{\PYZsh{} Mas facil de leer que el beeswarm para ver rapidamente que variable pesa mas}
\PY{n}{shap}\PY{o}{.}\PY{n}{summary\PYZus{}plot}\PY{p}{(}\PY{n}{shap\PYZus{}values}\PY{o}{.}\PY{n}{values}\PY{p}{,} \PY{n}{X\PYZus{}test\PYZus{}clf\PYZus{}df}\PY{p}{,} \PY{n}{plot\PYZus{}type}\PY{o}{=}\PY{l+s+s1}{\PYZsq{}}\PY{l+s+s1}{bar}\PY{l+s+s1}{\PYZsq{}}\PY{p}{)}
\end{Verbatim}
\end{tcolorbox}

    \begin{center}
    \adjustimage{max size={0.9\linewidth}{0.9\paperheight}}{UT9_autor_XXXXXXXXX_files/UT9_autor_XXXXXXXXX_74_0.png}
    \end{center}
    { \hspace*{\fill} \\}
    
    \begin{tcolorbox}[breakable, size=fbox, boxrule=1pt, pad at break*=1mm,colback=cellbackground, colframe=cellborder]
\prompt{In}{incolor}{42}{\boxspacing}
\begin{Verbatim}[commandchars=\\\{\}]
\PY{c+c1}{\PYZsh{} Waterfall plot \PYZhy{} muestra paso a paso como se construye la prediccion}
\PY{c+c1}{\PYZsh{} Partiendo del valor base, va sumando/restando la contribucion de cada variable}
\PY{n}{shap}\PY{o}{.}\PY{n}{waterfall\PYZus{}plot}\PY{p}{(}\PY{n}{shap\PYZus{}values}\PY{p}{[}\PY{l+m+mi}{0}\PY{p}{]}\PY{p}{)}
\end{Verbatim}
\end{tcolorbox}

    \begin{center}
    \adjustimage{max size={0.9\linewidth}{0.9\paperheight}}{UT9_autor_XXXXXXXXX_files/UT9_autor_XXXXXXXXX_75_0.png}
    \end{center}
    { \hspace*{\fill} \\}
    
    \subsubsection{3.2.2 Conclusiones - Dataset
completo}\label{conclusiones---dataset-completo}

Del análisis SHAP global del modelo XGBoost sobre el dataset Social
Network Ads podemos concluir:

\begin{enumerate}
\def\labelenumi{\arabic{enumi}.}
\item
  \textbf{Importancia de variables:} El summary plot nos muestra de un
  vistazo qué variable pesa más en las decisiones del modelo. Podemos
  ver si la edad del usuario influye más que su salario (o al revés) a
  la hora de predecir si comprará el producto.
\item
  \textbf{Dirección del impacto:} Los puntos rojos a la derecha nos
  dicen que valores altos de esa variable empujan hacia ``sí compra'', y
  los azules a la izquierda hacia ``no compra''. Esto nos permite
  entender la lógica del modelo: por ejemplo, si usuarios con mayor edad
  y mayor salario tienen más probabilidad de comprar.
\item
  \textbf{De caja negra a modelo explicable:} Sin SHAP, XGBoost es una
  ``caja negra'' que nos da un resultado pero no sabemos por qué. Con
  SHAP podemos justificar cada predicción. Esto es muy importante en la
  vida real: imagina que una empresa de marketing necesita explicar a su
  cliente \emph{por qué} el modelo recomienda mostrar publicidad a un
  usuario concreto.
\item
  \textbf{Más allá de las métricas globales:} En la UT6 evaluamos los
  modelos solo con métricas como accuracy o precision. SHAP nos da un
  nivel de detalle mucho mayor: no solo sabemos \emph{cuánto} acierta el
  modelo, sino \emph{por qué} toma cada decisión.
\end{enumerate}

\begin{center}\rule{0.5\linewidth}{0.5pt}\end{center}

\section{4. Conclusiones finales}\label{conclusiones-finales}

\subsection{¿Qué hemos hecho en este
trabajo?}\label{quuxe9-hemos-hecho-en-este-trabajo}

Este trabajo conecta lo aprendido en la \textbf{UT5} (regresión), la
\textbf{UT6} (clasificación) y la \textbf{UT9} (XGBoost, AutoML y SHAP).
Hemos usado los mismos datasets de las prácticas anteriores para poder
comparar resultados directamente.

\subsection{4.1 Sobre XGBoost}\label{sobre-xgboost}

XGBoost es un algoritmo de \emph{gradient boosting} que construye muchos
árboles de decisión, donde cada nuevo árbol intenta corregir los errores
del anterior. Es como un equipo donde cada miembro se especializa en
resolver lo que los demás fallan.

\textbf{En regresión (House Prices):} - Obtuvimos un R² competitivo para
predecir precios de vivienda. El modelo captura relaciones no lineales
(por ejemplo, que tener garaje sube el precio, pero no de forma
proporcional al número de plazas). - En la UT5 ya vimos que los modelos
ensemble (Random Forest, Stacking) superaban a la regresión lineal.
XGBoost sigue esa misma línea: combinar varios modelos simples da mejor
resultado que uno solo complejo. - Hemos probado 3 configuraciones de
hiperparámetros con validación cruzada para no quedarnos con valores por
defecto sin justificar.

\textbf{En clasificación (Social Network Ads):} - XGBoost alcanza un
accuracy comparable a los mejores modelos de la UT6 (KNN y SVM RBF con
0.93). - La ventaja de XGBoost sobre SVM es que es más interpretable
(podemos ver la importancia de cada variable) y escala mejor a datasets
grandes.

\subsection{4.2 Sobre AutoML con FLAML}\label{sobre-automl-con-flaml}

La idea de AutoML es sencilla: en vez de que nosotros probemos a mano
diferentes algoritmos y configuraciones (como hicimos en la UT6 probando
KNN, SVM, Naive Bayes, etc.), dejamos que un sistema automático lo haga
por nosotros.

FLAML ha probado automáticamente múltiples algoritmos (Random Forest,
XGBoost, LightGBM, Extra Trees\ldots) y ha buscado los mejores
hiperparámetros para cada uno, todo en 120 segundos. Lo que a nosotros
nos llevó varias prácticas, FLAML lo hace en 2 minutos.

\textbf{¿Significa que AutoML reemplaza al data scientist?} No.~AutoML
automatiza la parte más mecánica (probar modelos), pero sigue siendo
necesario: - Entender el problema y elegir bien los datos - Preparar y
limpiar el dataset - Interpretar los resultados (ahí es donde entra
SHAP)

\textbf{Sobre la elección de FLAML:} El enunciado proponía auto-sklearn
o PyCaret, pero ambas librerías son incompatibles con Python 3.12 (la
versión actual de Google Colab). auto-sklearn fue abandonado en 2022 y
PyCaret 3.3.2 bloquea explícitamente Python 3.12. FLAML de Microsoft es
la alternativa más actual y funcional.

\subsection{4.3 Sobre SHAP y la
interpretabilidad}\label{sobre-shap-y-la-interpretabilidad}

SHAP responde a la pregunta más importante en Machine Learning aplicado:
\textbf{¿por qué el modelo ha tomado esta decisión?}

\begin{itemize}
\tightlist
\item
  Los \textbf{force plots} nos permiten explicar una predicción
  concreta, tanto aciertos como fallos. Hemos visto que cuando el modelo
  falla, SHAP nos ayuda a entender que el usuario tenía características
  ``fronterizas'' donde la decisión no era clara.
\item
  Los \textbf{summary plots} nos dan la foto global: qué variables son
  las más importantes y cómo influyen en general.
\item
  Los \textbf{dependence plots} muestran si la relación es lineal (a más
  edad, más compra) o tiene matices (compra más a partir de cierta edad,
  pero luego se estabiliza).
\end{itemize}

\subsection{4.4 Sobre los datasets
utilizados}\label{sobre-los-datasets-utilizados}

{\def\LTcaptype{none} % do not increment counter
\begin{longtable}[]{@{}
  >{\raggedright\arraybackslash}p{(\linewidth - 6\tabcolsep) * \real{0.2500}}
  >{\raggedright\arraybackslash}p{(\linewidth - 6\tabcolsep) * \real{0.2500}}
  >{\raggedright\arraybackslash}p{(\linewidth - 6\tabcolsep) * \real{0.2500}}
  >{\raggedright\arraybackslash}p{(\linewidth - 6\tabcolsep) * \real{0.2500}}@{}}
\toprule\noalign{}
\begin{minipage}[b]{\linewidth}\raggedright
Dataset
\end{minipage} & \begin{minipage}[b]{\linewidth}\raggedright
Usado en
\end{minipage} & \begin{minipage}[b]{\linewidth}\raggedright
Problema
\end{minipage} & \begin{minipage}[b]{\linewidth}\raggedright
Objetivo
\end{minipage} \\
\midrule\noalign{}
\endhead
\bottomrule\noalign{}
\endlastfoot
House Prices & UT5 y UT9 & Regresión & Predecir el precio de venta de
una vivienda \\
Social Network Ads & UT6 y UT9 & Clasificación & Predecir si un usuario
comprará un producto \\
\end{longtable}
}

Ambos se descargan automáticamente desde Kaggle con \texttt{kagglehub}.
El dataset de clasificación fue proporcionado en la UT6 como ``Ensayo de
Motores'' con columnas renombradas (Age→RPM, EstimatedSalary→Par de
carga), pero al descargarlo de Kaggle descubrí que es el dataset Social
Network Ads original. He optado por usar los nombres originales porque
tienen más sentido con los datos reales.

\subsection{4.5 Resumen: evolución a lo largo del
curso}\label{resumen-evoluciuxf3n-a-lo-largo-del-curso}

{\def\LTcaptype{none} % do not increment counter
\begin{longtable}[]{@{}
  >{\raggedright\arraybackslash}p{(\linewidth - 4\tabcolsep) * \real{0.3333}}
  >{\raggedright\arraybackslash}p{(\linewidth - 4\tabcolsep) * \real{0.3333}}
  >{\raggedright\arraybackslash}p{(\linewidth - 4\tabcolsep) * \real{0.3333}}@{}}
\toprule\noalign{}
\begin{minipage}[b]{\linewidth}\raggedright
UT
\end{minipage} & \begin{minipage}[b]{\linewidth}\raggedright
Qué hicimos
\end{minipage} & \begin{minipage}[b]{\linewidth}\raggedright
Qué aprendimos
\end{minipage} \\
\midrule\noalign{}
\endhead
\bottomrule\noalign{}
\endlastfoot
\textbf{UT5} & Regresión con varios modelos (Lineal, SVR, Random Forest)
& A predecir valores numéricos y comparar modelos con RMSE y R² \\
\textbf{UT6} & Clasificación con 7 modelos diferentes (KNN, SVM, etc.) &
A predecir categorías y evaluar con accuracy, precision, recall \\
\textbf{UT9} & XGBoost + FLAML AutoML + SHAP & Gradient boosting (mejor
rendimiento), AutoML (automatizar la búsqueda) y SHAP (explicar las
decisiones) \\
\end{longtable}
}

La progresión es clara: primero aprendimos a entrenar modelos
manualmente, luego a automatizar ese proceso con AutoML, y finalmente a
explicar \emph{por qué} los modelos toman las decisiones que toman. En
un proyecto real de Machine Learning, las tres habilidades son
necesarias.


    % Add a bibliography block to the postdoc
    
    
    
\end{document}
